% Sensibilisation dans le cadre des catastrophes liées aux mouvements atmosphérique et océanique — SVT 1ère D — Cours signé OnBuch+
% Programme officiel MINESEC. Compiler avec : tectonic ND12-catastrophes-mouvements-atmospheriques-oceaniques.tex
\newcommand{\DOCMATIERE}{SVT}
\newcommand{\DOCNIVEAU}{1ère D}
\newcommand{\DOCMODULE}{Module 3 : L'Éducation à l'Environnement et au Développement Durable}
\newcommand{\DOCLECON}{ND12}
\newcommand{\DOCTITRE}{Sensibilisation dans le cadre des catastrophes liées aux mouvements atmosphérique et océanique}
% OnBuch+ — préambule « cours signés » ENRICHI (compilé avec tectonic / XeLaTeX)
% Système visuel « Pop » d'OnBuch+ : fond crème (jamais blanc), bordures encre
% épaisses, ombres « dures » décalées, titres Archivo Black en capitales.
% Version MATHÉMATIQUES : graphes (pgfplots/tikz), boîtes pédagogiques
% (définition, propriété, méthode, attention, le savais-tu, exercice résolu,
% exercice, TP, bilan), page de garde et signature OnBuch+ ancrée en bas.
%
% Macros attendues dans main.tex AVANT \input{preamble} :
%   \DOCMATIERE  \DOCNIVEAU  \DOCMODULE  \DOCLECON (numéro)  \DOCTITRE
\documentclass[11pt]{article}

\usepackage{fontspec}
\usepackage[french]{babel}
\usepackage[a4paper,margin=2cm,top=3.4cm,bottom=2.8cm,headheight=1.4cm,footskip=1.3cm]{geometry}
\usepackage{amsmath,amssymb,amsfonts}
\usepackage{mathtools} % \xmapsto, \coloneqq... souvent utilisés par les modèles
\usepackage{cancel} % simplifications barrées (bilans de forces)
% \abs{z} (module, valeur absolue) et \norm{u} (norme), utilisés par les modèles
\providecommand{\abs}{}\let\abs\relax\DeclarePairedDelimiter{\abs}{\lvert}{\rvert}
\providecommand{\norm}{}\let\norm\relax\DeclarePairedDelimiter{\norm}{\lVert}{\rVert}
\usepackage[table]{xcolor} % [table] : fournit aussi \rowcolors (alternance de lignes)
\usepackage{pagecolor}
\usepackage{tikz}
\usetikzlibrary{arrows.meta,positioning,calc,shapes.geometric,shapes.misc,shapes.arrows,decorations.pathmorphing,decorations.pathreplacing,decorations.markings,patterns,shadows,mindmap,trees,fit,backgrounds,matrix,intersections,angles,quotes}
\usepackage{pgfplots}
\usepackage[version=4]{mhchem} % équations nucléaires \ce{^{235}_{92}U}
\usepackage[european,siunitx]{circuitikz} % schémas électriques
\pgfplotsset{compat=1.18}
\usepgfplotslibrary{fillbetween,groupplots} % aires entre courbes (calcul intégral)
\usepackage{enumitem}
\usepackage{fancyhdr}
% [most] charge toutes les bibliothèques tcolorbox (listings, minted, raster,
% documentation...) et gonfle inutilement la mémoire TeX sur un document long
% (~300 boîtes) : on ne charge que ce qui est réellement utilisé.
\usepackage[skins,breakable]{tcolorbox}
\usepackage{graphicx}
\usepackage{colortbl}
\usepackage{booktabs}
\usepackage{tabularx}
\usepackage{diagbox} % case d'en-tête diagonale des tableaux à double entrée
% Colonnes extensibles centrée / gauche / droite (C, L, R), souvent utilisées par les modèles
\newcolumntype{C}{>{\centering\arraybackslash}X}
\newcolumntype{L}{>{\raggedright\arraybackslash}X}
\newcolumntype{R}{>{\raggedleft\arraybackslash}X}
\usepackage{array}
\usepackage{multicol}
\usepackage{siunitx}
% parse-numbers=false : accepte aussi \SI{2,5 \times 10^{-3}}{...}
\sisetup{locale=FR,output-decimal-marker={,},group-minimum-digits=4,parse-numbers=false}
\DeclareSIUnit\ohm{\ensuremath{\symup{\Omega}}} % U+2126 absent de la police
\DeclareSIUnit\division{div} % réglages d'oscilloscope (ms/div, V/div)
\DeclareSIUnit\tr{tr} % tours (vitesse de rotation, ex. \SI{1500}{\tr\per\minute})
\DeclareSIUnit\ch{ch} % cheval-vapeur (puissance moteur, unité usuelle hors SI)
\DeclareSIUnit\franccfa{FCFA} % franc CFA (coûts, contexte camerounais)
\DeclareSIUnit\dioptre{\ensuremath{\delta}} % vergence d'une lentille (dioptrie)
\usepackage{qrcode}
% \degree : commande fréquemment utilisée par les modèles pour un angle ou une
% température en dehors de \SI{}{} (non fournie par les paquets chargés).
\providecommand{\degree}{^{\circ}}
% \qed (fin de démonstration), utilisé par les modèles sans amsthm
\providecommand{\qed}{\hfill$\square$}
% \barre : double barre des valeurs interdites dans un tableau de variations (tkz-tab)
\providecommand{\barre}{\ensuremath{\|}}
% environnement proof (amsthm non chargé) utilisé par les modèles
\ifdefined\proof\else\newenvironment{proof}[1][Démonstration]{\par\noindent\textit{#1.}\ }{\qed\par}\fi
\usepackage{lastpage}
\usepackage{hyperref}
\hypersetup{hidelinks}

% ============================================================
% Palette « Pop » OnBuch+ — mêmes hex que la classe P (plus_ui.dart)
% ============================================================
\definecolor{poporange}{HTML}{F59321}
\definecolor{popdark}{HTML}{E07A0C}
\definecolor{popcream}{HTML}{FFF6E8}
\definecolor{popink}{HTML}{1C1714}
\definecolor{poppurple}{HTML}{7A5AE0}
\definecolor{popgreen}{HTML}{2E9E6A}
\definecolor{poppink}{HTML}{E0457B}
\definecolor{popblue}{HTML}{2F7DE1}
\definecolor{popgold}{HTML}{C98A12}
\definecolor{popmuted}{HTML}{8A7F76}
\definecolor{popline}{HTML}{EADFD2}
\definecolor{popgray}{HTML}{8A7F76}
\colorlet{popgrey}{popgray}
% Alias tolérants (le modèle nomme parfois les couleurs différemment)
\colorlet{popwhite}{white}
\colorlet{popblack}{popink}
\colorlet{popred}{poppink}
\colorlet{popyellow}{popgold}
% Teintes claires pour fonds de tableaux / zones
\colorlet{popblueL}{popblue!10!white}
\colorlet{poporangeL}{poporange!14!white}
\colorlet{popgreenL}{popgreen!12!white}
\colorlet{poppurpleL}{poppurple!12!white}
\colorlet{poppinkL}{poppink!12!white}
\colorlet{popgoldL}{popgold!14!white}

\pagecolor{popcream}

% ============================================================
% Typographie
% ============================================================
\defaultfontfeatures{Ligatures=TeX}
\setmainfont{PlusJakartaSans-Medium.ttf}[Path=../../../fonts/,BoldFont=PlusJakartaSans-Bold.ttf]
% Maths en sans-serif (Fira Math) avec chiffres et lettres droites de Plus
% Jakarta, pour que les formules se fondent dans le texte.
\usepackage[math-style=ISO,bold-style=ISO]{unicode-math}
\setmathfont{FiraMath-Regular.otf}[Path=../../../fonts/]
\setmathfont{PlusJakartaSans-Medium.ttf}[Path=../../../fonts/,range={up/num,up/{Latin,latin},\mathup}]
\usepackage{icomma} % virgule décimale sans espace en mode math
% Caractères mathématiques Unicode absents des polices de texte (Plus Jakarta,
% Archivo Black) : rendus par leur équivalent mathématique, en texte comme en
% formule — sinon ils s'affichent en carré vide (ex. « Arithmétique dans ℤ »).
\usepackage{newunicodechar}
\newunicodechar{ℝ}{\ensuremath{\mathbb{R}}}
\newunicodechar{ℤ}{\ensuremath{\mathbb{Z}}}
\newunicodechar{ℂ}{\ensuremath{\mathbb{C}}}
\newunicodechar{ℕ}{\ensuremath{\mathbb{N}}}
\newunicodechar{ℚ}{\ensuremath{\mathbb{Q}}}
\newunicodechar{𝔻}{\ensuremath{\mathbb{D}}}
\newunicodechar{α}{\ensuremath{\alpha}}
\newunicodechar{β}{\ensuremath{\beta}}
\newunicodechar{γ}{\ensuremath{\gamma}}
\newunicodechar{δ}{\ensuremath{\delta}}
\newunicodechar{ε}{\ensuremath{\varepsilon}}
\newunicodechar{θ}{\ensuremath{\theta}}
\newunicodechar{λ}{\ensuremath{\lambda}}
\newunicodechar{μ}{\ensuremath{\mu}}
\newunicodechar{σ}{\ensuremath{\sigma}}
\newunicodechar{τ}{\ensuremath{\tau}}
\newunicodechar{φ}{\ensuremath{\varphi}}
\newunicodechar{ψ}{\ensuremath{\psi}}
\newunicodechar{ω}{\ensuremath{\omega}}
\newunicodechar{Σ}{\ensuremath{\Sigma}}
% majuscules produites par \MakeUppercase dans les titres (α-aminés -> Α)
\newunicodechar{Α}{\ensuremath{\alpha}}
\newunicodechar{Β}{\ensuremath{\beta}}
\newunicodechar{Γ}{\ensuremath{\Gamma}}
\newunicodechar{Δ}{\ensuremath{\Delta}}
\newunicodechar{Ε}{\ensuremath{\varepsilon}}
\newunicodechar{Θ}{\ensuremath{\Theta}}
\newunicodechar{Λ}{\ensuremath{\Lambda}}
\newunicodechar{Μ}{\ensuremath{\mu}}
\newunicodechar{Τ}{\ensuremath{\tau}}
\newunicodechar{Φ}{\ensuremath{\Phi}}
\newunicodechar{Ψ}{\ensuremath{\Psi}}
\newunicodechar{Ω}{\ensuremath{\Omega}}
\newunicodechar{↦}{\ensuremath{\mapsto}}
\newunicodechar{∈}{\ensuremath{\in}}
\newunicodechar{∉}{\ensuremath{\notin}}
\newunicodechar{∧}{\ensuremath{\wedge}}
\newunicodechar{⇔}{\ensuremath{\Leftrightarrow}}
\newunicodechar{⟺}{\ensuremath{\Longleftrightarrow}}
\newunicodechar{⇒}{\ensuremath{\Rightarrow}}
\newunicodechar{⟹}{\ensuremath{\Longrightarrow}}
\newunicodechar{∘}{\ensuremath{\circ}}
\newunicodechar{∪}{\ensuremath{\cup}}
\newunicodechar{∩}{\ensuremath{\cap}}
\newunicodechar{≡}{\ensuremath{\equiv}}
\newunicodechar{⊂}{\ensuremath{\subset}}
\newunicodechar{⊥}{\ensuremath{\perp}}
\newunicodechar{∃}{\ensuremath{\exists}}
\newunicodechar{∀}{\ensuremath{\forall}}
\newunicodechar{∛}{\ensuremath{\sqrt[3]{\,}}}
\newunicodechar{∜}{\ensuremath{\sqrt[4]{\,}}}
\newunicodechar{⃗}{\ensuremath{{}^{\to}}}
\newunicodechar{ⁿ}{\ifmmode^{n}\else\textsuperscript{\ensuremath{n}}\fi}
\newunicodechar{ˣ}{\ifmmode^{x}\else\textsuperscript{\ensuremath{x}}\fi}
\newunicodechar{ᵇ}{\ifmmode^{b}\else\textsuperscript{\ensuremath{b}}\fi}
\newunicodechar{ʳ}{\ifmmode^{r}\else\textsuperscript{\ensuremath{r}}\fi}
\newunicodechar{ᵘ}{\ifmmode^{u}\else\textsuperscript{\ensuremath{u}}\fi}
\newunicodechar{ʸ}{\ifmmode^{y}\else\textsuperscript{\ensuremath{y}}\fi}
\newunicodechar{ᵃ}{\ifmmode^{a}\else\textsuperscript{\ensuremath{a}}\fi}
\newunicodechar{ᵉ}{\ifmmode^{e}\else\textsuperscript{\ensuremath{e}}\fi}
\newunicodechar{ᵏ}{\ifmmode^{k}\else\textsuperscript{\ensuremath{k}}\fi}
\newunicodechar{ᵖ}{\ifmmode^{p}\else\textsuperscript{\ensuremath{p}}\fi}
\newunicodechar{ᴺ}{\ifmmode^{N}\else\textsuperscript{\ensuremath{N}}\fi}
\newunicodechar{ᶜ}{\ifmmode^{c}\else\textsuperscript{\ensuremath{c}}\fi}
\newunicodechar{ʰ}{\ifmmode^{h}\else\textsuperscript{\ensuremath{h}}\fi}
\newunicodechar{ᵠ}{\ifmmode^{q}\else\textsuperscript{\ensuremath{q}}\fi}
\newunicodechar{ₙ}{\ifmmode_{n}\else\textsubscript{\ensuremath{n}}\fi}
\newunicodechar{ₐ}{\ifmmode_{a}\else\textsubscript{\ensuremath{a}}\fi}
\newunicodechar{ᵢ}{\ifmmode_{i}\else\textsubscript{\ensuremath{i}}\fi}
\newunicodechar{ᵣ}{\ifmmode_{r}\else\textsubscript{\ensuremath{r}}\fi}
\newunicodechar{ₖ}{\ifmmode_{k}\else\textsubscript{\ensuremath{k}}\fi}
\newunicodechar{ᵦ}{\ifmmode_{\beta}\else\textsubscript{\ensuremath{\beta}}\fi}
\newunicodechar{ₓ}{\ifmmode_{x}\else\textsubscript{\ensuremath{x}}\fi}
\newfontfamily\popdisplay{ArchivoBlack-Regular.ttf}[Path=../../../fonts/]
\newfontfamily\poplogo{SpaceGrotesk-Bold.ttf}[Path=../../../fonts/]
\newfontfamily\popsemi{PlusJakartaSans-SemiBold.ttf}[Path=../../../fonts/]
\newfontfamily\popxbold{PlusJakartaSans-ExtraBold.ttf}[Path=../../../fonts/]
\newfontfamily\popxboldls{PlusJakartaSans-ExtraBold.ttf}[Path=../../../fonts/,LetterSpace=3.0]

\setlist[enumerate]{leftmargin=1.7em,itemsep=5pt,topsep=4pt}
\setlist[itemize]{leftmargin=1.7em,itemsep=4pt,topsep=4pt,label={\color{poporange}\textbullet}}
\setlength{\parindent}{0pt}
\setlength{\parskip}{5pt}
\renewcommand{\arraystretch}{1.25}

% pgfplots : style Pop par défaut
\pgfplotsset{
  every axis/.append style={axis line style={popink,line width=1pt},tick style={popink},
    grid style={popline,line width=0.5pt},label style={font=\small},tick label style={font=\footnotesize}},
  popcurve/.style={poporange,line width=1.8pt,no markers,smooth},
}
% Même style disponible dans un \draw TikZ simple (hors pgfplots)
\tikzset{popcurve/.style={poporange,line width=1.8pt}}
\tikzset{popgrid/.style={popline,very thin}}
\colorlet{popcurve}{poporange} % le nom du style est parfois utilisé comme couleur

% ============================================================
% Pill / squiggle
% ============================================================
\newcommand{\poppill}[3]{%
  \tikz[baseline=-0.55ex]\node[draw=popink,line width=1.2pt,rounded corners=2.6pt,%
    fill=#2,inner xsep=6.5pt,inner ysep=3pt]{%
    \popxboldls\fontsize{7}{7}\selectfont\color{#3}\MakeUppercase{#1}};%
}
\newcommand{\popsquiggle}[1]{%
  \tikz[baseline=-0.4ex]{
    \draw[#1,line width=2.6pt,line cap=round]
      plot [smooth] coordinates {(0,0.09) (0.34,0.01) (0.68,0.21) (1.02,0.01) (1.36,0.10)};
  }%
}
% Mot-clé surligné (marqueur orange sous le texte)
\newcommand{\cle}[1]{{\popxbold\color{popdark}#1}}

% ============================================================
% En-tête / pied
% ============================================================
\pagestyle{fancy}
\fancyhf{}
\renewcommand{\headrulewidth}{0pt}
\renewcommand{\footrulewidth}{0pt}
\lhead{\raisebox{-0.15ex}{\poplogo\fontsize{13}{13}\selectfont\color{popink}OnBuch\color{poporange}+}}
\rhead{\poppill{\DOCMATIERE\ \textbullet\ \DOCNIVEAU\ \textbullet\ Leçon \DOCLECON}{white}{popink}}
\lfoot{\popxbold\fontsize{7.6}{7.6}\selectfont\color{popmuted}\MakeUppercase{Cours signé OnBuch+}\ \textbullet\ {\popsemi\DOCTITRE}}
\rfoot{\tikz[baseline=-0.6ex]\node[rounded corners=8.5pt,fill=popink,minimum height=17pt,minimum width=17pt,inner xsep=5pt,inner sep=0]{\popxbold\fontsize{8}{8}\selectfont\color{popcream}\thepage\,/\,\pageref*{LastPage}};}

% ============================================================
% Titres
% ============================================================
\newcommand{\chaptitre}[2]{%
  \begin{tcolorbox}[enhanced,colback=poporange,colframe=popink,boxrule=2.6pt,arc=11pt,
      drop shadow=popink,left=15pt,right=15pt,top=12pt,bottom=11pt,before skip=3pt,after skip=18pt]
    \poppill{#2}{popink}{popcream}\par\vspace{9pt}
    {\raggedright\hyphenpenalty=10000\exhyphenpenalty=10000\popdisplay\fontsize{22}{24}\selectfont\color{popcream}\MakeUppercase{#1}\par}
  \end{tcolorbox}
}
% Section numérotée : pastille orange avec le numéro + titre Archivo
\newcounter{popsec}
\newcommand{\coursec}[1]{%
  \stepcounter{popsec}\setcounter{popsub}{0}%
  \par\vspace{16pt}\noindent
  \begin{minipage}{\linewidth}
  \tikz[baseline=-0.7ex]\node[draw=popink,line width=1.6pt,fill=poporange,rounded corners=4pt,minimum size=22pt,inner sep=0]{\popdisplay\fontsize{12}{12}\selectfont\color{popcream}\thepopsec};%
  \hspace{8pt}{\popdisplay\fontsize{15}{17}\selectfont\color{popink}\MakeUppercase{#1}}\par
  \vspace{4pt}\noindent{\color{poporange}\rule{36pt}{3.6pt}}
  \end{minipage}\par\nopagebreak\vspace{8pt}
}
\newcounter{popsub}
\newcommand{\courssub}[1]{%
  \stepcounter{popsub}%
  \par\vspace{9pt}\noindent{\popxbold\fontsize{11.5}{13}\selectfont\color{popink}{\color{poporange}\thepopsec.\thepopsub}\hspace{6pt}#1}\par\nopagebreak\vspace{4pt}
}

% ============================================================
% Boîtes pédagogiques — même recette : fond blanc, bordure épaisse colorée,
% ombre dure encre, badge plein. Titre optionnel [#1].
% ============================================================
\tcbset{popbase/.style={breakable,enhanced,colback=white,boxrule=2.3pt,arc=9pt,
  shadow={3pt}{-3pt}{0pt}{popink},left=14pt,right=14pt,top=11pt,bottom=11pt,before skip=11pt,after skip=15pt,
  fontupper=\popsemi\fontsize{10.6}{14.5}\selectfont\color{popink},
  fontlower=\popsemi\fontsize{10.6}{14.5}\selectfont\color{popink}}}
% Coche dessinée (le glyphe ✓ n'existe pas dans les polices du document)
\newcommand{\popcheck}[1]{\tikz[baseline=-0.5ex]\draw[#1,line width=1.6pt,line cap=round,line join=round](-0.13,0.0)--(-0.04,-0.09)--(0.13,0.1);}
\providecommand{\coche}{\popcheck{popgreen}}
\providecommand{\resultat}[1]{\par\smallskip\noindent\fcolorbox{popgreen}{popgreenL}{\parbox{\dimexpr\linewidth-2\fboxsep-2\fboxrule\relax}{\textbf{Résultat.} #1}}\par\smallskip}
\newcommand{\popboxhead}[3]{\poppill{#1}{#2}{white}\ifx\relax#3\relax\else\hspace{6pt}{\popxbold\fontsize{10.6}{12}\selectfont\color{popink}#3}\fi\par\vspace{7pt}}

\newtcolorbox{aretenir}[1][]{popbase,colframe=popgreen,before upper={\popboxhead{À retenir}{popgreen}{#1}}}
\newtcolorbox{exemplebox}[1][]{popbase,colframe=poporange,before upper={\popboxhead{Exemple}{poporange}{#1}}}
\newtcolorbox{definition}[1][]{popbase,colframe=popblue,before upper={\popboxhead{Définition}{popblue}{#1}}}
\newtcolorbox{propriete}[1][]{popbase,colframe=poppurple,before upper={\popboxhead{Propriété}{poppurple}{#1}}}
\newtcolorbox{methode}[1][]{popbase,colframe=popgold,before upper={\popboxhead{Méthode}{popgold}{#1}}}
\newtcolorbox{attention}[1][]{popbase,colframe=poppink,before upper={\popboxhead{Attention}{poppink}{#1}}}
\newtcolorbox{experience}[1][]{popbase,colframe=popblue,colback=popblueL,before upper={\popboxhead{Expérience}{popblue}{#1}}}
% « Le savais-tu ? » : carte violette pleine (contexte, culture, Cameroun)
\newtcolorbox{savaistu}[1][]{popbase,colback=poppurple,colframe=popink,
  fontupper=\popsemi\fontsize{10.4}{14.2}\selectfont\color{white},
  before upper={\poppill{Le savais-tu\,?}{white}{poppurple}\ifx\relax#1\relax\else\hspace{6pt}{\popxbold\color{white}#1}\fi\par\vspace{7pt}}}
% Exercice résolu : énoncé en haut, \tcblower puis la solution sur fond crème
\newcounter{popexo}\newcounter{popexr}
\newtcolorbox{exoresolu}[1][]{popbase,colframe=poporange,colbacklower=popcream,
  segmentation style={popink,line width=1.2pt,dashed},
  before upper={\stepcounter{popexr}\popboxhead{Exercice résolu \thepopexr}{poporange}{#1}},
  before lower={\poppill{Solution}{popink}{popcream}\par\vspace{6pt}}}
% Exercice (non résolu en place) — la correction est dans la partie Corrigés
\newtcolorbox{exercice}[1][]{popbase,colframe=popink,
  before upper={\stepcounter{popexo}\popboxhead{Exercice \thepopexo}{popink}{#1}}}
% Corrigé
\newtcolorbox{corrige}[1][]{popbase,colframe=popgreen,colback=popgreenL,
  before upper={\popboxhead{Corrigé}{popgreen}{#1}}}
% Objectifs / prérequis / bilan
\newtcolorbox{objectifs}{popbase,colframe=popink,colback=poporangeL,
  before upper={\popboxhead{Objectifs de la leçon}{popink}{}}}
\newtcolorbox{prerequis}{popbase,colframe=popink,colback=popblueL,
  before upper={\popboxhead{Prérequis}{popblue}{}}}
\newtcolorbox{bilan}{popbase,colframe=popink,colback=white,boxrule=2.8pt,
  before upper={\popboxhead{Fiche bilan}{poporange}{L'essentiel à connaître pour l'examen}}}
% Figure encadrée avec légende
\newtcolorbox{popfigure}[1][]{enhanced,breakable=false,colback=white,colframe=popline,boxrule=1.4pt,arc=8pt,
  left=10pt,right=10pt,top=10pt,bottom=8pt,before skip=10pt,after skip=12pt,halign=center,
  lower separated=false,halign lower=center,
  fontlower=\popsemi\fontsize{9}{11.5}\selectfont\color{popmuted},
  #1}
\newcommand{\legende}[1]{\par\vspace{6pt}{\popsemi\fontsize{9}{11.5}\selectfont\color{popmuted}#1}}

% ============================================================
% Signature OnBuch+
% ============================================================
\newcommand{\poplogoinline}[1]{{\poplogo\fontsize{#1}{#1}\selectfont\color{popink}OnBuch\color{poporange}+}}
\newcommand{\signatureonbuch}{%
  \begin{tcolorbox}[enhanced,colback=popink,colframe=popink,boxrule=2.2pt,arc=10pt,
      drop shadow=poporange,left=14pt,right=14pt,top=10pt,bottom=10pt,before skip=0pt,after skip=0pt]
    \tikz[baseline=-0.7ex]\node[circle,fill=popgreen,draw=popcream,line width=1.3pt,minimum size=22pt,inner sep=0]{\popcheck{white}};%
    \hspace{9pt}\begin{minipage}[c]{0.62\linewidth}
      {\popxbold\fontsize{12}{13}\selectfont\color{popcream}Signé\ \textbullet\ {\poplogo OnBuch\color{poporange}+}}\par\vspace{2pt}
      {\raggedright\popsemi\fontsize{8}{10.5}\selectfont\color{popline}\MakeUppercase{Équipe pédagogique OnBuch+}\\\MakeUppercase{Conforme au programme officiel MINESEC}\par}
    \end{minipage}\hfill
    {\poplogo\fontsize{22}{22}\selectfont\color{popcream}OnBuch\color{poporange}+}
  \end{tcolorbox}%
}

% ============================================================
% Page de garde
% #1 titre · #2 matière · #3 niveau · #4 module · #5 n° leçon · #6 durée ·
% #7 description / objectifs (texte ou itemize)
% ============================================================
\newcommand{\pagegarde}[7]{%
  \thispagestyle{empty}
  \newgeometry{a4paper,margin=1.7cm,top=1.6cm,bottom=1.4cm}%
  \begin{tikzpicture}[remember picture,overlay]
    \fill[poporange!18] ([xshift=-3.2cm,yshift=-3.6cm]current page.north east) circle (4.6cm);
    \fill[poppurple!14] ([xshift=2.4cm,yshift=7.6cm]current page.south west) circle (3.8cm);
  \end{tikzpicture}%
  {\poplogo\fontsize{30}{30}\selectfont\color{popink}OnBuch\color{poporange}+}\hfill
  \raisebox{0.6ex}{\poppill{Cours signé \textbullet\ Premium}{popink}{popcream}}\par
  \vspace{1pt}\popsquiggle{poppurple}\par
  \vspace{8pt}
  \begin{tcolorbox}[enhanced,colback=poporange,colframe=popink,boxrule=2.8pt,arc=13pt,
      drop shadow=popink,left=18pt,right=18pt,top=12pt,bottom=13pt,after skip=10pt]
    \poppill{#2 \textbullet\ #3}{popink}{popcream}\hspace{5pt}\poppill{#4}{white}{popink}\par\vspace{8pt}
    {\popdisplay\fontsize{14}{14}\selectfont\color{popink}LEÇON #5}\par\vspace{4pt}
    {\raggedright\hyphenpenalty=10000\exhyphenpenalty=10000\popdisplay\fontsize{25}{27}\selectfont\color{popcream}\MakeUppercase{#1}\par}
    \vspace{8pt}
    {\popxbold\fontsize{9.5}{9.5}\selectfont\color{popink}\MakeUppercase{Durée conseillée : #6}}
  \end{tcolorbox}
  \begin{tcolorbox}[enhanced,colback=white,colframe=popink,boxrule=2.2pt,arc=10pt,
      drop shadow=popink,left=16pt,right=16pt,top=10pt,bottom=10pt,after skip=8pt]
    \poppill{Dans cette leçon}{poppurple}{white}\par\vspace{5pt}
    \popsemi\fontsize{9.3}{12.6}\selectfont\color{popink}#7
  \end{tcolorbox}
  \noindent\begin{minipage}[t]{0.487\linewidth}
    \begin{tcolorbox}[enhanced,colback=popgreen,colframe=popink,boxrule=2.2pt,arc=10pt,
        drop shadow=popink,left=11pt,right=11pt,top=10pt,bottom=10pt,height=3.1cm,valign=center]
      \tcbox[colback=white,colframe=popink,boxrule=1.3pt,arc=5pt,boxsep=0pt,nobeforeafter,
        left=3pt,right=3pt,top=3pt,bottom=3pt,valign=center]{\qrcode[height=1.9cm]{https://play.google.com/store/apps/details?id=com.luvvix.onbuch&hl=fr}}%
      \hspace{8pt}\begin{minipage}[c]{0.5\linewidth}
        {\popdisplay\fontsize{11}{12.5}\selectfont\color{white}\MakeUppercase{Télécharge OnBuch}}\par\vspace{4pt}
        {\raggedright\popsemi\fontsize{8.4}{11}\selectfont\color{white}Gratuit \textbullet\ Google Play\\Tuteur IA, annales, concours, résultats\par}
      \end{minipage}
    \end{tcolorbox}
  \end{minipage}\hfill
  \begin{minipage}[t]{0.487\linewidth}
    \begin{tcolorbox}[enhanced,colback=poppurple,colframe=popink,boxrule=2.2pt,arc=10pt,
        drop shadow=popink,left=11pt,right=11pt,top=10pt,bottom=10pt,height=3.1cm,valign=center]
      \tikz[baseline=-0.7ex]\node[circle,fill=white,draw=popink,line width=1.3pt,minimum size=1.6cm,inner sep=0]{\popdisplay\fontsize{24}{24}\selectfont\color{poppurple}+};%
      \hspace{8pt}\begin{minipage}[c]{0.6\linewidth}
        {\popdisplay\fontsize{11}{12.5}\selectfont\color{white}\MakeUppercase{Passe à OnBuch+}}\par\vspace{4pt}
        {\raggedright\popsemi\fontsize{8.4}{11}\selectfont\color{white}Tous les cours signés, Tuteur IA illimité, fiches PDF — onglet OnBuch+ de l'app\par}
      \end{minipage}
    \end{tcolorbox}
  \end{minipage}
  \vspace{8pt}
  \popcontenu
  \vfill
  \signatureonbuch
  \restoregeometry
  \newpage
}

% Rangée « Ce cours contient » (page de garde)
\newcommand{\poptile}[3]{% #1 couleur #2 gros texte #3 légende
  \begin{tcolorbox}[enhanced,colback=#1,colframe=popink,boxrule=2pt,arc=9pt,shadow={2.5pt}{-2.5pt}{0pt}{popink},
      left=6pt,right=6pt,top=9pt,bottom=9pt,halign=center,valign=center,height=1.75cm,width=\linewidth,nobeforeafter]
    {\popdisplay\fontsize{12.5}{13}\selectfont\color{white}\MakeUppercase{#2}}\par\vspace{3pt}
    {\centering\popsemi\fontsize{7.8}{9.5}\selectfont\color{white}#3\par}
  \end{tcolorbox}}
\newcommand{\popcontenu}{%
  \par\noindent{\popxboldls\fontsize{7.5}{7.5}\selectfont\color{popmuted}CE COURS SIGNÉ CONTIENT}\par\vspace{7pt}
  \noindent\begin{minipage}[t]{0.235\linewidth}\poptile{popblue}{Cours}{complet, illustré, pas à pas}\end{minipage}\hfill
  \begin{minipage}[t]{0.235\linewidth}\poptile{poporange}{Méthodes}{et exercices résolus}\end{minipage}\hfill
  \begin{minipage}[t]{0.235\linewidth}\poptile{poppink}{Exercices}{type examen + corrigés}\end{minipage}\hfill
  \begin{minipage}[t]{0.235\linewidth}\poptile{popgreen}{Fiche bilan}{l'essentiel à retenir}\end{minipage}\par}

% Sommaire
\newtcolorbox{sommaire}{enhanced,colback=white,colframe=popink,boxrule=2.2pt,arc=10pt,
  drop shadow=popink,left=18pt,right=18pt,top=13pt,bottom=13pt,before skip=6pt,after skip=20pt,
  before upper={\poppill{Sommaire}{poppurple}{white}\par\vspace{9pt}\popsemi\fontsize{10.6}{14.5}\selectfont\color{popink}}}

% Signature de fin de cours — poussée tout en bas de la dernière page
\newcommand{\findecours}{%
  \par\vfill\nopagebreak
  \begin{center}\popsquiggle{poporange}\end{center}\vspace{4pt}
  \signatureonbuch
}
\AtBeginDocument{\def\llbracket{\mathopen{[\mkern-3.5mu[}}\def\rrbracket{\mathclose{]\mkern-3.5mu]}}} % intervalles d'entiers (glyphe ⟦ absent de la police)
% Glyphes absents de Fira Math : redéfinis à partir de glyphes disponibles
\AtBeginDocument{%
  \def\setminus{\mathbin{\backslash}}\let\smallsetminus\setminus
  \def\vdots{\mathord{\vbox{\baselineskip3.2pt\lineskiplimit0pt\kern4pt\hbox{.}\hbox{.}\hbox{.}}}}%
  \def\ddots{\mathinner{\mkern1mu\raise6.5pt\hbox{.}\mkern2mu\raise3.6pt\hbox{.}\mkern2mu\raise0.7pt\hbox{.}\mkern1mu}}%
  \def\triangle{\mathord{\scalebox{0.9}{$\symup{\Delta}$}}}%
  \def\nabla{\mathord{\raisebox{\depth}{\scalebox{1}[-1]{$\symup{\Delta}$}}}}%
  \def\checkmark{\popcheck{popdark}}%
  \def\star{\mathbin{\ast}}\def\bigstar{\mathord{\ast}}%
  \def\diamond{\mathbin{\scalebox{0.6}{$\Diamond$}}}%
  \def\bigcup{\mathop{\mathchoice{\raisebox{-0.3em}{\scalebox{1.7}{$\cup$}}}{\scalebox{1.25}{$\cup$}}{\cup}{\cup}}\displaylimits}%
  \def\bigcap{\mathop{\mathchoice{\raisebox{-0.3em}{\scalebox{1.7}{$\cap$}}}{\scalebox{1.25}{$\cap$}}{\cap}{\cap}}\displaylimits}%
  \def\lesseqqgtr{\mathrel{\lessgtr}}\def\bigoplus{\mathop{\oplus}\displaylimits}%
}
\providecommand{\ssi}{si et seulement si} % abréviation fréquente
\AtBeginDocument{\providecommand{\wideparen}{\overparen}} % arc de cercle

%% FIGFIX-BEGIN v3 — correctifs globaux des figures (docs/onbuch-generation/tools/figfix)
\usepackage{adjustbox}\usepackage{etoolbox}
% toute figure TikZ/pgfplots plus large que la ligne est réduite à la largeur disponible
\BeforeBeginEnvironment{tikzpicture}{\begin{adjustbox}{max width=\linewidth}}
\AfterEndEnvironment{tikzpicture}{\end{adjustbox}}
% légendes pgfplots lisibles par-dessus les courbes
\pgfplotsset{every axis/.append style={legend style={fill=white,fill opacity=0.92,text opacity=1,draw=black!15,inner sep=3pt}}}
% étiquettes d'arêtes (syntaxe edge["..."]) avec fond blanc
\tikzset{every edge quotes/.append style={fill=white,inner sep=1.2pt}}
% bulles de cartes mentales : texte à la ligne dans la bulle, bulle un peu plus grande
\tikzset{every concept/.append style={text width=2.0cm,align=center,minimum size=2.3cm,inner sep=2pt}}
%% FIGFIX-END
\begin{document}

\pagegarde{Sensibilisation dans le cadre des catastrophes liées aux mouvements atmosphérique et océanique}{SVT}{1ère D}{Module 3 : L'Éducation à l'Environnement et au Développement Durable}{ND12}{8 h}{Cette leçon explique les mécanismes de la circulation atmosphérique et océanique à l'échelle du globe, puis montre comment ces mouvements engendrent des phénomènes climatiques dangereux (tempêtes, cyclones, inondations). Elle développe enfin une culture de prévention et de protection face aux catastrophes liées à ces mouvements, dans le contexte camerounais et mondial.
\vspace{4pt}
\begin{multicols}{2}\small
\begin{itemize}
\item À la fin de cette leçon, je sais décrire sur un schéma les grandes cellules de convection atmosphérique et les zones de pression associées.
\item À la fin de cette leçon, je sais expliquer l'origine des vents planétaires (alizés, moussons, harmattan) et le rôle de la force de Coriolis.
\item À la fin de cette leçon, je sais décrire sur une carte les principaux courants marins et distinguer courants chauds et courants froids.
\item À la fin de cette leçon, je sais relier les interactions océan-atmosphère à des phénomènes climatiques (mousson, El Niño).
\item À la fin de cette leçon, je sais relier un phénomène météorologique dangereux (cyclone, tempête, inondation) à son origine dynamique.
\item À la fin de cette leçon, je sais identifier les conditions de formation d'un cyclone tropical et expliquer pourquoi le Cameroun en est épargné.
\end{itemize}\end{multicols}}

\begin{sommaire}
\begin{multicols}{2}
\begin{enumerate}
\item La circulation générale de l'atmosphère
\item La circulation océanique : les courants marins
\item Interactions océan-atmosphère
\item Les phénomènes climatiques dangereux
\item Prévention et protection face aux catastrophes
\item Activité d'intégration
\item Méthodes et exercices résolus
\item Exercices
\item Corrigés des exercices
\item Fiche bilan
\end{enumerate}
\end{multicols}
\end{sommaire}

\begin{objectifs}
\begin{itemize}
\item À la fin de cette leçon, je sais décrire sur un schéma les grandes cellules de convection atmosphérique et les zones de pression associées.
\item À la fin de cette leçon, je sais expliquer l'origine des vents planétaires (alizés, moussons, harmattan) et le rôle de la force de Coriolis.
\item À la fin de cette leçon, je sais décrire sur une carte les principaux courants marins et distinguer courants chauds et courants froids.
\item À la fin de cette leçon, je sais relier les interactions océan-atmosphère à des phénomènes climatiques (mousson, El Niño).
\item À la fin de cette leçon, je sais relier un phénomène météorologique dangereux (cyclone, tempête, inondation) à son origine dynamique.
\item À la fin de cette leçon, je sais identifier les conditions de formation d'un cyclone tropical et expliquer pourquoi le Cameroun en est épargné.
\item À la fin de cette leçon, je sais proposer des mesures de prévention et de protection adaptées face aux catastrophes liées aux mouvements atmosphériques et océaniques.
\end{itemize}
\end{objectifs}

\begin{prerequis}
\begin{itemize}
\item Notions de pression atmosphérique et de température (physique de Seconde / géographie de Première)
\item Répartition inégale de l'énergie solaire à la surface du globe (climat et milieux, classe de Première)
\item Rotation de la Terre et ses conséquences (jour/nuit, fuseaux horaires)
\item Structure verticale de l'atmosphère (troposphère) et cycle de l'eau (évaporation, condensation, précipitations)
\end{itemize}
\end{prerequis}

\newpage

\coursec{La circulation générale de l'atmosphère}

En octobre 2019, des pluies diluviennes provoquent des glissements de terrain meurtriers à Bafoussam, ensevelissant des quartiers entiers sous la boue. Chaque année, entre décembre et février, les tempêtes de sable du harmattan couvrent Garoua d'une poussière ocre qui réduit la visibilité et perturbe le trafic aérien. À Yaoundé comme à Douala, les inondations frappent régulièrement les quartiers bas lors des fortes pluies. D'où viennent ces vents et ces pluies extrêmes ? Pourquoi certaines régions du monde subissent-elles des cyclones dévastateurs alors que le littoral camerounais en est épargné ? Pour comprendre ces catastrophes et apprendre à s'en protéger, il faut d'abord comprendre comment l'air se déplace à l'échelle de la planète : c'est l'objet de cette première partie, consacrée à la \cle{circulation générale de l'atmosphère}.

\courssub{Le moteur de la circulation : l'inégale répartition de l'énergie solaire}

\textbf{Observation.} Place la main au-dessus d'une lampe allumée : tu sens la chaleur monter. À l'échelle de la Terre, le Soleil joue le rôle de la lampe, mais il ne chauffe pas toutes les régions de la même façon. À l'équateur, les rayons arrivent presque verticalement : ils concentrent leur énergie sur une petite surface. Aux pôles, les mêmes rayons arrivent très inclinés : leur énergie se répartit sur une grande surface et traverse une couche d'atmosphère plus épaisse. Résultat mesuré par les satellites : la zone équatoriale reçoit en moyenne plus d'énergie qu'elle n'en réémet vers l'espace (\cle{surplus énergétique}), tandis que les régions polaires perdent plus d'énergie qu'elles n'en reçoivent (\cle{déficit énergétique}).

\textbf{Hypothèse.} Si l'équateur accumulait indéfiniment cette chaleur, il deviendrait invivable et les pôles se refroidiraient sans fin. Or les températures moyennes restent stables d'année en année. Il doit donc exister un mécanisme de \emph{transport de chaleur} de l'équateur vers les pôles.

\textbf{Interprétation et conclusion.} Ce mécanisme, c'est le mouvement de l'air (et, nous le verrons plus tard, celui des océans). L'air chaud de l'équateur s'élève, se déplace en altitude vers les pôles, tandis que l'air froid des pôles revient en surface vers l'équateur. La circulation atmosphérique est donc un gigantesque système de redistribution thermique.

\begin{definition}[Circulation générale de l'atmosphère]
La \cle{circulation générale de l'atmosphère} est l'ensemble des mouvements d'air à grande échelle (continentale et planétaire) qui redistribuent la chaleur de l'équateur, où elle est en excès, vers les pôles, où elle est en déficit. Elle résulte de l'inégale répartition de l'énergie solaire à la surface du globe et de la rotation de la Terre.
\end{definition}

\courssub{La pression atmosphérique : dépressions et anticyclones}

\textbf{Situation concrète.} Au-dessus d'une marmite d'eau bouillante, l'air chaud s'élève ; dans une chambre froide, l'air glacial « tombe » au sol. Ces deux observations quotidiennes illustrent le principe fondamental : l'\cle{air chaud} se dilate, devient \emph{moins dense} que l'air environnant et s'élève ; l'\cle{air froid} se contracte, devient \emph{plus dense} et descend.

Lorsque l'air s'élève en un lieu, il « allège » la colonne d'air qui pèse sur le sol : la pression atmosphérique y diminue, on parle de \cle{basse pression} ou \cle{dépression}. Lorsque l'air descend, il « alourdit » la colonne : la pression augmente, on parle de \cle{haute pression} ou \cle{anticyclone}. En surface, l'air s'écoule toujours des hautes pressions vers les basses pressions, comme l'eau coule du haut vers le bas d'une pente : c'est le \cle{vent}.

\begin{attention}[Erreur fréquente : pourquoi l'air chaud monte-t-il ?]
L'air chaud ne « monte » pas parce qu'il serait \emph{attiré vers le haut} ou parce que « la chaleur monte » par magie. Il s'élève parce qu'il est \textbf{moins dense} que l'air froid qui l'entoure : celui-ci, plus lourd, s'engouffre dessous et le pousse vers le haut. C'est la \cle{poussée d'Archimède}, exactement comme une bulle d'air qui remonte dans l'eau ou un ballon gonflé à l'hélium qui s'élève dans le ciel.
\end{attention}

\begin{experience}[Mettre en évidence la convection de l'air]
\textbf{Matériel :} une bougie allumée, un serpentin de papier léger suspendu à un fil, une boîte à deux cheminées (ou une boîte percée de deux trous avec deux tubes).

\textbf{Protocole :} 1) Suspendre le serpentin de papier au-dessus de la flamme, sans la toucher. 2) Placer la bougie allumée sous l'une des cheminées de la boîte et approcher un bâtonnet d'encens fumant de l'autre cheminée.

\textbf{Observations :} le serpentin tourne sous l'effet de l'air chaud ascendant ; la fumée de l'encens est aspirée dans la cheminée froide et ressort par la cheminée chaude.

\textbf{Interprétation :} l'air chauffé par la flamme devient moins dense et s'élève (ascendance) ; l'air froid environnant s'engouffre pour le remplacer (subsidence et vent de surface). On a créé en miniature une \emph{cellule de convection}, modèle réduit de la circulation atmosphérique planétaire.
\end{experience}

\courssub{Les trois cellules de convection par hémisphère}

\begin{definition}[Cellule de convection]
Une \cle{cellule de convection} est une boucle fermée de mouvement de l'air comprenant quatre étapes : une \textbf{ascendance} d'air chaud en un lieu, un \textbf{déplacement en altitude} vers des latitudes plus froides, une \textbf{subsidence} (descente) d'air refroidi, et un \textbf{retour en surface} vers le point de départ.
\end{definition}

Si la Terre ne tournait pas, une seule cellule géante relierait l'équateur au pôle dans chaque hémisphère. Mais la rotation terrestre fragmente cette circulation : on observe en réalité \textbf{trois cellules de convection} par hémisphère.

\begin{enumerate}
\item \textbf{La cellule de Hadley (0°--30° de latitude).} C'est la plus puissante et la plus directement thermique. À l'équateur, l'air très chauffé s'élève (ascendance), s'écoule en altitude vers les pôles, se refroidit et redescend vers 30° de latitude (subsidence), puis revient en surface vers l'équateur. C'est le retour de surface de cette cellule qui alimente les alizés.
\item \textbf{La cellule de Ferrel (30°--60° de latitude).} C'est une cellule \emph{indirecte}, entraînée par les deux autres comme un engrenage intermédiaire. En surface, l'air s'y déplace des 30° vers les 60° de latitude ; en altitude, le flux est inverse.
\item \textbf{La cellule polaire (60°--90° de latitude).} Au pôle, l'air très froid et dense descend (subsidence), s'écoule en surface vers les latitudes plus basses jusqu'à environ 60°, où il rencontre l'air plus doux venu du sud et s'élève (ascendance), avant de revenir en altitude vers le pôle.
\end{enumerate}

\begin{popfigure}
\begin{center}
\begin{tikzpicture}[scale=1.0, font=\small]
% Sol
\draw[thick, popdark] (0,0) -- (12,0);
\fill[popgoldL] (0,0) rectangle (12,-0.25);
% Latitudes
\foreach \x/\lat in {0/Équateur 0°, 3/30°N, 6/60°N, 9/Pôle 90°N}{
  \draw[dashed, popmuted] (\x,0) -- (\x,5.2);
  \node[below] at (\x,-0.28) {\lat};
}
% Cellule de Hadley (0-3)
\draw[-{Stealth[length=3mm]}, very thick, poporange] (0.35,0.3) -- (0.35,4.4);
\draw[-{Stealth[length=3mm]}, very thick, poporange] (0.35,4.6) -- (2.65,4.6);
\draw[-{Stealth[length=3mm]}, very thick, poporange] (2.65,4.4) -- (2.65,0.3);
\draw[-{Stealth[length=3mm]}, very thick, poporange] (2.65,0.18) -- (0.55,0.18);
\node[poporange, font=\bfseries] at (1.5,2.6) {Hadley};
% Cellule de Ferrel (3-6) - Indirecte
\draw[-{Stealth[length=3mm]}, very thick, popblue] (3.35,0.18) -- (5.45,0.18);
\draw[-{Stealth[length=3mm]}, very thick, popblue] (5.65,0.3) -- (5.65,4.4);
\draw[-{Stealth[length=3mm]}, very thick, popblue] (5.65,4.6) -- (3.35,4.6);
\draw[-{Stealth[length=3mm]}, very thick, popblue] (3.35,4.4) -- (3.35,0.3);
\node[popblue, font=\bfseries] at (4.5,2.6) {Ferrel};
% Cellule Polaire (6-9) - Directe
\draw[-{Stealth[length=3mm]}, very thick, poppurple] (8.8,4.4) -- (8.8,0.3);
\draw[-{Stealth[length=3mm]}, very thick, poppurple] (8.8,0.18) -- (6.2,0.18);
\draw[-{Stealth[length=3mm]}, very thick, poppurple] (6.2,0.3) -- (6.2,4.4);
\draw[-{Stealth[length=3mm]}, very thick, poppurple] (6.2,4.6) -- (8.8,4.6);
\node[poppurple, font=\bfseries] at (7.5,2.6) {Polaire};
% Zones de pression
\node[fill=popgreenL, rounded corners, inner sep=2pt] at (0,5.55) {\textbf{B} (dépression équatoriale)};
\node[fill=poppinkL, rounded corners, inner sep=2pt] at (3,5.55) {\textbf{H} (anticyclone subtropical)};
\node[fill=popgreenL, rounded corners, inner sep=2pt] at (6,5.55) {\textbf{B} (dépression subpolaire)};
\node[fill=poppinkL, rounded corners, inner sep=2pt] at (9,5.55) {\textbf{H} (anticyclone polaire)};
% Légendes ascendance/subsidence
\node[poporange, rotate=90, font=\scriptsize] at (0.05,2.4) {ascendance};
\node[poporange, rotate=90, font=\scriptsize] at (2.95,2.4) {subsidence};
\node[popblue, rotate=90, font=\scriptsize] at (5.95,2.4) {ascendance};
\node[poppurple, rotate=90, font=\scriptsize] at (8.95,2.4) {subsidence};
\node[poppurple, rotate=90, font=\scriptsize] at (6.2,2.4) {ascendance};
\end{tikzpicture}
\end{center}
\legende{Coupe méridienne de l'hémisphère Nord montrant les trois cellules de convection (Hadley, Ferrel, polaire). B : basse pression (air ascendant) ; H : haute pression (air descendant). L'hémisphère Sud présente une organisation symétrique.}
\end{popfigure}

\courssub{Les zones de pression planétaires}

Des trois cellules découle une alternance régulière de ceintures de pression tout autour du globe :

\begin{itemize}
\item \textbf{Dépression équatoriale (vers 0°)} : l'air chaud s'élève en permanence ; c'est une zone de calmes, d'airs convergents et de fortes pluies (ascendance = refroidissement de l'air en altitude = condensation de la vapeur d'eau = nuages et précipitations).
\item \textbf{Anticyclones subtropicaux (vers 30°)} : l'air descend, se réchauffe et se dessèche ; ce sont des zones de temps sec et stable. C'est sous ces latitudes que se trouvent les grands déserts du monde : le \textbf{Sahara}, le désert d'Arabie, le Kalahari, le désert australien.
\item \textbf{Dépressions subpolaires (vers 60°)} : l'air remonte, le temps y est perturbé et pluvieux.
\item \textbf{Anticyclones polaires (vers 90°)} : air froid, dense, descendant, temps sec et glacial.
\end{itemize}

\begin{aretenir}[La règle d'or]
\textbf{Ascendance} d'air $\Rightarrow$ basse pression $\Rightarrow$ nuages et pluies. \textbf{Subsidence} d'air $\Rightarrow$ haute pression $\Rightarrow$ temps sec et ciel dégagé. C'est pourquoi il pleut abondamment à l'équateur et presque jamais au cœur du Sahara.
\end{aretenir}

\courssub{La force de Coriolis et les vents planétaires}

\textbf{Observation.} Si les vents de surface allaient simplement des hautes vers les basses pressions, les alizés souffleraient plein nord ou plein sud. Or les marins savent depuis des siècles que les alizés soufflent du \emph{nord-est} dans l'hémisphère Nord et du \emph{sud-est} dans l'hémisphère Sud. Les vents sont déviés.

\begin{propriete}[Force de Coriolis (admise)]
La \cle{force de Coriolis}, due à la rotation de la Terre sur elle-même, dévie tout mouvement d'air (ou d'eau) \textbf{vers la droite} de sa trajectoire dans l'hémisphère Nord et \textbf{vers la gauche} dans l'hémisphère Sud. Elle est \textbf{nulle à l'équateur} et maximale aux pôles. (Justification qualitative : un point de l'équateur tourne plus vite qu'un point de haute latitude ; une masse d'air qui change de latitude conserve sa vitesse de rotation initiale, ce qui la décale par rapport au sol.)
\end{propriete}

En combinant les zones de pression et la déviation de Coriolis, on obtient les grands \cle{vents planétaires} :

\begin{itemize}
\item \textbf{Les alizés (0°--30°)} : ils soufflent des anticyclones subtropicaux vers la dépression équatoriale. Déviés par Coriolis, ils deviennent des vents de \textbf{nord-est} dans l'hémisphère Nord et de \textbf{sud-est} dans l'hémisphère Sud. Ce sont des vents remarquablement réguliers en direction et en force.
\item \textbf{Les vents d'ouest des moyennes latitudes (30°--60°)} : ils soufflent des anticyclones subtropicaux vers les dépressions subpolaires et sont déviés en vents dominants d'ouest (sud-ouest dans l'hémisphère Nord). Ils apportent les perturbations pluvieuses de l'Europe tempérée.
\item \textbf{Les vents d'est polaires (60°--90°)} : ils s'échappent des anticyclones polaires et sont déviés en vents d'est, froids et secs.
\end{itemize}

\begin{popfigure}
\begin{center}
\begin{tikzpicture}[scale=1.0, font=\small]
% Globe stylisé : rectangle équateur-centré
\draw[thick, popdark, fill=popblueL!40] (0,0) rectangle (12,6);
% Latitudes
\draw[thick, popdark] (0,3) -- (12,3);
\node[left, popdark] at (0,3) {Équateur 0°};
\foreach \y/\lat in {4.5/30°N, 5.75/60°N, 1.5/30°S, 0.25/60°S}{
  \draw[dashed, popmuted] (0,\y) -- (12,\y);
  \node[left, popmuted] at (0,\y) {\lat};
}
% ZCIT
\draw[very thick, poporange, decorate, decoration={snake, amplitude=0.8mm}] (0,3.15) -- (12,3.15);
\node[poporange, font=\bfseries] at (10.6,3.45) {ZCIT};
% Alizés NE (hémisphère nord) : de 30°N vers équateur, déviés vers l'ouest
\foreach \x in {1.5, 4.5, 7.5}{
  \draw[-{Stealth[length=3mm]}, very thick, popgreen] (\x+0.8,4.4) .. controls (\x+0.5,4.0) .. (\x,3.3);
}
\node[popgreen, font=\bfseries] at (2.2,4.9) {Alizés de NE};
% Alizés SE
\foreach \x in {1.5, 4.5, 7.5}{
  \draw[-{Stealth[length=3mm]}, very thick, popgreen] (\x+0.8,1.6) .. controls (\x+0.5,2.0) .. (\x,2.7);
}
\node[popgreen, font=\bfseries] at (2.2,1.1) {Alizés de SE};
% Vents d'ouest nord
\foreach \x in {2, 5.5, 9}{
  \draw[-{Stealth[length=3mm]}, very thick, poppurple] (\x,4.6) .. controls (\x+0.5,5.1) .. (\x+0.9,5.6);
}
\node[poppurple, font=\bfseries] at (6.5,5.85) {Vents d'ouest (hém. N)};
% Vents d'ouest sud
\foreach \x in {2, 5.5, 9}{
  \draw[-{Stealth[length=3mm]}, very thick, poppurple] (\x,1.4) .. controls (\x+0.5,0.9) .. (\x+0.9,0.4);
}
\node[poppurple, font=\bfseries] at (6.5,0.12) {Vents d'ouest (hém. S)};
\end{tikzpicture}
\end{center}
\legende{Carte simplifiée des vents planétaires. Les alizés (verts) convergent vers la zone de convergence intertropicale (ZCIT, en orange) ; les vents d'ouest (violets) dominent aux moyennes latitudes. La courbure des flèches traduit la déviation de Coriolis : vers la droite au nord, vers la gauche au sud.}
\end{popfigure}

\begin{savaistu}[D'où vient le mot « alizé » ?]
Le mot « alizé » viendrait du portugais ancien \emph{alísio}, « régulier, sans obstacle ». Ces vents d'une constance remarquable ont fait l'histoire : ce sont eux qui ont permis aux caravelles européennes de Christophe Colomb d'atteindre les Amériques en suivant la route des alizés de nord-est, puis de revenir plus au nord en profitant des vents d'ouest. Avant la vapeur, toute la navigation atlantique à voile était organisée autour de ces autoroutes du vent.
\end{savaistu}

\courssub{La zone de convergence intertropicale (ZCIT) et les saisons au Cameroun}

Là où les alizés de nord-est et de sud-est se rencontrent, l'air converge et s'élève en masse : c'est la \cle{zone de convergence intertropicale} (\cle{ZCIT}, ou ITCZ en anglais), une ceinture de nuages, d'averses violentes et de calmes que les marins redoutent (le « pot au noir »). La ZCIT n'est pas fixe : elle suit avec retard le déplacement apparent du Soleil et se déplace donc vers le nord en juillet, vers le sud en janvier.

Ce balancement saisonnier explique directement l'alternance des saisons en Afrique de l'Ouest et au Cameroun :

\begin{itemize}
\item \textbf{En juillet}, la ZCIT remonte vers le nord (parfois jusqu'à 15--20°N). L'alizé de sud-est de l'hémisphère Sud franchit l'équateur, est dévié vers la droite (Coriolis) et devient un vent de sud-ouest chargé d'humidité atlantique : c'est la \cle{mousson}. Elle apporte la saison des pluies au Cameroun et au Sahel.
\item \textbf{En janvier}, la ZCIT redescend vers le sud. Le Cameroun septentrional est alors balayé par l'alizé de nord-est continental, venant du Sahara : c'est le \cle{harmattan}, vent sec, chaud et chargé de poussière, qui assèche la végétation, irrite les voies respiratoires et voile le ciel de Garoua à Maroua.
\end{itemize}

\begin{popfigure}
\begin{center}
\begin{tikzpicture}[scale=0.92, font=\small]
% --- Janvier ---
\begin{scope}
\draw[thick, popdark, fill=popgoldL!50] (0,0) rectangle (5,3.4);
\draw[thick, popblue!60, fill=popblueL!60] (0,0) rectangle (5,1.1);
\node[popblue!70!black] at (2.5,0.55) {Océan Atlantique};
\node[popdark] at (2.5,3.1) {Sahara};
\draw[dashed, popmuted] (0,1.6) -- (5,1.6) node[right, popmuted]{Équateur};
\draw[very thick, poporange] (0,1.35) -- (5,1.35);
\node[poporange, font=\bfseries, anchor=west] at (0.1,1.6) {ZCIT (sud)};
\draw[-{Stealth[length=3mm]}, very thick, popgold!80!black] (3.8,3.0) -- (2.2,1.6);
\node[popgold!80!black, font=\bfseries, align=center] at (4.1,2.3) {Harmattan\\ (sec, NE)};
\node[font=\bfseries] at (2.5,-0.5) {JANVIER : saison sèche au Nord};
\end{scope}
% --- Juillet ---
\begin{scope}[xshift=7cm]
\draw[thick, popdark, fill=popgoldL!50] (0,0) rectangle (5,3.4);
\draw[thick, popblue!60, fill=popblueL!60] (0,0) rectangle (5,1.1);
\node[popblue!70!black] at (2.5,0.55) {Océan Atlantique};
\node[popdark] at (2.5,3.1) {Sahara};
\draw[dashed, popmuted] (0,1.6) -- (5,1.6) node[right, popmuted]{Équateur};
\draw[very thick, poporange] (0,2.6) -- (5,2.6);
\node[poporange, font=\bfseries, anchor=west] at (0.1,2.85) {ZCIT (nord)};
\draw[-{Stealth[length=3mm]}, very thick, popgreen!60!black] (1.2,0.35) .. controls (1.6,1.2) .. (2.6,2.3);
\node[popgreen!60!black, font=\bfseries, align=center] at (1.0,2.2) {Mousson\\ (humide, SO)};
\node[font=\bfseries] at (2.5,-0.5) {JUILLET : saison des pluies};
\end{scope}
\end{tikzpicture}
\end{center}
\legende{Position de la ZCIT au-dessus de l'Afrique de l'Ouest en janvier (à gauche) et en juillet (à droite). En janvier, le harmattan sec domine ; en juillet, la mousson humide venue de l'Atlantique apporte les pluies jusqu'au Sahel.}
\end{popfigure}

\begin{methode}[Décrire un mouvement atmosphérique sur un schéma ou une carte]
\begin{enumerate}
\item \textbf{Repérer les zones de pression} : localiser les hautes pressions (H, anticyclones) et les basses pressions (B, dépressions).
\item \textbf{Identifier le mouvement vertical} : ascendance (air chaud, basse pression, pluies probables) ou subsidence (air froid, haute pression, temps sec).
\item \textbf{Déterminer le sens horizontal} : le vent de surface va toujours des H vers les B.
\item \textbf{Appliquer la déviation de Coriolis} : courber la trajectoire vers la droite dans l'hémisphère Nord, vers la gauche dans l'hémisphère Sud.
\item \textbf{Conclure} en nommant le vent obtenu (alizé, vent d'ouest, mousson, harmattan...) et en précisant ses conséquences climatiques (sec ou humide).
\end{enumerate}
\end{methode}

\begin{exoresolu}[Décrire le vent qui balaie Garoua en janvier]
Énoncé : En janvier, la ville de Garoua (Nord-Cameroun, environ 9°N) est balayée par un vent sec et poussiéreux. À l'aide de tes connaissances sur la circulation générale de l'atmosphère, décris l'origine et le trajet de ce vent, puis explique pourquoi il est sec.
\tcblower
Solution détaillée :
\begin{enumerate}
\item \textbf{Zones de pression} : en janvier, la ZCIT (dépression équatoriale) est descendue au sud de l'équateur, tandis que l'anticyclone subtropical de l'hémisphère Nord domine le Sahara vers 20--30°N.
\item \textbf{Mouvement vertical} : sous l'anticyclone saharien, l'air descend (subsidence), se réchauffe et se dessèche.
\item \textbf{Mouvement horizontal} : l'air de surface s'écoule de l'anticyclone saharien (H) vers la ZCIT (B), donc globalement du nord vers le sud.
\item \textbf{Déviation de Coriolis} : dans l'hémisphère Nord, ce flux est dévié vers la droite et devient un vent de \textbf{nord-est} : c'est l'alizé continental de NE, appelé \textbf{harmattan}.
\item \textbf{Conclusion} : ayant traversé des milliers de kilomètres de désert sans rencontrer d'océan, cet air n'a pas pu se recharger en vapeur d'eau ; il arrive à Garoua chaud, très sec et chargé de poussière saharienne, d'où la sécheresse, le ciel voilé et les gênes respiratoires observées.
\end{enumerate}
\end{exoresolu}

\begin{exercice}[À toi de jouer]
En juillet, Douala reçoit des pluies abondantes apportées par un vent de sud-ouest. En appliquant la méthode en cinq étapes, décris l'origine de ce vent, explique pourquoi il est humide, et précise le rôle de la force de Coriolis dans sa direction. Pourquoi ce vent provoque-t-il des pluies en atteignant les reliefs comme le Mont Cameroun ?
\end{exercice}

\begin{corrige}[Exercice : la mousson de juillet à Douala]
1) \textbf{Pressions} : en juillet, la ZCIT est remontée au nord du Cameroun ; l'anticyclone subtropical de l'hémisphère Sud (vers 30°S, sur l'Atlantique) alimente l'alizé de sud-est. 2) \textbf{Mouvement vertical} : l'air s'élève le long de la ZCIT, d'où nuages et pluies. 3) \textbf{Mouvement horizontal} : l'alizé de SE souffle de l'anticyclone sud-atlantique vers la ZCIT, en traversant l'océan Atlantique équatorial, ce qui le charge fortement en vapeur d'eau. 4) \textbf{Coriolis} : en franchissant l'équateur, ce vent passe dans l'hémisphère Nord où la déviation s'exerce vers la droite : il vire et devient un vent de \textbf{sud-ouest}, la mousson. 5) \textbf{Conclusion} : arrivant sur la côte, cet air humide est forcé de s'élever le long des flancs du Mont Cameroun ; en montant, il se refroidit, sa vapeur d'eau se condense : il en résulte des pluies orographiques abondantes (Debundscha, au pied du Mont Cameroun, est l'un des lieux les plus arrosés d'Afrique, avec environ \SI{10000}{mm} de pluie par an).
\end{corrige}

\begin{aretenir}[L'essentiel de la section]
\begin{itemize}
\item Le moteur de la circulation atmosphérique est le \textbf{surplus d'énergie solaire à l'équateur} et le \textbf{déficit aux pôles}.
\item Air chaud = moins dense = \textbf{ascendance} = basse pression = pluies ; air froid = plus dense = \textbf{subsidence} = haute pression = temps sec.
\item Chaque hémisphère compte \textbf{trois cellules} : Hadley (0--30°), Ferrel (30--60°), polaire (60--90°).
\item La \textbf{force de Coriolis} dévie les vents à droite au Nord, à gauche au Sud ; elle est nulle à l'équateur.
\item Il en résulte les \textbf{alizés}, les \textbf{vents d'ouest} et les \textbf{vents polaires d'est}, ainsi que la \textbf{ZCIT}, dont le balancement saisonnier commande la mousson humide et le harmattan sec au Cameroun.
\end{itemize}
\end{aretenir}

\coursec{La circulation océanique : les courants marins}

Dans la section précédente, nous avons vu que l'atmosphère est parcourue par de grands mouvements organisés : les cellules de convection et les vents dominants (alizés, vents d'ouest) redistribuent la chaleur entre l'équateur et les pôles. Mais l'atmosphère n'est pas seule à transporter l'énergie thermique sur la planète. Sous ses vents, l'océan — qui recouvre environ 71 \% de la surface du globe — se déplace lui aussi, lentement mais puissamment, formant de véritables « fleuves » au milieu des mers : ce sont les \cle{courants marins}. Comprendre ces mouvements est indispensable pour expliquer le climat du littoral camerounais, la richesse de certaines pêcheries, ou encore la formation des cyclones que nous étudierons plus loin.

\courssub{Qu'est-ce qu'un courant marin ?}

\textbf{Une observation pour commencer.} En 1992, un conteneur transportant des milliers de canards en plastique s'est déversé dans le Pacifique Nord lors d'une tempête. Pendant plus de quinze ans, ces jouets se sont échoués sur les côtes de l'Alaska, d'Hawaï, du Japon, puis de l'Europe occidentale, suivant des trajets précis et prévisibles. Comment des objets flottants peuvent-ils parcourir des milliers de kilomètres en suivant toujours les mêmes routes ? Parce qu'ils sont transportés par des déplacements d'ensemble des masses d'eau : les courants marins.

\begin{definition}[Courant marin]
Un \cle{courant marin} est un \textbf{déplacement permanent ou saisonnier d'une masse d'eau océanique}, en surface ou en profondeur. Il est caractérisé par :
\begin{itemize}
\item sa \textbf{direction} et son \textbf{sens} de déplacement ;
\item sa \textbf{vitesse} (de quelques centaines de mètres par jour à plus de \SI{2}{m/s} pour les plus rapides) ;
\item sa \textbf{température} (courant chaud ou courant froid) ;
\item son \textbf{origine géographique} (latitude de formation).
\end{itemize}
\end{definition}

\begin{attention}[Ce qu'un courant n'est pas]
Un courant marin n'est ni une vague ni la marée. Les vagues sont des mouvements de surface provoqués localement par le vent, sans transport net d'eau sur de longues distances ; les marées sont des oscillations du niveau de la mer dues à l'attraction de la Lune et du Soleil. Un courant, lui, \emph{transporte} réellement l'eau sur des milliers de kilomètres, de façon durable.
\end{attention}

\courssub{Les courants de surface : l'océan entraîné par les vents}

\textbf{Intuition.} Souffle sur la surface d'une tasse de thé : le liquide se met en mouvement dans le sens de ton souffle. À l'échelle de l'océan, les vents dominants jouent exactement ce rôle.

Les \cle{courants de surface} (jusqu'à environ 400 m de profondeur) sont provoqués par le \textbf{frottement des vents dominants} sur la surface de l'eau :
\begin{itemize}
\item les \textbf{alizés}, qui soufflent des tropiques vers l'équateur, poussent les eaux de surface d'est en ouest dans la zone intertropicale : ce sont les courants équatoriaux ;
\item les \textbf{vents d'ouest} des latitudes moyennes poussent les eaux d'ouest en est.
\end{itemize}

Ces courants, déviés par les continents et par la rotation de la Terre, dessinent des circuits fermés. On distingue alors deux grandes familles de courants selon leur origine :

\begin{definition}[Courants chauds et courants froids]
Un \cle{courant chaud} provient des \textbf{basses latitudes} (régions tropicales et équatoriales) et transporte de l'eau chaude vers les latitudes plus élevées. Exemple : le \textbf{Gulf Stream} dans l'Atlantique Nord.
Un \cle{courant froid} provient des \textbf{hautes latitudes} (régions polaires ou subpolaires) ou des profondeurs, et transporte de l'eau froide vers l'équateur. Exemple : le \textbf{courant de Benguela}, qui longe la côte sud-ouest de l'Afrique.
\end{definition}

\begin{exemplebox}[Le Gulf Stream, un fleuve dans l'océan]
Le Gulf Stream naît dans le golfe du Mexique, longe la côte est des États-Unis puis traverse l'Atlantique Nord vers l'Europe occidentale. Il transporte entre \textbf{30 et 150 millions de mètres cubes d'eau par seconde} (soit 30 à 150 Sverdrups), soit plus de \textbf{100 fois le débit de tous les fleuves du monde réunis}. Grâce à lui, le climat de l'Europe de l'Ouest est beaucoup plus doux que celui des régions situées à la même latitude au Canada : le port de Bergen (Norvège, 60° N) reste libre de glace en hiver, alors que des ports canadiens à la même latitude sont pris dans les glaces.
\end{exemplebox}

\courssub{Les courants de profondeur : la circulation thermohaline}

\textbf{Problème.} Les vents n'agissent que sur les couches superficielles. Pourtant, des mesures de température montrent que l'eau des grands fonds océaniques, même sous les tropiques, est très froide (environ \SI{2}{\celsius} à \SI{4}{\celsius}) : elle vient donc des régions polaires. Quel moteur peut faire circuler l'eau en profondeur, là où le vent ne souffle pas ?

\begin{experience}[Comprendre les courants de densité]
\textbf{Matériel :} un aquarium transparent, de l'eau froide colorée en bleu, de l'eau chaude colorée en rouge, une cloison amovible.
\textbf{Protocole :} remplir l'aquarium en séparant les deux eaux par la cloison (eau froide d'un côté, eau chaude de l'autre), puis retirer délicatement la cloison.
\textbf{Observation :} l'eau froide bleue, plus dense, s'écoule vers le fond et glisse sous l'eau chaude ; l'eau chaude rouge, moins dense, s'étale en surface. Deux courants de sens opposés s'établissent.
\textbf{Interprétation :} les différences de \textbf{densité} suffisent à mettre l'eau en mouvement, sans aucun vent. La densité de l'eau de mer dépend de sa \textbf{température} (l'eau froide est plus dense) et de sa \textbf{salinité} (l'eau plus salée est plus dense).
\end{experience}

\begin{definition}[Circulation thermohaline]
La \cle{circulation thermohaline} est l'ensemble des courants de profondeur provoqués par les \textbf{différences de densité} des masses d'eau, elles-mêmes dues aux différences de \textbf{température} (thermo-) et de \textbf{salinité} (-haline).
\end{definition}

\textbf{Mécanisme.} Aux hautes latitudes (Atlantique Nord, Antarctique), l'eau de surface se refroidit fortement ; la formation de la glace de mer, qui rejette du sel, augmente encore la salinité de l'eau restante. Cette eau, devenue très dense, \textbf{plonge vers les grands fonds} (formation des eaux profondes : NADW en Atlantique Nord, AABW en Antarctique) puis s'écoule lentement vers l'équateur. Ailleurs, de l'eau profonde remonte vers la surface pour compenser (remontée diffuse ou upwelling côtier). L'ensemble forme un vaste circuit planétaire, parfois appelé « tapis roulant » océanique, dont le cycle complet dure environ un millier d'années.

\courssub{Les grands gyres océaniques}

\textbf{Observation.} Sur une carte mondiale des courants de surface, on remarque que dans chaque grand bassin océanique (Atlantique Nord, Atlantique Sud, Pacifique Nord, etc.), les courants dessinent d'immenses boucles fermées appelées \cle{gyres}.

\begin{propriete}[Sens de rotation des gyres]
Sous l'effet combiné des vents dominants, de la position des continents et de la \textbf{force de Coriolis} (déviation due à la rotation de la Terre), les gyres tournent :
\begin{itemize}
\item dans le \textbf{sens horaire} dans l'\textbf{hémisphère Nord} ;
\item dans le \textbf{sens antihoraire} dans l'\textbf{hémisphère Sud}.
\end{itemize}
\end{propriete}

\begin{popfigure}
\begin{center}
\begin{tikzpicture}[scale=1,>=stealth]
% Océan
\fill[popblue!15] (0,0) rectangle (10,5);
\draw[popline] (0,0) rectangle (10,5);
% Continents stylisés
\fill[popgreen!35] (0,0.6) -- (0,4.6) -- (1.4,4.2) -- (1.0,3.0) -- (1.6,2.0) -- (1.1,0.8) -- cycle;
\fill[popgreen!35] (10,0.7) -- (10,4.5) -- (8.7,4.3) -- (9.1,3.1) -- (8.5,1.9) -- (9.0,0.9) -- cycle;
\node[popdark,font=\small] at (0.8,2.6) {Amérique};
\node[popdark,font=\small] at (9.3,2.6) {Europe--Afrique};
% Équateur
\draw[dashed,popmuted] (0,1.7) -- (10,1.7);
\node[popmuted,font=\scriptsize,anchor=west] at (0.1,1.85) {Équateur};
% Gyre horaire (hémisphère Nord)
\draw[->,very thick,poporange] (2.2,3.9) .. controls (4.5,4.5) and (6.5,4.5) .. (8.3,3.8);
\draw[->,very thick,poporange] (8.3,3.8) .. controls (8.6,3.2) and (8.5,2.9) .. (8.2,2.6);
\draw[->,very thick,popblue] (8.2,2.4) .. controls (6.5,1.95) and (4.5,1.95) .. (2.4,2.4);
\draw[->,very thick,popblue] (2.4,2.4) .. controls (2.0,2.9) and (2.0,3.4) .. (2.2,3.9);
\node[poporange,font=\small\bfseries] at (5.3,4.75) {Gulf Stream (chaud)};
\node[popblue,font=\small\bfseries] at (5.3,2.05) {Courant des Canaries (froid)};
% Sens horaire indiqué
\draw[->,thick,popdark] (5.3,3.3) arc (90:-180:0.55);
\node[popdark,font=\scriptsize] at (5.3,3.9) {sens horaire};
% Courant de Guinée
\draw[->,very thick,poporange] (8.3,1.4) .. controls (7.0,1.0) and (6.0,1.0) .. (5.0,1.3);
\node[poporange,font=\scriptsize] at (6.6,0.75) {Courant de Guinée (chaud)};
\end{tikzpicture}
\end{center}
\legende{Le gyre de l'Atlantique Nord : la rotation horaire associe un courant chaud sur le bord ouest (Gulf Stream, en orange) et un courant froid sur le bord est (courant des Canaries, en bleu). Au sud, le courant de Guinée, chaud, longe le golfe de Guinée.}
\end{popfigure}

\textbf{Conséquence remarquable :} dans un gyre, le courant situé du côté \emph{ouest} de l'océan remonte des tropiques vers les pôles (courant \textbf{chaud} : Gulf Stream, Kuro Shio), tandis que le courant du côté \emph{est} redescend des hautes latitudes vers l'équateur (courant \textbf{froid} : courant des Canaries, courant de Benguela, courant de Californie, courant de Humboldt).

\begin{popfigure}
\begin{center}
\begin{tikzpicture}[scale=0.95,>=stealth]
% Fond océan simplifié
\fill[popblue!12] (0,0) rectangle (12,6);
% Amériques
\fill[popgreen!30] (1.2,5.4) -- (2.6,5.6) -- (3.0,4.4) -- (2.4,3.4) -- (2.8,2.6) -- (2.2,1.2) -- (1.6,0.6) -- (1.2,2.0) -- (0.8,3.6) -- cycle;
% Afrique-Europe
\fill[popgreen!30] (5.6,5.6) -- (7.2,5.6) -- (7.4,4.6) -- (6.8,3.8) -- (7.2,2.8) -- (6.6,1.4) -- (6.0,0.8) -- (5.6,2.2) -- (5.2,3.6) -- (5.4,4.6) -- cycle;
% Asie-Océanie
\fill[popgreen!30] (7.8,5.4) -- (10.6,5.6) -- (11.2,4.4) -- (10.4,3.6) -- (10.8,2.8) -- (10.0,2.4) -- (9.4,3.2) -- (8.6,3.8) -- (7.8,4.4) -- cycle;
% Équateur
\draw[dashed,popmuted] (0,3.0) -- (12,3.0);
\node[popmuted,font=\scriptsize,anchor=west] at (0.1,3.12) {Équateur};
% Courants Atlantique Nord
\draw[->,very thick,poporange] (2.5,3.6) .. controls (3.2,4.4) and (4.4,5.0) .. (5.6,5.1);
\node[poporange,font=\scriptsize\bfseries] at (4.1,5.35) {Gulf Stream};
\draw[->,very thick,popblue] (5.5,4.6) .. controls (5.0,4.0) and (4.9,3.5) .. (4.6,3.1);
\node[popblue,font=\scriptsize\bfseries] at (5.0,4.0) {Canaries};
% Courants Atlantique Sud / Golfe de Guinée
\draw[->,very thick,poporange] (4.4,2.75) .. controls (5.2,2.55) and (5.9,2.6) .. (6.4,2.75);
\node[poporange,font=\scriptsize\bfseries] at (5.3,2.3) {Guinée};
\draw[->,very thick,popblue] (6.3,1.2) .. controls (5.9,1.7) and (5.7,2.1) .. (5.6,2.5);
\node[popblue,font=\scriptsize\bfseries] at (6.9,1.5) {Benguela};
% Courants Pacifique Nord
\draw[->,very thick,poporange] (9.0,2.9) .. controls (9.4,3.7) and (9.9,4.3) .. (10.4,4.6);
\node[poporange,font=\scriptsize\bfseries] at (10.9,3.9) {Kuro Shio};
\draw[->,very thick,popblue] (1.6,4.9) .. controls (1.4,4.2) and (1.5,3.6) .. (1.8,3.1);
\node[popblue,font=\scriptsize\bfseries] at (0.9,4.3) {Californie};
% Courants Pacifique Sud
\draw[->,very thick,popblue] (2.3,1.0) .. controls (2.0,1.6) and (1.95,2.2) .. (2.0,2.7);
\node[popblue,font=\scriptsize\bfseries] at (3.2,1.3) {Humboldt};
% Légende couleurs
\draw[->,very thick,poporange] (8.0,0.6) -- (8.9,0.6);
\node[anchor=west,font=\scriptsize,popdark] at (9.0,0.6) {courant chaud};
\draw[->,very thick,popblue] (10.4,0.6) -- (11.3,0.6);
\node[anchor=west,font=\scriptsize,popdark] at (11.4,0.6) {courant froid};
\end{tikzpicture}
\end{center}
\legende{Carte simplifiée des principaux courants marins de surface. Les courants chauds (orange) partent des basses latitudes ; les courants froids (bleu) viennent des hautes latitudes ou des remontées d'eau profonde.}
\end{popfigure}

\begin{methode}[Décrire un courant marin sur une carte]
Pour décrire correctement un courant à partir d'un schéma ou d'une carte (compétence exigible au Probatoire), indique dans l'ordre :
\begin{enumerate}
\item son \textbf{nom} (ex. : courant de Benguela) ;
\item son \textbf{origine} : latitude de formation (ex. : hautes latitudes de l'Atlantique Sud) ;
\item son \textbf{sens de déplacement} (ex. : du sud vers le nord, le long de la côte sud-ouest africaine) ;
\item sa \textbf{nature thermique} : chaud ou froid, en comparant sa température à celle des eaux environnantes ;
\item son \textbf{influence sur les côtes bordières} : climat, pêche, navigation.
\end{enumerate}
\end{methode}

\courssub{Le rôle des courants marins}

\textbf{1. Redistribution de la chaleur et influence sur les climats côtiers.} Les courants marins transportent d'énormes quantités de chaleur : avec l'atmosphère, ils assurent le transfert d'énergie de l'équateur vers les pôles. Leur influence sur le climat des côtes qu'ils longent est directe.

\begin{propriete}[Influence climatique des courants]
Les \textbf{courants chauds} réchauffent et humidifient l'air des côtes qu'ils longent : le climat y est plus \textbf{doux et pluvieux}. Les \textbf{courants froids} refroidissent la couche d'air au contact de l'eau, ce qui la \textbf{stabilise} (l'air froid, plus dense, ne s'élève pas ; peu de convection, donc peu de nuages de pluie) : le climat côtier devient \textbf{sec}, avec des \textbf{brouillards} fréquents. C'est pourquoi de grands \textbf{déserts côtiers} (Namib face au Benguela, Atacama face au courant de Humboldt) se trouvent face à des courants froids.
\end{propriete}

\begin{savaistu}[Le courant de Guinée et la pluie de Debundscha]
Le \cle{courant de Guinée}, courant \textbf{chaud}, longe le littoral du golfe de Guinée, donc celui du Cameroun. En réchauffant et en humidifiant l'air marin, il contribue à l'humidité exceptionnelle de notre côte. Lorsque cet air très humide rencontre les pentes du mont Cameroun (4 095 m), il est forcé de s'élever (ascendance orographique), se refroidit adiabatiquement et déverse des pluies abondantes : \textbf{Debundscha}, au pied du mont Cameroun, reçoit \textbf{plus de 10 mètres de pluie par an}, ce qui en fait l'un des endroits les plus pluvieux du monde. Inversement, à la même latitude que le Sahel, des côtes baignées par des courants froids sont presque désertiques : la nature du courant fait toute la différence.
\end{savaistu}

\textbf{2. Les remontées d'eau froide (upwelling) et les pêcheries.} Le long de certaines côtes (bord est des océans), les vents dominants (alizés) poussent les eaux de surface vers le large (transport d'Ekman). Pour les remplacer, de l'eau \textbf{profonde, froide et riche en sels nutritifs} (nitrates, phosphates, silicates) remonte vers la surface : c'est l'\cle{upwelling}. Ces nutriments nourrissent le phytoplancton, base de la chaîne alimentaire marine : les zones d'upwelling comptent parmi les \textbf{pêcheries les plus riches du monde} (côtes du Pérou, de Namibie, du Sénégal, de Mauritanie).

\begin{popfigure}
\begin{center}
\begin{tikzpicture}[scale=1,>=stealth]
% Mer
\fill[popblue!20] (0,1.5) rectangle (10,4.5);
% Fond en pente (talus continental)
\fill[popgold!40] (0,0) -- (10,0) -- (10,1.5) -- (3.5,1.5) -- (0,1.5) -- cycle;
\draw[thick,popgold!60] (3.5,1.5) -- (10,1.5); % trait de côte/sous-marin
% Côte / continent
\fill[popgreen!35] (0,1.5) -- (3.5,1.5) -- (3.5,4.5) -- (0,4.5) -- cycle;
\node[popdark,font=\small] at (1.6,3.8) {Continent};
% Surface mer
\draw[thick,popblue] (3.5,4.5) -- (10,4.5);
% Vent vers le large
\draw[->,very thick,popmuted] (4.5,5.0) -- (7.0,5.0);
\draw[->,very thick,popmuted] (4.5,5.4) -- (7.0,5.4);
\node[popmuted,font=\scriptsize,anchor=west] at (7.1,5.2) {Vent côtier (alizé)};
% Transport d'Ekman (eau de surface chassée à 90° à gauche dans hémisphère Sud, ici schématisé vers le large)
\draw[->,thick,popblue!70] (4.0,4.2) .. controls (6.0,4.1) and (8.0,4.1) .. (9.5,4.2);
\node[popblue!70,font=\scriptsize] at (6.8,3.8) {Eau de surface chassée au large};
% Remontée d'eau froide (upwelling)
\draw[->,very thick,popblue] (4.5,0.8) .. controls (4.3,1.8) and (4.1,2.8) .. (4.0,3.8);
\draw[->,very thick,popblue] (5.5,0.8) .. controls (5.1,1.9) and (4.7,2.9) .. (4.5,3.7);
\node[popblue,font=\scriptsize\bfseries,align=center] at (7.0,1.2) {Eau profonde froide\\(riche en nutriments)};
% Phytoplancton / poissons
\node[popgreen!70!black,font=\scriptsize] at (4.5,4.8) {Bloom phytoplanctonique};
\foreach \x in {5.5,6.2,6.9} {\node[popdark,font=\scriptsize] at (\x,3.5) {\textbullet};}
\node[popdark,font=\scriptsize] at (8.0,3.5) {Poissons (pêcheries)};
\end{tikzpicture}
\end{center}
\legende{Schéma en coupe d'un upwelling côtier (hémisphère Sud) : le vent (alizé) chasse l'eau de surface vers le large (transport d'Ekman) ; une eau profonde, froide et chargée en nutriments, remonte le long du talus continental et enrichit la zone euphotique, favorisant le phytoplancton puis les poissons.}
\end{popfigure}

\begin{exoresolu}[Décrire un courant et relier un phénomène à son origine]
\textbf{Énoncé.} Sur une carte de l'Atlantique Sud, on repère le long de la côte namibienne un courant remontant du sud vers le nord. La côte bordière est occupée par le désert du Namib, alors que l'intérieur des terres, à la même latitude, reçoit des pluies estivales. 1) Décris ce courant. 2) Explique l'aridité de la côte.
\tcblower
\textbf{Solution.} 1) Il s'agit du \textbf{courant de Benguela}. Il prend naissance aux \textbf{hautes latitudes} de l'Atlantique Sud (région subantarctique), se déplace \textbf{du sud vers le nord} le long de la côte sud-ouest africaine, et constitue un courant \textbf{froid}. Il appartient au gyre de l'Atlantique Sud, qui tourne dans le sens antihoraire (hémisphère Sud).
2) L'air au contact de ce courant froid est \textbf{refroidi et stabilisé} : ne pouvant s'élever par convection, il ne forme pas de nuages de pluie (seulement des brouillards d'advection). La côte namibienne reçoit donc très peu de précipitations, d'où le \textbf{désert côtier du Namib}. L'aridité côtière est ainsi directement reliée à la présence du courant froid (et de l'upwelling associé qui renforce le refroidissement).
\end{exoresolu}

\begin{attention}[Erreurs fréquentes]
\begin{itemize}
\item Ne dis pas qu'un courant froid « apporte le froid » comme un vent glacial : il \textbf{refroidit l'air} au contact de l'eau, ce qui \emph{empêche les ascendances convectives} et donc les pluies — c'est ce mécanisme qui explique les déserts côtiers.
\item Ne confonds pas le sens de rotation des gyres : \textbf{horaire au Nord, antihoraire au Sud}. Une inversion est une erreur classique au Probatoire.
\item Un courant chaud ne vient pas forcément d'une mer « chaude » : c'est son \textbf{origine en basse latitude} qui le définit (il transporte de la chaleur relative vers les pôles).
\item Ne confonds pas \emph{upwelling} (remontée d'eau profonde) et \emph{formation d'eau profonde} (enfoncement en haute latitude).
\end{itemize}
\end{attention}

\begin{aretenir}[L'essentiel de la section]
\begin{itemize}
\item Un courant marin est un déplacement durable d'une masse d'eau, défini par sa direction, sa vitesse et sa température.
\item Les courants de \textbf{surface} sont entraînés par les \textbf{vents dominants} ; les courants de \textbf{profondeur} (circulation thermohaline) résultent des différences de \textbf{densité} (température, salinité).
\item Les gyres tournent dans le sens \textbf{horaire au Nord} et \textbf{antihoraire au Sud} (force de Coriolis).
\item Courants \textbf{chauds} (Gulf Stream, courant de Guinée, Kuro Shio) : côtes douces et humides. Courants \textbf{froids} (Benguela, Canaries, Humboldt, Californie) : côtes sèches, brouillards, mais riches pêcheries grâce à l'\textbf{upwelling}.
\end{itemize}
\end{aretenir}

\begin{exercice}[Pour t'entraîner]
1) Explique, en t'appuyant sur la circulation thermohaline, pourquoi l'eau des grands fonds océaniques est froide même sous l'équateur.
2) Le climat de Douala est chaud et très humide, celui de la côte namibienne (Walvis Bay) est sec et frais. Relie chacun de ces climats au courant qui borde la côte correspondante et précise le mécanisme physique en jeu.
3) Reproduis le schéma du gyre de l'Atlantique Nord en plaçant correctement le Gulf Stream et le courant des Canaries, avec le sens de rotation. Indique la nature thermique de chaque courant.
\end{exercice}

\begin{corrige}[Exercices d'entraînement]
\textbf{1)} La circulation thermohaline est un circuit planétaire. Aux pôles (Atlantique Nord, Antarctique), l'eau de surface se refroidit et se concentre en sel (rejet saumure lors de la formation de la banquise). Elle devient très dense et \textbf{plonge vers le fond} (formation des eaux profondes : NADW, AABW). Ces masses d'eau froides et denses s'écoulent ensuite lentement dans les grands fonds de tous les océans, y compris sous l'équateur. C'est donc une eau d'origine polaire qui remplit les abysses.
\textbf{2)} 
\begin{itemize}
\item \textbf{Douala (Cameroun) :} Bordée par le \textbf{courant de Guinée (chaud)}. Mécanisme : le courant réchauffe la couche d'air au contact $\rightarrow$ l'air se charge en humidité (évaporation accrue) et devient instable $\rightarrow$ ascendances convectives et orographiques (mont Cameroun) $\rightarrow$ pluies abondantes, climat chaud et humide.
\item \textbf{Walvis Bay (Namibie) :} Bordée par le \textbf{courant de Benguela (froid)}. Mécanisme : le courant refroidit la couche d'air au contact $\rightarrow$ l'air se stabilise (inversion de température, air froid sous air chaud) $\rightarrow$ inhibition de la convection $\rightarrow$ absence de nuages de pluie, seulement brouillards $\rightarrow$ climat sec, désert côtier (Namib).
\end{itemize}
\textbf{3)} Schéma attendu : Bassin atlantique Nord. Flèches formant une boucle horaire. Bord ouest (côté Amérique) : flèche vers le Nord, étiquetée \textbf{Gulf Stream (chaud)}. Bord est (côté Europe/Afrique) : flèche vers le Sud, étiquetée \textbf{Courant des Canaries (froid)}. Bord sud (équateur) : flèche vers l'Ouest (Courant équatorial Nord, chaud). Bord nord : dérive nord-atlantique vers l'Est. Indication explicite : \textbf{Rotation horaire}.
\end{corrige}

\coursec{Interactions océan-atmosphère}

Dans la section précédente, nous avons décrit les \cle{courants marins} et vu que l'océan n'est pas une masse d'eau immobile : il circule, entraîné par les vents et la rotation terrestre. Mais la relation entre l'océan et l'atmosphère n'est pas à sens unique. L'atmosphère entraîne l'océan, certes, mais l'océan, en retour, chauffe et humidifie l'atmosphère. C'est ce dialogue permanent, ces \cle{échanges continus} d'énergie et de matière entre les deux milieux, que nous allons maintenant étudier. Ils expliquent certains des phénomènes climatiques les plus importants de la planète : la \cle{mousson} qui rythme les saisons en Afrique de l'Ouest, et le phénomène \cle{El Niño} qui bouleverse le climat du monde entier.

\courssub{L'océan, réservoir de chaleur et d'humidité}

\textbf{Une observation de départ.}\par
En août, à Douala, l'air est lourd, saturé d'humidité, et les pluies sont abondantes. D'où vient toute cette eau ? Elle ne tombe pas du vide : elle provient de l'\cle{évaporation} de l'océan Atlantique, tout proche. Sous l'effet du rayonnement solaire, l'eau de surface de l'océan s'évapore en permanence ; cette vapeur d'eau s'incorpore à l'atmosphère, est transportée par les vents, puis se condense en nuages et retombe en \cle{précipitations}. L'océan est ainsi la source quasi unique de l'eau des pluies : environ 86 \% de la vapeur d'eau atmosphérique provient de l'évaporation océanique.

Mais l'océan ne fournit pas seulement de l'eau : il fournit aussi de la \cle{chaleur}. L'eau possède une capacité thermique considérable : elle se réchauffe lentement, mais elle stocke énormément d'énergie et la restitue lentement. C'est pourquoi, en bord de mer, les nuits restent douces même après une journée chaude : l'océan relâche vers l'atmosphère la chaleur accumulée le jour (on parle de \cle{flux de chaleur} océan-atmosphère).

\begin{propriete}[L'océan, régulateur thermique de la planète (admise)]
L'océan stocke environ \cle{1000 fois plus de chaleur} que l'atmosphère. Il absorbe la majeure partie de l'énergie solaire reçue par la Terre et la redistribue lentement, avec une \cle{inertie de plusieurs années}. L'océan régule donc le climat de la planète : sans lui, les écarts de température entre le jour et la nuit, entre les saisons et entre les latitudes seraient extrêmes et invivables.
\end{propriete}

\textbf{Les échanges sont continus et réciproques.}\par
On peut résumer le dialogue océan-atmosphère en deux flèches :

\begin{itemize}
\item \textbf{L'océan agit sur l'atmosphère} : il lui fournit de la chaleur (flux de chaleur sensible et latent lié à l'évaporation) et de la vapeur d'eau, qui alimentera nuages et précipitations.
\item \textbf{L'atmosphère agit sur l'océan} : les \cle{vents}, par friction sur la surface de l'eau, entraînent les \cle{courants marins} de surface (alizés $\rightarrow$ courants équatoriaux, comme nous l'avons vu dans la section précédente) et brassent les eaux superficielles.
\end{itemize}

\begin{exemplebox}[Pourquoi pleut-il davantage sur les côtes camerounaises qu'au Sahel ?]
L'air qui arrive sur le littoral camerounais vient de l'océan Atlantique : il est chargé de vapeur d'eau. En rencontrant les reliefs (Mont Cameroun, \SI{4095}{m}), il s'élève, se refroidit, et la vapeur se condense : il en résulte des pluies abondantes (Debundscha, au pied du Mont Cameroun, est l'une des localités les plus arrosées du monde, avec près de \SI{10000}{mm} de pluie par an). À l'inverse, l'air qui parvient à Maroua ou au Sahel a perdu son humidité en chemin ou provient du continent sec : les précipitations y sont faibles. L'océan est bien la \emph{source} de l'eau des pluies.
\end{exemplebox}

\courssub{La mousson : un basculement saisonnier des vents}

\textbf{Situation concrète.}\par
Tout habitant de l'Afrique de l'Ouest connaît le rythme des saisons : une saison des pluies marquée par des vents venus de l'océan, humides et frais, puis une saison sèche dominée par l'\cle{harmattan}, vent chaud, sec et poussiéreux venu du Sahara. Comment expliquer ce basculement régulier de la direction des vents ?

\begin{definition}[La mousson]
La \cle{mousson} est un \cle{vent saisonnier} qui \cle{inverse sa direction} entre l'été et l'hiver, sous l'effet du \cle{contraste de réchauffement} entre les continents et les océans. En été, elle souffle de l'océan vers le continent (mousson humide) ; en hiver, elle souffle du continent vers l'océan (mousson sèche).
\end{definition}

\textbf{Le mécanisme : le continent et l'océan ne se réchauffent pas à la même vitesse.}\par
La roche et le sol se réchauffent vite et se refroidissent vite ; l'eau, grâce à sa forte capacité thermique, se réchauffe lentement et se refroidit lentement. De là découlent deux situations opposées :

\begin{itemize}
\item \textbf{En été (mousson humide)} : le continent se réchauffe plus vite que l'océan. L'air chaud au-dessus de la terre se dilate et s'élève : il se forme une \cle{dépression} (zone de basse pression) sur le continent, tandis que l'océan plus frais reste surmonté d'un \cle{anticyclone} (zone de haute pression). L'air circule des hautes vers les basses pressions : un flux d'air \cle{humide} souffle de l'océan vers le continent, apportant nuages et \cle{précipitations}.
\item \textbf{En hiver (mousson sèche)} : le continent se refroidit plus vite que l'océan. La situation s'inverse : anticyclone sur le continent froid, pressions relativement basses sur l'océan plus doux. Le vent souffle du continent vers l'océan : c'est un air \cle{sec}, qui ne produit aucune pluie.
\end{itemize}

\begin{exemplebox}[Mousson et harmattan en Afrique de l'Ouest]
En été boréal (juin à septembre), le Sahara et le Sahel surchauffent : la dépression sahélienne aspire l'air humide de l'océan Atlantique Sud. C'est la \cle{mousson d'Afrique de l'Ouest} : des vents de sud-ouest chargés d'humidité remontent du golfe de Guinée vers l'intérieur des terres et déclenchent la saison des pluies, du littoral camerounais jusqu'au Sahel. En hiver boréal (décembre à février), le continent se refroidit, l'anticyclone s'installe sur le Sahara et le flux s'inverse : l'\cle{harmattan}, vent de nord-est sec et chargé de poussières, descend du désert vers le golfe de Guinée. C'est la saison sèche, avec son ciel voilé et son air desséchant.
\end{exemplebox}

\begin{popfigure}
\begin{tikzpicture}[scale=1.0, font=\small]
% Océan
\fill[popblue!35] (0,0) rectangle (5,1.2);
\node[popdark] at (2.5,0.55) {\textbf{Océan plus froid}};
% Continent
\fill[popgold!40] (5,0) rectangle (11,1.2);
\node[popdark] at (8,0.55) {\textbf{Continent chaud (été)}};
% Soleil
\draw[poporange, fill=popyellow, line width=1pt] (8,5.6) circle (0.45);
\foreach \a in {0,45,...,315} \draw[poporange, line width=1pt] (8,5.6) ++(\a:0.55) -- ++(\a:0.35);
\node[poporange] at (8,6.5) {Rayonnement solaire};
% Anticyclone sur océan
\node[popblue, font=\bfseries] at (2.5,4.6) {Anticyclone (hautes pressions)};
% Dépression sur continent
\node[poppink, font=\bfseries] at (8,4.6) {Dépression (basses pressions)};
% Flux d'air en surface : océan -> continent
\draw[->, popblue, line width=2pt] (1.2,1.7) -- (6.8,1.7) node[midway, below, popdark] {Flux d'air humide (mousson)};
% Air descendant sur océan
\draw[->, popblue, line width=1.6pt] (2.5,4.2) -- (2.5,2.0);
\node[popblue, right] at (2.6,3.1) {Air descendant};
% Air ascendant sur continent
\draw[->, poppink, line width=1.6pt] (8,1.8) -- (8,4.2);
\node[poppink, right] at (8.1,3.0) {Air chaud ascendant};
% Retour en altitude
\draw[->, popmuted, line width=1.4pt] (7.6,4.35) -- (2.9,4.35);
% Évaporation
\foreach \x in {1.0,2.0,3.0,4.0} \draw[->, popblue!70, line width=0.9pt] (\x,1.25) .. controls (\x+0.15,1.5) and (\x-0.15,1.7) .. (\x,1.95);
\node[popblue!80] at (1.0,2.35) {Évaporation};
% Nuages et pluie
\draw[popmuted, fill=popcream, line width=0.8pt] (9.2,3.9) ellipse (1.1 and 0.45);
\draw[popmuted, fill=popcream, line width=0.8pt] (8.5,4.15) ellipse (0.7 and 0.35);
\foreach \x in {8.6,9.0,9.4,9.8} \draw[popblue, line width=1pt] (\x,3.4) -- (\x-0.15,2.8);
\node[popblue] at (9.6,2.5) {Précipitations};
\end{tikzpicture}
\legende{Schéma en coupe de la mousson d'été. Le continent, réchauffé plus vite que l'océan, devient une zone de basse pression qui aspire l'air humide océanique ; cet air s'élève, se refroidit et provoque les précipitations de la saison des pluies.}
\end{popfigure}

\courssub{El Niño : quand le Pacifique se dérègle}

\textbf{Une observation troublante.}\par
Certaines années, les pêcheurs du Pérou constatent un désastre : les eaux côtières, normalement froides et poissonneuses, deviennent chaudes, les poissons disparaissent, et des pluies diluviennes s'abattent sur des côtes d'ordinaire désertiques. Au même moment, de l'autre côté du Pacifique, l'Indonésie et l'Australie subissent des sécheresses. Tous ces événements, séparés par des milliers de kilomètres, ont une cause commune : \cle{El Niño}.

\begin{methode}[Démarche d'investigation : comprendre d'abord la situation normale]
Pour saisir l'anomalie, il faut connaître le fonctionnement habituel du Pacifique équatorial :
\begin{enumerate}
\item \textbf{Observation (situation normale)} : les \cle{alizés} soufflent d'est en ouest le long de l'équateur. Ils poussent les eaux chaudes de surface vers l'ouest (vers l'Indonésie), où la mer est chaude (environ \SI{28}{\celsius}) et où l'air humide s'élève : pluies abondantes. À l'est (côtes du Pérou), les eaux de surface chassées sont remplacées par une remontée d'eau froide profonde (\cle{upwelling}) riche en nutriments : mer froide (environ \SI{20}{\celsius}), air stable, temps sec, mais pêche très riche.
\item \textbf{Observation (année El Niño)} : les alizés faiblissent ou s'inversent. Les eaux chaudes ne sont plus repoussées vers l'ouest : elles \emph{reviennent} vers l'est. La surface du Pacifique équatorial Est se réchauffe anormalement ; l'upwelling péruvien est bloqué.
\item \textbf{Interprétation} : la zone de pluies, qui suit toujours les eaux les plus chaudes, se déplace avec elles vers l'est. L'ouest (Indonésie, Australie) se dessèche ; l'est (Pérou, Équateur) est inondé. Sans upwelling, les nutriments disparaissent, les poissons aussi : la pêche péruvienne s'effondre.
\item \textbf{Conclusion} : un simple affaiblissement des alizés suffit à déplacer le « moteur thermique » du Pacifique et à bouleverser les climats à l'échelle de la planète entière.
\end{enumerate}
\end{methode}

\begin{definition}[El Niño]
\cle{El Niño} est un phénomène climatique \cle{périodique} (se produisant tous les \cle{2 à 7 ans}, sans régularité stricte) caractérisé par un \cle{réchauffement anormal des eaux superficielles du Pacifique équatorial Est}, accompagné d'un affaiblissement des alizés et d'un déplacement des zones de précipitations. La phase inverse, marquée par un refroidissement anormal de ces mêmes eaux, est appelée \cle{La Niña}.
\end{definition}

\begin{popfigure}
\begin{tikzpicture}[scale=0.92, font=\small]
% ===== Situation normale =====
\begin{scope}
% Océan
\fill[popblue!30] (0,0) rectangle (11,1.4);
% Eaux chaudes à l'ouest
\fill[poporange!45] (0,0.7) rectangle (4,1.4);
% Continents
\fill[popgold!50] (-1.6,0) rectangle (0,2.6);
\fill[popgold!50] (11,0) rectangle (12.6,2.6);
\node[popdark, rotate=90] at (-1.15,1.3) {Asie / Indonésie};
\node[popdark, rotate=90] at (12.15,1.3) {Pérou};
% Alizés
\draw[->, popgreen, line width=1.8pt] (10.5,2.3) -- (1.5,2.3) node[midway, above, popdark] {Alizés forts (E $\rightarrow$ O)};
% Upwelling
\draw[->, popblue, line width=1.6pt] (10.2,0.15) -- (10.2,1.15);
\node[popblue, left] at (10.0,0.65) {Upwelling d'eau froide};
% Pluies à l'ouest
\draw[popmuted, fill=popcream] (2,3.4) ellipse (1.0 and 0.4);
\foreach \x in {1.5,1.9,2.3,2.7} \draw[popblue, line width=1pt] (\x,2.95) -- (\x-0.12,2.5);
\node[popblue] at (2,3.95) {Pluies abondantes};
\node[popdark] at (2,0.35) {Eaux chaudes ($\approx$ \SI{28}{\celsius})};
\node[popdark] at (9,0.35) {Eaux froides ($\approx$ \SI{20}{\celsius})};
\node[font=\bfseries, popdark] at (5.5,-0.7) {Situation normale};
\end{scope}
% ===== El Niño =====
\begin{scope}[yshift=-5.6cm]
% Océan
\fill[popblue!30] (0,0) rectangle (11,1.4);
% Eaux chaudes étendues vers l'est
\fill[poporange!45] (0,0.7) rectangle (9.5,1.4);
% Continents
\fill[popgold!50] (-1.6,0) rectangle (0,2.6);
\fill[popgold!50] (11,0) rectangle (12.6,2.6);
\node[popdark, rotate=90] at (-1.15,1.3) {Asie / Indonésie};
\node[popdark, rotate=90] at (12.15,1.3) {Pérou};
% Alizés faibles
\draw[->, popgreen, dashed, line width=1.4pt] (8.5,2.3) -- (4.5,2.3) node[midway, above, popdark] {Alizés affaiblis};
% Upwelling bloqué
\draw[popblue, line width=1.6pt] (9.9,0.2) -- (10.5,0.8);
\draw[popblue, line width=1.6pt] (10.5,0.2) -- (9.9,0.8);
\node[popblue, left] at (9.7,0.5) {Upwelling bloqué};
% Pluies à l'est
\draw[popmuted, fill=popcream] (9,3.4) ellipse (1.0 and 0.4);
\foreach \x in {8.5,8.9,9.3,9.7} \draw[popblue, line width=1pt] (\x,2.95) -- (\x-0.12,2.5);
\node[popblue] at (9,3.95) {Pluies diluviennes};
% Sécheresse à l'ouest
\draw[poporange, fill=popyellow] (2,3.3) circle (0.3);
\node[poporange] at (2,3.95) {Sécheresse};
\node[popdark] at (4.5,0.35) {Eaux anormalement chaudes vers l'est};
\node[font=\bfseries, popdark] at (5.5,-0.7) {Année El Niño};
\end{scope}
\end{tikzpicture}
\legende{Comparaison du Pacifique équatorial en situation normale et en année El Niño. Normalement, les alizés accumulent les eaux chaudes et les pluies à l'ouest, tandis qu'un upwelling froid et riche en poissons alimente les côtes du Pérou. Lors d'El Niño, les alizés faiblissent, les eaux chaudes et les pluies se déplacent vers l'est, et l'upwelling est bloqué.}
\end{popfigure}

\textbf{Les conséquences d'El Niño à l'échelle mondiale.}\par
Parce que l'atmosphère et l'océan forment un système couplé, le déplacement des eaux chaudes du Pacifique perturbe la circulation atmosphérique de toute la planète (on parle de \cle{téléconnexions climatiques}) :

\begin{itemize}
\item \textbf{Sécheresses} en Indonésie, en Australie, en Inde (mousson affaiblie) et en \cle{Afrique australe} (Mozambique, Zimbabwe, Afrique du Sud), avec des récoltes compromises et des crises alimentaires ;
\item \textbf{Pluies diluviennes et inondations} sur les côtes normalement arides du \cle{Pérou} et de l'Équateur ;
\item \textbf{Effondrement des pêcheries} péruviennes (anchois), faute d'upwelling ;
\item Hivers plus doux dans certaines régions d'Amérique du Nord, modification des trajectoires des cyclones, etc.
\end{itemize}

\cle{La Niña}, phase froide de l'oscillation, produit des effets globalement inverses : renforcement des alizés, eaux de surface du Pacifique Est anormalement froides, pluies renforcées en Indonésie et sécheresses accrues sur les côtes américaines du Pacifique. L'ensemble El Niño / La Niña constitue l'oscillation australe (ENSO, \emph{El Niño -- Southern Oscillation}), la principale source de variabilité climatique d'année en année sur Terre.

\begin{attention}[Erreur fréquente : ce qu'El Niño n'est pas]
El Niño n'est \textbf{ni un courant marin permanent}, ni une \textbf{conséquence directe du réchauffement climatique}. C'est une \cle{oscillation naturelle} du système océan-atmosphère, connue et décrite depuis des siècles, qui alterne avec La Niña. Le réchauffement climatique pourrait modifier sa fréquence ou son intensité, mais il ne l'a pas créé. Autre piège : El Niño ne se produit pas dans l'Atlantique ni au Cameroun ; ses effets se font sentir à distance, par téléconnexions.
\end{attention}

\begin{savaistu}[Pourquoi « El Niño » ?]
Le nom « El Niño » signifie « l'Enfant » (l'enfant Jésus) en espagnol. Il vient des \cle{pêcheurs péruviens} qui, depuis des siècles, observaient l'arrivée des eaux chaudes côtières — et la disparition des poissons — autour de la période de \cle{Noël}. Ce phénomène saisonnier local a ensuite donné son nom au grand dérèglement planétaire que les scientifiques étudient aujourd'hui par satellites et par bouées ancrées dans tout le Pacifique (réseau d'observation TAO/TRITON).
\end{savaistu}

\courssub{Réchauffement climatique et interactions océan-atmosphère}

\textbf{Un océan qui chauffe, une atmosphère qui s'emballe.}\par
Puisque l'océan stocke l'essentiel de la chaleur du système climatique, le \cle{réchauffement climatique} le concerne au premier chef : plus de 90 \% de l'excès de chaleur piégé par les gaz à effet de serre est absorbé par les océans. Or, un océan plus chaud modifie les échanges avec l'atmosphère :

\begin{itemize}
\item \textbf{Plus d'évaporation} : une atmosphère plus chaude peut contenir davantage de vapeur d'eau (environ 7 \% de plus par degré de réchauffement, relation de Clausius-Clapeyron), ce qui alimente des \cle{précipitations plus intenses} et des inondations plus violentes ;
\item \textbf{Plus d'énergie disponible} pour les phénomènes extrêmes : les cyclones et tempêtes, qui puisent leur énergie dans la chaleur des eaux de surface (au-delà d'environ \SI{26}{\celsius}), peuvent devenir \cle{plus intenses} sur des océans plus chauds ;
\item \textbf{Élévation du niveau des mers} (dilatation thermique de l'eau et fonte des glaces), qui aggrave les submersions côtières lors des tempêtes — un risque réel pour les villes littorales comme Douala ou Limbé ;
\item Possible modification de la fréquence et de l'intensité des épisodes El Niño et La Niña, sujet encore débattu par les climatologues.
\end{itemize}

\begin{exoresolu}[Relier un phénomène à son origine dynamique]
Énoncé. En 2015-2016, l'Afrique australe a connu une sécheresse sévère ayant détruit les récoltes de maïs, tandis que le Pérou subissait des pluies torrentielles. (1) Identifier le phénomène climatique responsable. (2) Expliquer, à partir des échanges océan-atmosphère, l'origine de ces deux anomalies.
\tcblower
Solution. (1) Il s'agit d'un épisode \cle{El Niño} : le réchauffement anormal des eaux superficielles du Pacifique équatorial Est en 2015-2016 a été l'un des plus forts jamais enregistrés.
(2) Origine dynamique : les alizés se sont affaiblis, cessant de pousser les eaux chaudes de surface vers l'ouest du Pacifique. Les eaux chaudes se sont donc déplacées vers l'est, et avec elles la zone de convection et de pluies : les côtes du Pérou, habituellement sèches car surmontées d'eaux froides (upwelling), ont reçu des pluies diluviennes. À l'inverse, l'ouest du Pacifique et, par perturbation de la circulation atmosphérique globale (téléconnexions), l'Afrique australe ont connu un déficit de pluies : d'où la sécheresse. Ce même phénomène explique l'effondrement des pêches péruviennes, l'upwelling d'eau froide riche en nutriments étant bloqué par la couche d'eau chaude superficielle.
\end{exoresolu}

\begin{aretenir}[L'essentiel de la section]
\begin{itemize}
\item L'océan est un immense \cle{réservoir de chaleur et d'humidité} : il fournit la vapeur d'eau des précipitations et régule le climat grâce à sa capacité de stockage thermique (environ 1000 fois celle de l'atmosphère).
\item Les échanges sont réciproques : l'océan chauffe et humidifie l'atmosphère ; les vents entraînent les courants marins.
\item La \cle{mousson} est un vent saisonnier qui s'inverse entre l'été (océan $\rightarrow$ continent, humide, pluies) et l'hiver (continent $\rightarrow$ océan, sec : harmattan en Afrique de l'Ouest), à cause du contraste thermique continent/océan.
\item \cle{El Niño} est un réchauffement périodique (tous les 2 à 7 ans) des eaux du Pacifique équatorial Est, lié à l'affaiblissement des alizés ; il provoque sécheresses (Indonésie, Afrique australe) et pluies diluviennes (Pérou). \cle{La Niña} est sa phase froide inverse.
\item Le réchauffement climatique, en chauffant les océans, augmente l'énergie et l'humidité disponibles pour les phénomènes extrêmes.
\end{itemize}
\end{aretenir}

Ces interactions montrent que le système climatique fonctionne comme un tout : un affaiblissement des alizés dans le Pacifique peut provoquer une famine en Afrique australe. Dans la section suivante, nous verrons comment ces mêmes mécanismes — chaleur océanique, vents, contrastes de pression — engendrent les \cle{phénomènes climatiques dangereux} : tempêtes, cyclones et inondations.

\coursec{Les phénomènes climatiques dangereux}

Dans les sections précédentes, nous avons vu que l'océan et l'atmosphère forment un couple indissociable : l'océan chauffe et humidifie l'air, l'atmosphère entraîne les eaux de surface par les vents, et des phénomènes comme \emph{El Niño} montrent que toute anomalie de l'un se répercute sur l'autre. Mais ce moteur thermique planétaire peut aussi se dérégler localement et produire des événements d'une violence extrême. En mars 2019, le cyclone Idai s'abattait sur le Mozambique, faisant plus de 1\,300 morts ; chaque année, les pluies torrentielles transforment certains quartiers de Yaoundé et de Douala en véritables marécages ; en octobre 2019, un glissement de terrain à Bafoussam ensevelissait des habitations endormies. Tous ces drames ont une même origine : les \cle{mouvements de l'atmosphère et de l'océan} lorsqu'ils deviennent extrêmes. L'objectif de cette section est de comprendre \emph{comment} naissent ces phénomènes dangereux, afin de mieux s'en protéger.

\courssub{Les tempêtes et les cyclones tropicaux}

\textbf{Observation et questionnement.} Sur les images satellites, on observe parfois au-dessus des océans tropicaux d'immenses spirales de nuages, larges de plusieurs centaines de kilomètres, avec un petit trou clair au centre. D'où vient cette rotation ? Pourquoi ces monstres naissent-ils toujours au-dessus des mers chaudes, et jamais au-dessus du golfe de Guinée ?

\begin{definition}[Cyclone tropical]
Un \cle{cyclone tropical} est une \textbf{dépression très creuse} (pression centrale parfois inférieure à 950 hPa) qui se forme au-dessus des océans chauds intertropicaux. Il se caractérise par des \textbf{vents tournants en spirale} (rotation dite \cle{cyclonique}) dépassant \SI{118}{km/h} (soit 64 nœuds) et des \textbf{pluies torrentielles}. En dessous de ce seuil, on parle de tempête tropicale (vents 63--117 km/h) ou de dépression tropicale (vents < 63 km/h).
\end{definition}

\textbf{Les conditions de formation.} Un cyclone est une véritable machine thermique : il puise son énergie dans la chaleur de l'océan. Cinq conditions doivent être réunies simultanément :

\begin{enumerate}
\item une \textbf{eau de mer à plus de \num{26,5} \textcelsius{}} sur une grande épaisseur (au moins 50 m), qui fournit une \textbf{forte évaporation} : la vapeur d'eau est le « carburant » du cyclone ;
\item une \textbf{convergence des alizés} en surface (zone de convergence intertropicale ou perturbation ondulatoire), qui force l'air chaud et humide à s'élever ;
\item en s'élevant, la vapeur d'eau se \textbf{condense} en nuages : cette condensation \textbf{libère de la chaleur latente}, qui réchauffe encore l'air, accélère l'ascendance et creuse la dépression — c'est un cercle auto-entretenu (retroaction positive) ;
\item une \textbf{force de Coriolis suffisante} pour mettre l'air convergeant en rotation : elle n'est efficace qu'au-delà d'environ \textbf{5° de latitude} de part et d'autre de l'équateur ;
\item un \textbf{faible cisaillement vertical du vent} (différence de vitesse/direction entre le vent en surface et en altitude) : un cisaillement fort « décapite » le cyclone naissant et l'empêche de s'organiser verticalement.
\end{enumerate}

\begin{propriete}[Impossibilité de formation à l'équateur]
Un cyclone tropical \textbf{ne peut pas se former à l'équateur} (ni le traverser), car la force de Coriolis y est nulle. Justification guidée : (1) l'air en surface converge vers le centre dépressionnaire ; (2) la force de Coriolis dévie ce mouvement (vers la droite dans l'hémisphère Nord, vers la gauche dans l'hémisphère Sud) ; (3) cette déviation transforme la convergence simple en \textbf{rotation cyclonique} ; (4) à l'équateur, Coriolis étant nul, l'air convergeant remplit directement la dépression sans tourner : la dépression se comble et meurt. C'est pourquoi le \textbf{golfe de Guinée est une zone sans cyclones}.
\end{propriete}

\textbf{La structure d'un cyclone.} Un cyclone mûr présente une organisation remarquable, visible en coupe verticale :

\begin{itemize}
\item l'\cle{œil} : zone centrale de 20 à 50 km de diamètre, où l'air \textbf{descend} par subsidence adiabatique ; le temps y est étonnamment \textbf{calme}, le ciel parfois dégagé, la pression minimale ;
\item le \cle{mur de l'œil} : anneau de nuages gigantesques (cumulonimbus) entourant l'œil ; c'est là que règnent les \textbf{vents et les pluies les plus violents} (ascendances maximales, gradient de pression maximal) ;
\item les \cle{bandes spiralées} : bras nuageux qui s'enroulent vers le centre, apportant par rafales successives pluies et vents violents ;
\item en altitude (vers 12--15 km), l'air \textbf{diverge} et s'évacue vers l'extérieur, ce qui aspire l'air de surface et entretient la dépression (pompage).
\end{itemize}

\begin{popfigure}
\begin{center}
\begin{tikzpicture}[scale=1.0]
% sol océan
\fill[popblue!25] (-6.2,0) rectangle (6.2,-0.5);
\draw[thick,popblue] (-6.2,0) -- (6.2,0);
\node[popdark] at (-5.2,-0.28) {\small Océan chaud ($>26.5$~\textcelsius)};
% mur de l'oeil (nuages)
\fill[popmuted!40] (-2.6,0) -- (-2.2,0) .. controls (-1.9,2.5) and (-1.6,4.0) .. (-1.2,4.6) -- (-2.0,4.6) .. controls (-2.6,3.5) and (-2.9,1.5) .. (-2.6,0);
\fill[popmuted!40] (2.6,0) -- (2.2,0) .. controls (1.9,2.5) and (1.6,4.0) .. (1.2,4.6) -- (2.0,4.6) .. controls (2.6,3.5) and (2.9,1.5) .. (2.6,0);
% bandes spiralées
\fill[popmuted!25] (-5.8,0.2) .. controls (-5.0,1.2) and (-4.2,1.6) .. (-3.4,1.5) .. controls (-4.2,0.9) and (-5.0,0.5) .. (-5.8,0.2);
\fill[popmuted!25] (5.8,0.2) .. controls (5.0,1.2) and (4.2,1.6) .. (3.4,1.5) .. controls (4.2,0.9) and (5.0,0.5) .. (5.8,0.2);
\fill[popmuted!25] (-5.2,1.6) .. controls (-4.6,2.3) and (-3.9,2.6) .. (-3.2,2.5) .. controls (-3.9,2.0) and (-4.6,1.7) .. (-5.2,1.6);
\fill[popmuted!25] (5.2,1.6) .. controls (4.6,2.3) and (3.9,2.6) .. (3.2,2.5) .. controls (3.9,2.0) and (4.6,1.7) .. (5.2,1.6);
% oeil : air descendant
\draw[->,very thick,popblue] (0,4.2) -- (0,1.0);
\node[popblue,rotate=90] at (0.4,2.6) {\small Subsidence};
% ascendances dans le mur
\draw[->,very thick,poporange] (-2.35,0.3) .. controls (-2.1,2.0) .. (-1.6,4.2);
\draw[->,very thick,poporange] (2.35,0.3) .. controls (2.1,2.0) .. (1.6,4.2);
% convergence en surface
\draw[->,thick,popgreen] (-5.6,0.25) -- (-3.0,0.25);
\draw[->,thick,popgreen] (5.6,0.25) -- (3.0,0.25);
% divergence en altitude
\draw[->,thick,popgreen] (-1.2,4.9) -- (-4.4,4.9);
\draw[->,thick,popgreen] (1.2,4.9) -- (4.4,4.9);
% légendes pointées
\draw[thin,popdark] (0,2.4) -- (0.15,3.3) -- (1.0,3.3) node[right,popdark]{\small Œil (air descendant, calme)};
\draw[thin,popdark] (-2.35,2.6) -- (-3.6,3.6) node[left,popdark]{\small Mur de l'œil (vents/pluies max)};
\draw[thin,popdark] (4.6,1.9) -- (6.4,2.6) node[right,popdark]{\small Bandes spiralées};
\node[popgreen!60!popdark] at (-4.3,-0.75) {\small Convergence en surface};
\node[popgreen!60!popdark] at (0,5.25) {\small Divergence en altitude};
\end{tikzpicture}
\end{center}
\legende{Coupe verticale schématique d'un cyclone tropical. L'air chaud et humide converge en surface (flèches vertes), s'élève violemment dans le mur de l'œil (flèches orange), diverge en altitude, tandis que l'air descend doucement dans l'œil (flèche bleue), zone de calme relatif et de pression minimale.}
\end{popfigure}

\begin{attention}[Le calme trompeur de l'œil]
Erreur fréquente : croire que le centre du cyclone est la zone la plus dangereuse. C'est l'inverse ! Dans l'\textbf{œil}, le temps est \textbf{calme} et le ciel peut être dégagé. Le danger maximal se situe dans le \textbf{mur de l'œil}. Conséquence dramatique : lorsque l'œil passe sur une ville, les habitants croient la fin du danger et sortent... puis la seconde moitié du mur les frappe, avec des vents soufflant en sens inverse. Il faut rester abrité jusqu'à l'annonce officielle de fin d'alerte.
\end{attention}

\textbf{Nomenclature : un même phénomène, trois noms.} Selon le bassin océanique où il se forme, le cyclone tropical porte un nom différent :

\begin{center}
\begin{tabularx}{\linewidth}{|l|X|}\hline
\rowcolor{popblueL} \textbf{Nom} & \textbf{Zone de formation} \\ \hline
Cyclone (tropical) & Océan Indien, Pacifique Sud \\ \hline
Ouragan & Atlantique Nord, Pacifique Nord-Est \\ \hline
Typhon & Pacifique Nord-Ouest \\ \hline
\end{tabularx}
\end{center}

Il s'agit \textbf{rigoureusement du même phénomène météorologique} ; seul le vocabulaire régional change.

\textbf{L'échelle de Saffir-Simpson.} Pour classer l'intensité des cyclones, les météorologues utilisent l'échelle de Saffir-Simpson, fondée sur la vitesse des vents soutenus (moyenne sur 1 minute) :

\begin{center}
\begin{tabularx}{\linewidth}{|c|X|X|}\hline
\rowcolor{popblueL} \textbf{Catégorie} & \textbf{Vents (km/h)} & \textbf{Dégâts typiques} \\ \hline
1 & 119 -- 153 & Légers : toitures, arbres, panneaux \\ \hline
2 & 154 -- 177 & Modérés : toitures arrachées, coupures d'électricité prolongées \\ \hline
3 & 178 -- 208 & Importants : maisons endommagées, inondations côtières \\ \hline
4 & 209 -- 251 & Extrêmes : destructions majeures, zones inhabitables \\ \hline
5 & $>$ 251 & Catastrophiques : destruction totale de bâtiments \\ \hline
\end{tabularx}
\end{center}

\begin{exemplebox}[Le cyclone Idai (mars 2019) — étude de cas]
Le cyclone Idai, formé dans le canal du Mozambique, a frappé le Mozambique (ville de Beira), le Zimbabwe et le Malawi en mars 2019. Bilan : plus de \textbf{1\,300 morts}, des vents dépassant \SI{195}{km/h} (catégorie 3--4), des pluies diluviennes ayant provoqué une « mer intérieure » de près de 50 km de large autour de Beira, et des \textbf{centaines de milliers de déplacés}. C'est l'une des pires catastrophes climatiques de l'hémisphère Sud — et un rappel que l'Afrique australe et orientale, contrairement au golfe de Guinée, est exposée aux cyclones de l'océan Indien.
\end{exemplebox}

\begin{popfigure}
\begin{center}
\begin{tikzpicture}[scale=0.92]
% continents stylisés
\draw[fill=popgreen!20,draw=popgreen!50!popdark] (-5.6,0.2) .. controls (-5.2,1.6) and (-4.0,2.0) .. (-3.2,1.4) .. controls (-2.8,0.8) and (-3.4,0.2) .. (-3.6,-0.6) .. controls (-3.9,-1.6) and (-4.6,-2.0) .. (-5.0,-1.2) .. controls (-5.4,-0.6) and (-5.9,-0.4) .. (-5.6,0.2); % Amériques
\draw[fill=popgreen!20,draw=popgreen!50!popdark] (-2.4,1.2) .. controls (-1.6,1.9) and (-0.4,1.9) .. (0.2,1.3) .. controls (0.6,0.8) and (0.2,0.4) .. (-0.2,-0.2) .. controls (-0.5,-1.0) and (-0.3,-1.9) .. (0.1,-2.2) .. controls (-0.6,-2.3) and (-1.2,-1.6) .. (-1.4,-0.8) .. controls (-1.7,0.0) and (-2.8,0.5) .. (-2.4,1.2); % Europe-Afrique
\draw[fill=popgreen!20,draw=popgreen!50!popdark] (0.6,1.4) .. controls (1.8,2.0) and (3.4,1.8) .. (4.0,1.0) .. controls (4.4,0.4) and (3.6,0.0) .. (3.2,-0.5) .. controls (2.8,-1.0) and (3.4,-1.2) .. (3.0,-1.6) .. controls (2.2,-1.4) and (1.6,-0.6) .. (1.4,0.2) .. controls (1.2,0.8) and (0.2,0.9) .. (0.6,1.4); % Asie
\draw[fill=popgreen!20,draw=popgreen!50!popdark] (2.6,-1.6) .. controls (3.2,-1.2) and (4.0,-1.3) .. (4.2,-1.8) .. controls (3.8,-2.3) and (3.0,-2.3) .. (2.6,-1.6); % Australie
% équateur et tropiques
\draw[dashed,popdark] (-6.3,0) -- (6.3,0) node[right]{\scriptsize Équateur};
\draw[dotted,popmuted] (-6.3,1.15) -- (6.3,1.15) node[right]{\scriptsize Cancer};
\draw[dotted,popmuted] (-6.3,-1.15) -- (6.3,-1.15) node[right]{\scriptsize Capricorne};
% zones cycloniques (ovales rouges)
\draw[fill=poporange!45,draw=poporange,thick] (-4.6,0.75) ellipse (0.9 and 0.35);
\draw[fill=poporange!45,draw=poporange,thick] (-4.9,-0.8) ellipse (0.7 and 0.3);
\draw[fill=poporange!45,draw=poporange,thick] (0.9,0.85) ellipse (0.8 and 0.32);
\draw[fill=poporange!45,draw=poporange,thick] (1.3,-0.85) ellipse (0.9 and 0.32);
\draw[fill=poporange!45,draw=poporange,thick] (3.4,0.8) ellipse (0.9 and 0.35);
\draw[fill=poporange!45,draw=poporange,thick] (3.6,-0.85) ellipse (0.8 and 0.3);
% golfe de Guinée exempt
\draw[fill=popblue!35,draw=popblue,thick] (-0.9,0.25) ellipse (0.55 and 0.3);
\draw[thin,popdark] (-0.9,0.25) -- (-2.2,-1.3) node[below,align=center,popdark]{\scriptsize Golfe de Guinée :\\\scriptsize zone SANS cyclone};
\draw[thin,popdark] (-4.6,1.15) -- (-4.6,1.9) node[above,popdark]{\scriptsize Ouragans};
\draw[thin,popdark] (3.4,1.2) -- (3.4,1.9) node[above,popdark]{\scriptsize Typhons};
\draw[thin,popdark] (1.3,-1.2) -- (1.3,-2.1) node[below,popdark]{\scriptsize Cyclones (océan Indien)};
\end{tikzpicture}
\end{center}
\legende{Carte mondiale schématique des zones de formation des cyclones tropicaux (ovales orange). Ils naissent au-dessus des eaux chaudes, entre environ 5° et 25° de latitude. La bande équatoriale (entre les pointillés proches de l'équateur), dont le golfe de Guinée (ovale bleu), en est exempte car la force de Coriolis y est trop faible.}
\end{popfigure}

\textbf{Pourquoi le Cameroun est-il épargné par les cyclones ?} Deux raisons se combinent : (1) le Cameroun est situé à \textbf{proximité immédiate de l'équateur} (entre 2° et 13° N environ), où la force de Coriolis est quasi nulle : l'air ne peut pas acquérir de rotation cyclonique ; (2) les eaux du \textbf{golfe de Guinée} sont relativement \textbf{moins chaudes} que celles des grands bassins cycloniques, en partie à cause du courant froid de Benguela qui remonte le long des côtes africaines (remontée d'eau profonde). Le pays connaît donc des tempêtes et des pluies violentes, mais jamais de cyclones tropicaux.

\begin{savaistu}[Une énergie colossale]
Un cyclone tropical mûr libère en une seule journée, par condensation de la vapeur d'eau, une énergie équivalente à la \textbf{consommation électrique mondiale de plusieurs mois} (environ $2 \times 10^{13}$ W). C'est aussi ce qui le condamne : dès qu'il touche terre ou passe sur des eaux plus froides ($< 26,5$ \textcelsius{}), il est privé de son carburant (la vapeur d'eau chaude) et s'affaiblit rapidement — mais ses pluies torrentielles continuent de provoquer inondations et glissements de terrain loin à l'intérieur des terres.
\end{savaistu}

\courssub{Les inondations}

\textbf{Situation concrète.} Quartier Nkolondom à Yaoundé ou Makepe à Douala, un soir de pluie diluvienne : en moins d'une heure, l'eau monte dans les maisons, emporte véhicules et marchandises, coupe les routes. Ce scénario, vécu chaque année par des milliers de Camerounais, illustre le phénomène d'\cle{inondation}.

\begin{definition}[Inondation]
Une \cle{inondation} est la submersion temporaire, par de l'eau, de zones habituellement émergées. Elle devient une \textbf{catastrophe} lorsqu'elle affecte des populations, des cultures ou des infrastructures.
\end{definition}

\textbf{Les causes des inondations.} Elles sont souvent multiples et combinées :

\begin{itemize}
\item les \textbf{pluies intenses et prolongées} : le sol saturé ne peut plus absorber l'eau, qui ruisselle (c'est la cause principale au Cameroun, surtout pendant la grande saison des pluies) ;
\item les \textbf{crues} : le débit d'un cours d'eau dépasse la capacité de son lit, qui déborde dans la plaine inondable (ex. fleuve Sanaga, Logone et Chari dans l'Extrême-Nord) ;
\item les \textbf{ondes de tempête} (ou \cle{surcotes}) côtières : les vents violents d'une tempête ou d'un cyclone poussent l'eau de mer vers la côte et élèvent brutalement le niveau marin de plusieurs mètres, submergeant les plaines côtières basses ;
\item l'\textbf{urbanisation et l'imperméabilisation des sols} : le béton et le bitume empêchent l'infiltration ; l'eau ruisselle instantanément vers les points bas ;
\item l'\textbf{obstruction des caniveaux} par les déchets plastiques et les constructions anarchiques sur les axes d'écoulement, qui bloquent l'évacuation des eaux.
\end{itemize}

\begin{popfigure}
\begin{center}
\begin{tikzpicture}[scale=1.0]
% fond marin
\fill[popblue!20] (-6.3,0) rectangle (6.3,-1.6);
% niveau normal
\draw[thick,popblue] (-6.3,0) -- (6.3,0);
\node[left,popblue] at (-6.3,0) {\scriptsize Niveau normal};
% dôme de surcote
\fill[popblue!55] plot[smooth] coordinates {(-6.3,0) (-4,0.1) (-2,0.7) (0,1.3) (1.6,1.5) (3.0,1.1) (4.2,0.5) (4.8,0.2) (6.3,0)} -- (6.3,0) -- cycle;
\draw[thick,popblue!80!popdark] plot[smooth] coordinates {(-6.3,0) (-4,0.1) (-2,0.7) (0,1.3) (1.6,1.5) (3.0,1.1) (4.2,0.5) (4.8,0.2) (6.3,0)};
% côte basse
\fill[popgold!40] (4.8,0) -- (6.3,0) -- (6.3,0.9) -- (5.6,0.55) -- (5.0,0.2) -- cycle;
\draw[thick,popgold!60!popdark] (4.8,0.05) -- (5.0,0.2) -- (5.6,0.55) -- (6.3,0.9);
% maisons
\draw[fill=popcream,draw=popdark] (5.15,0.35) rectangle (5.55,0.75);
\draw[popdark] (5.1,0.75) -- (5.35,0.95) -- (5.6,0.75);
\draw[fill=popcream,draw=popdark] (5.75,0.62) rectangle (6.1,0.98);
\draw[popdark] (5.71,0.98) -- (5.925,1.16) -- (6.14,0.98);
% vent
\draw[->,very thick,poporange] (-5.2,2.2) -- (2.2,2.2);
\draw[->,very thick,poporange] (-4.4,2.7) -- (3.0,2.7);
\node[poporange] at (-1.2,3.0) {\small Vents violents poussant l'eau vers la côte};
% flèche submersion
\draw[->,thick,popblue!80!popdark] (4.0,0.9) .. controls (4.8,1.1) .. (5.4,1.05);
% légendes
\draw[thin,popdark] (1.6,1.5) -- (1.6,2.3) node[above,popdark]{\small Surcote : dôme d'eau élevé par le vent et la basse pression};
\node[popdark] at (5.6,-0.5) {\scriptsize Plaine côtière basse submergée};
\end{tikzpicture}
\end{center}
\legende{Coupe schématique d'une onde de tempête (surcote). Les vents violents et la très basse pression d'un cyclone soulèvent un dôme d'eau qui, en atteignant une côte basse, submerge les habitations. Ce mécanisme a causé l'essentiel des morts du cyclone Idai à Beira.}
\end{popfigure}

\begin{exemplebox}[Inondations et glissements de terrain au Cameroun]
À \textbf{Yaoundé} et \textbf{Douala}, la conjugaison de pluies intenses, de l'imperméabilisation des sols, de caniveaux obstrués par les déchets et de constructions dans les zones marécageuses provoque des inondations récurrentes, avec pertes en vies humaines et destructions de biens. En octobre 2019, après des pluies exceptionnelles, un \textbf{glissement de terrain à Bafoussam} (quartier Gouache) a enseveli des habitations et fait plus de 40 morts : le sol, gorgé d'eau, a vu sa cohésion diminuer et son poids augmenter, déclenchant une rupture de pente. L'eau de pluie est ici le facteur déclenchant d'une catastrophe géologique.
\end{exemplebox}

\courssub{Autres phénomènes dangereux liés à la circulation atmosphérique}

\begin{itemize}
\item \textbf{Les tornades} : colonnes d'air en rotation très rapide (jusqu'à 500 km/h), étroites (quelques dizaines à centaines de mètres) et de courte durée (minutes à une heure), descendant d'un nuage d'orage (cumulonimbus) très organisé (supercellule). Plus violentes mais beaucoup plus localisées qu'un cyclone, elles sont fréquentes aux États-Unis (\emph{Tornado Alley}) et rares au Cameroun.
\item \textbf{Les sécheresses} : elles résultent d'anomalies durables de la circulation atmosphérique (déplacement des zones de haute pression, affaiblissement de la mousson, phénomène \emph{El Niño} / \emph{La Niña}) qui privent une région de ses pluies habituelles. Au Sahel et dans l'Extrême-Nord camerounais, elles menacent l'agriculture, l'élevage et la sécurité alimentaire.
\item \textbf{Les tempêtes de sable du harmattan} : en saison sèche (décembre à février), le vent continental de nord-est (harmattan) transporte des poussières du Sahara jusqu'au golfe de Guinée. Visibilité réduite (danger pour la navigation aérienne et routière), sécheresse de la peau et des voies respiratoires, irritation des yeux, impact sur la photosynthèse : c'est un phénomène météorologique dangereux bien connu des Camerounais.
\end{itemize}

\courssub{Prévention et protection face aux catastrophes}

\textbf{Savoir-faire visé :} Proposer des mesures de protection face à ces phénomènes. La gestion du risque suit le cycle : \textbf{Prévention} $\rightarrow$ \textbf{Préparation} $\rightarrow$ \textbf{Réponse} $\rightarrow$ \textbf{Rétablissement}.

\begin{methode}[Relier un phénomène dangereux à son origine dynamique]
Face à un document (carte, schéma, article) décrivant une catastrophe, construis toujours la chaîne causale en quatre maillons :
\begin{enumerate}
\item \textbf{Source d'énergie} : eau de mer chaude, contraste thermique fort, sol saturé d'eau ;
\item \textbf{Mouvement de l'air ou de l'eau} : convergence, ascendance, rotation (Coriolis), vent poussant la mer ;
\item \textbf{Phénomène} : vents violents, pluies torrentielles, surcote, crue ;
\item \textbf{Dégâts} : inondations, destructions, glissements de terrain, déplacements de populations.
\end{enumerate}
Exemple appliqué au cyclone Idai : eau du canal du Mozambique à près de 30 \textcelsius{} $\rightarrow$ évaporation, convergence et rotation des alizés $\rightarrow$ cyclone de catégorie 4 avec pluies diluviennes $\rightarrow$ inondation de Beira, plus de 1\,300 morts.
\end{methode}

\begin{propriete}[Mesures de prévention et de protection (niveau individuel et collectif)]
\begin{itemize}
\item \textbf{Aménagement du territoire} : respect des zones non \ae{}dificandi (plaines inondables, pentes instables, littoral bas) ; restauration des zones humides et mangroves (éponges naturelles, brise-vagues) ; drainage urbain efficace et entretenu (curage régulier des caniveaux) ;
\item \textbf{Surveillance et alerte précoce} : stations météorologiques, radars, satellites (Météo-Cameroun, centres régionaux comme ACMAD) ; diffusion des alertes (radio, SMS, sirènes) ; plans d'urgence communaux (zones de refuge, évacuation) ;
\item \textbf{Adaptation des constructions} : maisons sur pilotis en zone inondable, toitures arrondies et ancrées au sol en zone cyclonique (Antilles, Réunion), matériaux résistants ;
\item \textbf{Gestion des déchets} : lutte contre les dépôts sauvages qui bouchent les exutoires ; tri et recyclage ;
\item \textbf{Éducation et sensibilisation} : exercices de simulation (évacuation), connaissance des consignes (ne pas traverser une route inondée, couper l'électricité, ne pas sortir dans l'œil du cyclone).
\end{itemize}
\end{propriete}

\begin{exemplebox}[Exemple camerounais : Plan de contingence inondations]
Le Cameroun, avec l'appui de partenaires (PNUD, Croix-Rouge), a élaboré des \textbf{plans de contingence} au niveau national et départemental. Ils prévoient : la cartographie des zones à risque (ex. lit majeur de la Bénoué, bas-fonds de Yaoundé), la mise en place de comités de veille villageois, le pré-positionnement de stocks (kits d'hygiène, moustiquaires, tentes) et la définition de circuits d'alerte (chef de village $\rightarrow$ sous-préfet $\rightarrow$ préfet $\rightarrow$ MINATD/MINEPDED). La participation des populations (savoir-être : vigilance, solidarité) est la clé de l'efficacité.
\end{exemplebox}

\begin{exoresolu}[Expliquer l'absence de cyclones au Cameroun]
\textbf{Énoncé.} Un élève affirme : « Le Cameroun a des côtes sur l'océan Atlantique tropical, il devrait donc subir des ouragans comme les Antilles. » À l'aide de vos connaissances, expliquez pourquoi cette affirmation est fausse (deux arguments attendus).
\tcblower
\textbf{Solution détaillée.}
\begin{enumerate}
\item \textbf{Argument de latitude (Coriolis).} Les côtes camerounaises bordent le golfe de Guinée, à moins de 5° de latitude de l'équateur (entre 2° et 4° N environ). Or la force de Coriolis, qui met en rotation l'air convergeant vers une dépression, y est quasi nulle : sans rotation, aucune dépression ne peut s'organiser en cyclone tropical.
\item \textbf{Argument thermique.} La formation d'un cyclone exige une eau de mer à plus de \num{26,5} \textcelsius{} sur une grande épaisseur (50 m). Les eaux du golfe de Guinée sont relativement moins chaudes en surface, notamment sous l'influence de la remontée d'eau froide (upwelling) liée au courant de Benguela et à la divergence équatoriale : la source d'énergie du cyclone fait donc défaut.
\end{enumerate}
\textbf{Conclusion.} L'affirmation est fausse : la proximité de l'équateur (Coriolis nul) et la température insuffisante des eaux expliquent que le Cameroun soit épargné par les cyclones tropicaux, contrairement aux Antilles situées entre 10° et 25° N au-dessus d'une mer des Caraïbes très chaude.
\end{exoresolu}

\begin{exercice}[Analyse d'un document : crue du Logone]
\textbf{Document.} Extrait d'un rapport OCHA 2022 : « Dans l'Extrême-Nord, les pluies exceptionnelles de septembre ont fait déborder le Logone et le Chari. 150\,000 personnes affectées, 25\,000 hectares de cultures détruits, 12\,000 maisons effondrées. Les digues de protection ont cédé par endroits. »
\begin{enumerate}
\item Identifie le phénomène climatique dangereux décrit.
\item En te basant sur la chaîne causale (Méthode), explique l'enchaînement : pluviométrie $\rightarrow$ crue $\rightarrow$ inondation $\rightarrow$ dégâts.
\item Propose deux mesures de prévention structurelles et deux mesures non structurelles adaptées à ce contexte sahélien.
\end{enumerate}
\end{exercice}

\begin{corrige}[Exercice : Crue du Logone]
\begin{enumerate}
\item Phénomène : \textbf{inondation par crue lente} (débordement de fleuves) suite à des pluies exceptionnelles.
\item Chaîne causale : \textbf{Source} = pluies exceptionnelles (mousson renforcée / anomalie circulation) sur le bassin versant (Adamoua, RCA, Tchad) $\rightarrow$ \textbf{Mouvement} = ruissellement massif, concentration dans le réseau hydrographique, onde de crue descendant le Logone/Chari $\rightarrow$ \textbf{Phénomène} = débit $>$ capacité du lit, rupture de digues, submersion de la plaine inondable (yaérés) $\rightarrow$ \textbf{Dégâts} = habitats en banco effondrés, cultures (riz, sorgho) perdues, populations déplacées, risque choléra.
\item Mesures : \textbf{Structurelles} : (1) Rehaussement/renforcement des digues existantes ; (2) Création de zones d'expansion de crue (bassins de rétention) en amont. \textbf{Non structurelles} : (1) Système d'alerte précoce basé sur la pluviométrie en amont (radar, stations) et diffusion par radio locale/SMS ; (2) Plan d'occupation des sols interdisant l'habitat permanent dans le lit majeur, promotion de l'agriculture récessive adaptée.
\end{enumerate}
\end{corrige}

\begin{aretenir}[L'essentiel de la section]
\begin{itemize}
\item Un cyclone tropical est une dépression très creuse des zones intertropicales : vents en spirale $>$ \SI{118}{km/h}, pluies torrentielles.
\item Il naît au-dessus d'une mer à plus de \num{26,5} \textcelsius{}, avec convergence des alizés, force de Coriolis suffisante (au-delà de 5° de latitude) et faible cisaillement vertical — jamais à l'équateur.
\item Structure : œil calme au centre, mur de l'œil (danger maximal), bandes spiralées, divergence en altitude.
\item Cyclone = ouragan = typhon : même phénomène, noms différents selon l'océan ; intensité classée par l'échelle de Saffir-Simpson (catégories 1 à 5).
\item Le Cameroun est épargné par les cyclones (équateur + eaux du golfe de Guinée moins chaudes), mais subit inondations, glissements de terrain, sécheresses (Extrême-Nord) et tempêtes de sable (harmattan).
\item Les inondations résultent de causes combinées : pluies extrêmes, crues, surcotes, imperméabilisation des sols et obstruction des caniveaux.
\item La prévention repose sur l'aménagement du territoire, l'alerte précoce, l'adaptation des constructions, la gestion des déchets et l'éducation des populations.
\end{itemize}
\end{aretenir}

\coursec{Prévention et protection face aux catastrophes}

Nous avons vu dans les sections précédentes que les mouvements de l'atmosphère et des océans engendrent des phénomènes dangereux : cyclones, tempêtes, inondations, sécheresses. Ces phénomènes sont \emph{naturels} : aucune technologie ne peut empêcher un cyclone de se former ou une pluie diluvienne de tomber sur Douala ou Yaoundé. Pourtant, deux pays frappés par un cyclone de même intensité ne comptent pas le même nombre de victimes. Pourquoi ? Parce que la catastrophe n'est pas seulement l'affaire de la nature : elle dépend aussi de la manière dont les sociétés se sont préparées. C'est l'objet de cette dernière section : comprendre comment l'être humain peut \cle{prévoir}, \cle{prévenir} et se \cle{protéger} face aux catastrophes liées aux mouvements atmosphériques et océaniques.

\courssub{Comprendre le risque : aléa, vulnérabilité et enjeu}

\textbf{Une situation concrète.} En octobre 2019, de fortes pluies s'abattent sur l'ouest du Cameroun. À Bafoussam, un glissement de terrain emporte des habitations construites sur un versant escarpé : des dizaines de personnes périssent. La même pluie tombe sur des quartiers voisins situés sur des plateaux stables : aucune victime. Le phénomène météorologique est identique ; les conséquences sont radicalement différentes. La différence ne tient pas au hasard, mais à la \emph{localisation} et à la \emph{fragilité} des populations exposées.

\begin{definition}[Le risque de catastrophe]
Un \cle{risque} de catastrophe résulte de la combinaison de trois éléments :
\begin{itemize}
\item l'\cle{aléa} : le phénomène dangereux lui-même (cyclone, inondation, glissement de terrain, onde de tempête), caractérisé par son intensité et sa probabilité d'occurrence ;
\item la \cle{vulnérabilité} : la fragilité des populations et des constructions (habitats précaires, constructions en zone inondable ou sur versant instable, pauvreté, manque d'information) ;
\item l'\cle{enjeu} : ce qui est exposé au danger (vies humaines, habitations, cultures, infrastructures, activités économiques).
\end{itemize}
On écrit souvent la relation fondamentale :
\[
\textbf{Risque} = \text{Aléa} \times \text{Vulnérabilité} \times \text{Enjeu}
\]
Si l'un des trois facteurs est nul (aucune construction en zone inondable, par exemple), le risque est nul, même si l'aléa est fort.
\end{definition}

\begin{exemplebox}[Deux quartiers, deux destins]
Considérons deux quartiers d'une ville camerounaise soumis aux mêmes pluies de mousson. Le quartier A est installé sur un bas-fond marécageux, avec des habitations en matériaux légers et des caniveaux obstrués : aléa fort + vulnérabilité forte + enjeu élevé = \textbf{catastrophe probable}. Le quartier B, sur un plateau drainé, avec des constructions conformes et des habitants informés : même aléa, mais vulnérabilité faible = \textbf{dégâts limités}. Réduire le risque, c'est donc agir sur la vulnérabilité et l'exposition des enjeux, puisqu'on ne peut pas supprimer l'aléa.
\end{exemplebox}

\courssub{Les trois piliers de la gestion des risques}

Face à un risque identifié, l'action des sociétés repose sur trois piliers complémentaires :

\begin{itemize}
\item la \cle{prévision} : anticiper le phénomène grâce à l'observation scientifique (météorologie, surveillance des océans) et diffuser des alertes ;
\item la \cle{prévention} : réduire durablement la vulnérabilité et l'exposition des enjeux (aménagement du territoire, constructions adaptées, information des populations) ;
\item la \cle{protection} : mettre les personnes à l'abri lorsque le phénomène survient (évacuation, abris, secours).
\end{itemize}

\begin{methode}[Raisonner en trois temps : avant, pendant, après]
Pour proposer des mesures de protection face à un phénomène climatique dangereux (compétence exigible au Probatoire), raisonne toujours en trois temps :
\begin{enumerate}
\item \textbf{Avant} : prévention et prévision — aménagement du territoire, entretien des ouvrages, surveillance météorologique, information et éducation des populations ;
\item \textbf{Pendant} : protection immédiate — alerte, mise à l'abri, évacuation ordonnée, respect des consignes ;
\item \textbf{Après} : réponse et résilience — secours aux victimes, soins, reconstruction plus sûre, retour d'expérience pour améliorer la prévention future.
\end{enumerate}
\end{methode}

\begin{popfigure}
\begin{center}
\begin{tikzpicture}[
    node distance=1.2cm,
    box/.style={rectangle, rounded corners=3pt, draw=popblue, fill=popblueL, thick, minimum width=2.2cm, minimum height=1cm, align=center, font=\small\bfseries},
    boxlast/.style={box, fill=popgreenL, draw=popgreen},
    fleche/.style={-{Stealth[length=3mm]}, very thick, poporange},
    flecheloop/.style={fleche, popgreen, dashed}
]
\node[box] (prev) {Prévision};
\node[box, right=of prev] (alerte) {Alerte};
\node[box, right=of alerte] (evac) {Évacuation};
\node[box, right=of evac] (secours) {Secours};
\node[boxlast, right=of secours] (recon) {Reconstruction};

\draw[fleche] (prev) -- (alerte);
\draw[fleche] (alerte) -- (evac);
\draw[fleche] (evac) -- (secours);
\draw[fleche] (secours) -- (recon);

\draw[flecheloop] (recon.south) to[bend right=45] node[below, font=\small\itshape, text=popdark, align=center]{retour d'expérience\\améliorer la prévention} (prev.south);
\end{tikzpicture}
\end{center}
\legende{La chaîne de gestion du risque : chaque maillon (prévision, alerte, évacuation, secours, reconstruction) doit fonctionner pour limiter les pertes. Le retour d'expérience après chaque catastrophe permet de renforcer les maillons amont.}
\end{popfigure}

\courssub{La prévision : voir venir le danger}

\textbf{Observation.} Un cyclone tropical met plusieurs jours à traverser un océan. Grâce à la surveillance continue de la planète, les météorologues peuvent détecter sa formation, suivre sa trajectoire et prévenir les populations menacées plusieurs jours avant son arrivée. C'est ce délai qui sauve des vies.

La prévision météorologique repose sur quatre outils complémentaires :

\begin{itemize}
\item les \cle{satellites météorologiques} : ils photographient en permanence la couverture nuageuse, mesurent la température de surface des océans et repèrent les zones de convection dès leur formation ;
\item les \cle{radars météorologiques} : ils détectent les précipitations, mesurent leur intensité et suivent leur déplacement à l'échelle régionale ;
\item les \cle{bouées océaniques} et stations au sol : elles mesurent en continu la température, la pression atmosphérique, la vitesse du vent, la hauteur des vagues ;
\item les \cle{modèles numériques} : de puissants ordinateurs simulent l'évolution de l'atmosphère à partir des données collectées et calculent les trajectoires probables des cyclones et des zones de fortes pluies.
\end{itemize}

Ces prévisions alimentent les \cle{systèmes d'alerte précoce} : bulletins météorologiques diffusés par la radio, la télévision, les téléphones mobiles, avec des niveaux de vigilance (vert, jaune, orange, rouge) qui déclenchent des consignes précises. Plus l'alerte est donnée tôt, plus l'évacuation peut être ordonnée et complète.

\begin{savaistu}[Le Bangladesh : la preuve que la prévention sauve des vies]
En 1970, le cyclone Bhola tue environ \num{300000} personnes au Bangladesh. En 1991, un autre cyclone en tue encore près de \num{140000}. Le pays investit alors massivement dans un système d'alerte précoce, des milliers d'abris cycloniques en béton surélevés et l'éducation des populations. Résultat : lors du cyclone Amphan en 2020, malgré une intensité comparable aux catastrophes passées, le bilan est d'environ une centaine de morts. La mortalité liée aux cyclones a été réduite de plus de 90~\% en cinquante ans. La leçon est universelle : \textbf{ce ne sont pas les cyclones qui tuent le plus, c'est l'absence de préparation}.
\end{savaistu}

\courssub{La prévention : réduire la vulnérabilité durablement}

La prévention agit \emph{avant} la catastrophe, en aménageant le territoire et en protégeant les milieux naturels qui nous protègent eux-mêmes.

\textbf{1. L'aménagement du territoire.} La règle d'or est simple : \textbf{ne pas construire en zone inondable ni sur les versants instables}. Les documents d'urbanisme doivent identifier les zones à risque (bas-fonds, lits majeurs des cours d'eau, pentes raides) et y interdire l'habitat. À Yaoundé comme à Douala, l'occupation anarchique des bas-fonds et des pentes est la première cause des drames liés aux pluies.

\textbf{2. Le reboisement des versants.} Les racines des arbres fixent le sol et le feuillage ralentit l'écoulement des eaux de pluie. Reboiser les versants réduit à la fois les glissements de terrain et la violence des crues en aval.

\textbf{3. L'entretien des ouvrages de drainage.} Caniveaux, canaux et égouts doivent être curés régulièrement, surtout avant la saison des pluies. Un caniveau obstrué par les déchets plastiques transforme une simple averse en inondation de quartier.

\textbf{4. Les défenses côtières naturelles et artificielles.} Les digues protègent les zones urbanisées, mais la nature offre des protections remarquables qu'il faut préserver.

\begin{propriete}[Rôle protecteur des mangroves et des récifs (admise)]
Les \cle{mangroves} (forêts de palétuviers des côtes tropicales, comme celles de l'estuaire du Wouri et du Rio del Rey au Cameroun) et les \cle{récifs coralliens} réduisent fortement l'énergie des vagues et des ondes de tempête avant qu'elles n'atteignent les côtes habitées. Leur destruction (coupe pour le bois, urbanisation, pollution) aggrave considérablement les dégâts côtiers lors des tempêtes. Protéger et restaurer les mangroves est donc une mesure de prévention à part entière.
\end{propriete}

\courssub{La protection et la réponse : agir au moment du danger}

Lorsque l'alerte est déclenchée, la priorité absolue est la sauvegarde des vies humaines.

\begin{itemize}
\item les \cle{plans d'évacuation} : chaque zone à risque doit disposer d'itinéraires d'évacuation balisés et de points de rassemblement connus de tous, situés en hauteur ;
\item les \cle{abris} : bâtiments publics solides (écoles, salles des fêtes) ou abris spécialement construits accueillent les populations évacuées ;
\item les \cle{secours d'urgence} : protection civile, sapeurs-pompiers, Croix-Rouge et services de santé organisent la recherche des victimes, les premiers soins et l'hébergement provisoire ;
\item l'\cle{éducation aux réflexes de survie} : une population informée réagit vite et correctement, sans panique.
\end{itemize}

\begin{attention}[Le piège mortel des eaux de crue]
Erreur fréquente et mortelle : vouloir traverser à pied ou en voiture une rue inondée. \textbf{30 cm d'eau en mouvement suffisent à emporter un adulte ; 60 cm suffisent à soulever et emporter une voiture.} L'eau de crue cache en outre les plaques d'égout soulevées, les fils électriques tombés et les objets tranchants. Règle absolue : on ne traverse \emph{jamais} une eau de crue, ni à pied, ni en véhicule ; on fait demi-tour et on cherche un itinéraire en hauteur.
\end{attention}

\begin{popfigure}
\begin{center}
\begin{tikzpicture}[scale=0.95, every node/.style={font=\small}, node distance=0.5cm]
% Zone inondable (bas-fond) - Left side
\coordinate (FL) at (0,0);
\fill[popblue!35] (FL) rectangle (4,3.5);
\draw[popblue, thick, dashed] (FL) rectangle (4,3.5);
\node[text=popblue, font=\small\bfseries, anchor=north] at (2,3.3) {Bas-fond inondable (aléa fort)};
% Cours d'eau
\draw[popblue, very thick, decorate, decoration={snake, amplitude=0.4mm, segment length=3mm}] (0,1.5) -- (4,1.5);
\node[text=popblue, font=\scriptsize] at (2,1.1) {Cours d'eau};
% Habitations interdites (hachurées)
\foreach \x/\y in {0.7/2.2, 1.8/2.4, 3.0/2.2}{
    \draw[fill=popmuted, pattern=north east lines, pattern color=popdark] (\x,\y) rectangle ++(0.5,0.4);
}
\node[text=popdark, font=\scriptsize, align=center] at (2,0.3) {Habitat spontané\\à interdire / reloger};

% Zone haute sûre - Right side
\coordinate (SH) at (5.5,0);
\fill[popgreen!15] (SH) rectangle ++(5.5,4.5);
\draw[popgreen, thick] (SH) rectangle ++(5.5,4.5);
\node[text=popgreen, font=\small\bfseries, anchor=north] at (8.25,4.3) {Zone haute sûre (hors aléa)};
% Habitations sûres
\foreach \x/\y in {6.0/0.5, 7.2/0.5, 9.0/0.5, 6.6/1.6, 8.6/1.6, 7.0/2.7, 9.0/2.7}{
    \draw[fill=popgoldL, draw=popgold] (\x,\y) rectangle ++(0.6,0.5);
}
% Point de rassemblement
\node[star, star points=5, fill=poporange, draw=popdark, minimum size=0.7cm, label={[font=\small\bfseries, text=popdark]below:Point de rassemblement}] (pr) at (8.25,3.6) {};
% Itinéraires d'évacuation (flèches courbes)
\draw[-{Stealth[length=3mm]}, line width=2pt, poporange] (2,2.5) .. controls (3.5,3.5) and (5.5,3.8) .. (7.0,3.7);
\draw[-{Stealth[length=3mm]}, line width=2pt, poporange] (3.2,1.8) .. controls (4.5,2.5) and (5.5,3.0) .. (7.0,3.4);
\node[text=poporange, font=\scriptsize\itshape, rotate=15] at (4.8,3.5) {Itinéraires d'évacuation balisés};
\end{tikzpicture}
\end{center}
\legende{Zonage simplifié d'un quartier exposé aux inondations : le bas-fond (bleu) est une zone inondable où l'habitat doit être interdit ; la zone haute (vert) accueille les constructions conformes et le point de rassemblement (étoile orange) ; les flèches orange indiquent les itinéraires d'évacuation à emprunter dès l'alerte.}
\end{popfigure}

\courssub{Conduites à tenir au Cameroun}

\textbf{Pendant de fortes pluies (saison des pluies) :}
\begin{itemize}
\item s'éloigner immédiatement des bas-fonds, des berges des cours d'eau et des pentes instables ;
\item ne jamais traverser les eaux de crue, à pied comme en voiture (rappel : 30 cm emportent un adulte, 60 cm une voiture) ;
\item se réfugier dans un bâtiment solide, de préférence en étage, et couper l'électricité si l'eau menace de monter ;
\item rester informé via la radio ou le téléphone et suivre scrupuleusement les consignes des autorités (protection civile, administration territoriale).
\end{itemize}

\textbf{Pendant le harmattan} (vent sec et poussiéreux soufflant du Sahara vers le Cameroun, généralement entre novembre et février) :
\begin{itemize}
\item boire régulièrement de l'eau pour lutter contre la déshydratation ;
\item protéger les voies respiratoires (foulard ou masque contre la poussière), surtout pour les enfants, les personnes âgées et les asthmatiques ;
\item humidifier les yeux (sérum physiologique) et la peau, limiter les efforts physiques intenses à l'extérieur.
\end{itemize}

\textbf{En toute saison :} suivre les bulletins météorologiques de la Direction de la Météorologie Nationale, particulièrement en début de saison des pluies, et prendre au sérieux les messages de vigilance (jaune, orange, rouge).

\courssub{Les acteurs de la gestion des risques au Cameroun}

\begin{center}
\begin{tabularx}{\linewidth}{|l|X|}
\hline
\rowcolor{popblueL} \textbf{Acteur} & \textbf{Rôle} \\
\hline
Direction de la Météorologie Nationale & Observation, prévision du temps, émission des bulletins et avis de vigilance (niveaux vert/jaune/orange/rouge) \\
\hline
Ministère de l'Administration Territoriale (Protection Civile) & Coordination des secours, organisation des évacuations, gestion des crises, élaboration des plans ORSEC \\
\hline
Croix-Rouge camerounaise & Premiers secours, aide aux sinistrés (hébergement, vivres, non-vivres), sensibilisation des communautés aux gestes qui sauvent \\
\hline
Radios communautaires & Diffusion rapide des alertes et des consignes dans les langues locales, jusqu'aux villages les plus reculés \\
\hline
\end{tabularx}
\end{center}

\courssub{Sensibilisation et éducation aux risques : le rôle de l'école}

L'école est un acteur central de la culture du risque. Un élève informé devient un relais d'information pour toute sa famille. Les établissements scolaires doivent :
\begin{itemize}
\item organiser régulièrement des \cle{exercices d'évacuation} (similaires aux exercices incendie), pour que chacun connaisse les consignes, les itinéraires et le point de rassemblement ;
\item afficher les \cle{consignes de sécurité} (numéros d'urgence, conduite à tenir) dans les salles de classe et les lieux de vie ;
\item intégrer l'éducation aux risques dans l'enseignement, comme le fait cette leçon, pour former des citoyens responsables et résilients.
\end{itemize}

\begin{exoresolu}[Proposer des mesures de protection pour un quartier inondable]
\textbf{Énoncé.} Un quartier de Douala, installé dans un bas-fond, est inondé chaque saison des pluies. En 2022, une crue a causé la mort de trois personnes emportées par les eaux et détruit de nombreuses habitations. Propose un plan d'action complet pour réduire ce risque, en distinguant les actions à mener avant, pendant et après les pluies.
\tcblower
\textbf{Solution détaillée.} On applique la méthode en trois temps (Avant / Pendant / Après).

\emph{Avant (Prévention et Prévision) :}
\begin{itemize}
\item Cartographier précisément les zones les plus inondables (lits majeurs, bas-fonds) et interdire toute nouvelle construction par un arrêté municipal ;
\item Curer les caniveaux, canaux et regards d'égout avant l'arrivée des pluies (campagne « ville propre ») et sanctionner le dépôt de déchets plastiques dans les drains ;
\item Reboiser les versants environnants (ex. pentes de Bonamoussadi, Ndogpassi) pour limiter le ruissellement et l'érosion ;
\item Engager un programme progressif de relogement des habitants des zones à risque très élevé vers des sites viabilisés hors aléa ;
\item Informer la population via les radios communautaires, les chefs de quartier et les SMS d'alerte sur les zones interdites et les itinéraires d'évacuation ;
\item Suivre quotidiennement les bulletins de la Direction de la Météorologie Nationale et activer la veille cyclonique/pluies intenses.
\end{itemize}

\emph{Pendant (Protection) :}
\begin{itemize}
\item Déclencher l'alerte (sirènes, SMS, radio) dès l'avis de vigilance \emph{orange} ou \emph{rouge} « Pluies intenses / Inondations » ;
\item Évacuer préventivement les habitants vers le point de rassemblement identifié (ex. école du plateau voisin, salle des fêtes) par les itinéraires balisés, en priorisant les personnes vulnérables (enfants, personnes âgées, handicapés) ;
\item Interdire strictement et physiquement (barrières, présence de forces de l'ordre) la traversée des rues inondées ; rappeler la règle : \textbf{30 cm d'eau en mouvement emportent une personne, 60 cm une voiture} ;
\item Couper l'électricité dans le quartier inondé pour éviter les électrocutions.
\end{itemize}

\emph{Après (Réponse et Résilience) :}
\begin{itemize}
\item Organiser les secours (Protection civile, Croix-Rouge, SAMU) : recherche de victimes, premiers soins, évacuation sanitaire ;
\item Héberger, nourrir et assurer un soutien psychosocial aux sinistrés dans les centres d'hébergement d'urgence ;
\item Déblayer les décombres, désinfecter les lieux (pulvérisation larvicide) pour éviter les épidémies post-inondation (choléra, fièvre typhoïde, paludisme, leptospirose) ;
\item Reconstruire \emph{hors} des zones inondables identifiées, avec des matériaux durables et sur pilotis si nécessaire, en appliquant les normes parasismiques et paracycloniques ;
\item Dresser un bilan (retour d'expérience) : efficacité de l'alerte, fluidité de l'évacuation, points de blocage, pour mettre à jour le plan communal de sauvegarde.
\end{itemize}
\end{exoresolu}

\courssub{Un enjeu de développement durable}

Avec le changement climatique, les phénomènes extrêmes (pluies diluviennes, sécheresses prolongées, tempêtes violentes, élévation du niveau de la mer) deviennent plus fréquents et plus intenses. Adapter l'urbanisation (construire hors des zones à risque, préserver les zones humides qui absorbent les crues, végétaliser les villes) et l'agriculture (choix de variétés résistantes, calendriers culturaux décalés, stockage de l'eau de pluie, agroforesterie) aux risques climatiques croissants est une condition du \cle{développement durable} : un développement qui répond aux besoins du présent sans compromettre la capacité des générations futures à faire face aux aléas. Protéger les mangroves, reboiser les collines, éduquer les populations aux gestes qui sauvent : autant d'actions qui relient directement cette leçon à l'ensemble du module d'Éducation à l'Environnement et au Développement Durable.

\begin{aretenir}[L'essentiel de la section]
\begin{itemize}
\item Risque = Aléa $\times$ Vulnérabilité $\times$ Enjeu : on ne peut pas supprimer l'aléa naturel, mais on peut et on doit réduire la vulnérabilité (fragilité des biens et des personnes) et l'exposition des enjeux (localisation).
\item La gestion des risques repose sur trois piliers indissociables : \textbf{prévision} (satellites, radars, bouées, modèles numériques $\rightarrow$ alerte précoce), \textbf{prévention} (aménagement du territoire, reboisement, drainage, protection des mangroves/récifs), \textbf{protection} (évacuation, abris, secours, éducation aux réflexes).
\item La démarche de résolution d'un problème de protection suit toujours trois temps : \textbf{Avant} (prévenir/prévoir), \textbf{Pendant} (protéger/évacuer), \textbf{Après} (secourir/reconstruire/tirer les leçons).
\item \textbf{Règle de survie absolue :} on ne traverse \emph{jamais} une eau de crue (30 cm = adulte emporté ; 60 cm = voiture emportée).
\item Au Cameroun, les acteurs clés sont la Direction de la Météorologie Nationale, la Protection Civile (MINAT), la Croix-Rouge et les radios communautaires ; l'école forme les citoyens aux réflexes de survie.
\item La prévention sauve des vies : l'exemple du Bangladesh (réduction de 90~\% de la mortalité cyclonique en 50 ans) prouve que l'investissement dans l'alerte précoce et les abris est rentable en vies humaines.
\end{itemize}
\end{aretenir}

\newpage

\coursec{Activité d'intégration}

\courssub{Objectifs de l'activité}

À l'issue de cette activité d'intégration, tu seras capable de :

\begin{itemize}
\item \cle{Décrire} un mouvement atmosphérique (mousson, ZCIT) à partir d'un schéma ou d'une carte ;
\item \cle{Relier} un phénomène météorologique dangereux (pluies diluviennes, inondations) à son origine dynamique (circulation atmosphérique et océanique) ;
\item \cle{Analyser} des données chiffrées (hauteurs de pluie, nombre de sinistrés) pour identifier les facteurs aggravants d'une catastrophe ;
\item \cle{Proposer} des mesures de prévention et de protection réalistes, chiffrées et argumentées ;
\item Développer des \cle{savoir-être} : esprit critique, travail en équipe, sens de la responsabilité citoyenne face aux risques naturels.
\end{itemize}

\courssub{Situation}

\textbf{Commission d'alerte : protéger le quartier de Ndogpassi face aux inondations.}

La ville de Douala, capitale économique du Cameroun, est située au bord de l'estuaire du Wouri, à très faible altitude (entre 0 et \SI{10}{m} au-dessus du niveau de la mer). Chaque année, pendant la grande saison des pluies, plusieurs quartiers connaissent des inondations dramatiques. Le quartier de Ndogpassi (Douala III) est l'un des plus touchés.

Le maire de l'arrondissement décide de créer une \cle{commission municipale de prévention des risques}. Ta classe est invitée à jouer le rôle de cette commission. Chaque groupe d'élèves reçoit le dossier suivant.

\textbf{Document 1 : carte simplifiée du quartier de Ndogpassi.} Le quartier s'étend sur \SI{2}{km^2}. On y distingue : une zone haute (cote $+8$ m) au nord, densément construite ; une zone basse (cote $+2$ m) au centre, où vivent environ \num{12000} personnes ; un marigot (ancien bras du Wouri) au sud, partiellement comblé par des constructions informelles ; un réseau de caniveaux dont 70 \% sont obstrués par des déchets plastiques et des sédiments.

\textbf{Document 2 : schéma de la circulation atmosphérique en saison des pluies.} En juillet-août, la \cle{Zone de Convergence InterTropicale} (ZCIT) se déplace vers le nord et se positionne au-dessus du Cameroun septentrional. L'\cle{alizé de sud-ouest}, chargé d'humidité après avoir traversé l'océan Atlantique équatorial (température de surface de l'eau : $27$ à \SI{29}{\celsius} dans le golfe de Guinée), devient la \cle{mousson humide} qui pénètre profondément dans les terres. En rencontrant l'air continental plus sec (alizé de nord-est, l'harmattan), l'air humide est forcé de s'élever, se refroidit et provoque des précipitations abondantes et durables.

\begin{popfigure}
\begin{center}
\begin{tikzpicture}[scale=0.95, every node/.style={font=\small}]
% Océan
\fill[popblue!25] (-6,-1.2) rectangle (0,0);
\node[popdark] at (-3,-0.9) {Océan Atlantique équatorial};
\node[popdark] at (-3,-1.55) {\footnotesize eau de surface : $27$--\SI{29}{\celsius}};
% Continent
\fill[popgold!35] (0,-1.2) -- (0,0) -- (6,0.8) -- (6,-1.2) -- cycle;
\node[popdark] at (3.2,-0.85) {Continent africain (surchauffé)};
% Évaporation
\draw[-{Stealth}, popblue, thick] (-4.5,0) .. controls (-4.5,0.8) .. (-4.2,1.4);
\draw[-{Stealth}, popblue, thick] (-2.5,0) .. controls (-2.5,0.9) .. (-2.2,1.5);
\node[popblue, align=center] at (-3.4,2.0) {\footnotesize évaporation intense};
% Mousson
\draw[-{Stealth}, popgreen, very thick] (-4.8,0.6) .. controls (-2,1.0) and (0,1.2) .. (2.2,1.6);
\node[popgreen!60!black] at (-1.2,1.75) {\footnotesize mousson humide (alizé de SO)};
% Harmattan
\draw[-{Stealth}, poporange, very thick] (5.6,2.6) -- (3.4,2.0);
\node[poporange!80!black] at (5.0,3.0) {\footnotesize harmattan (alizé de NE, sec)};
% Ascendance ZCIT
\draw[-{Stealth}, poppink, very thick] (2.6,1.7) .. controls (2.7,2.6) .. (2.6,3.6);
\node[poppink!70!black, align=center] at (2.6,4.0) {\footnotesize ZCIT : ascendance,};
\node[poppink!70!black, align=center] at (2.6,3.55) {};
\node[poppink!70!black] at (4.6,3.9) {\footnotesize refroidissement, condensation};
% Nuage et pluie
\draw[fill=popmuted!40, draw=popmuted] (1.6,4.4) ellipse (1.1 and 0.45);
\node[popdark] at (1.6,4.4) {\footnotesize cumulonimbus};
\foreach \x in {1.0,1.4,1.8,2.2} \draw[popblue, thick] (\x,3.95) -- (\x-0.15,3.45);
\node[popblue] at (0.4,3.5) {\footnotesize pluies};
\end{tikzpicture}
\end{center}
\legende{Mécanisme de la mousson humide en été boréal : l'air chargé de vapeur d'eau au-dessus de l'Atlantique équatorial chaud est aspiré vers le continent surchauffé ; à la ZCIT, il est soulevé par l'air sec de l'harmattan, se refroidit, et la vapeur d'eau se condense en pluies abondantes.}
\end{popfigure}

\textbf{Document 3 : données pluviométriques mensuelles moyennes à Douala (station de l'aéroport).}

\begin{center}
\begin{tabularx}{\linewidth}{|l|X|X|X|X|X|X|}
\hline
\rowcolor{popblueL}
Mois & Avril & Mai & Juin & Juillet & Août & Sept. \\
\hline
Pluie (mm) & 220 & 280 & 350 & 690 & 700 & 560 \\
\hline
Jours de pluie & 17 & 20 & 22 & 26 & 27 & 25 \\
\hline
\end{tabularx}
\end{center}

Douala reçoit en moyenne \SI{4000}{mm} de pluie par an, ce qui en fait l'une des villes les plus arrosées d'Afrique.

\textbf{Document 4 : bilan des inondations d'août de l'année écoulée à Ndogpassi.} Hauteur d'eau maximale dans la zone basse : \SI{1,20}{m} ; \num{3500} sinistrés ; 12 cases effondrées ; 3 décès par noyade ; épidémie de choléra ayant touché 87 personnes dans les 3 semaines suivantes (eaux stagnantes contaminées) ; fermeture de 4 écoles pendant 15 jours. Coût estimé des dégâts : 800 millions de FCFA. Budget annuel de prévention de la mairie : 150 millions de FCFA.

\textbf{Problème posé à la commission :} comment expliquer l'origine de ces pluies diluviennes et de ces inondations, et quel plan de protection chiffré proposer au conseil municipal pour limiter les dégâts lors de la prochaine saison des pluies ?

\courssub{Consignes}

\begin{enumerate}
\item À partir du document 2, décris le trajet de la mousson humide depuis l'océan Atlantique jusqu'à Douala, en précisant le rôle de la ZCIT et le mécanisme de formation des pluies (évaporation, advection, ascendance, condensation, précipitations).
\item Explique pourquoi les pluies sont maximales en juillet-août à Douala, en reliant la position de la ZCIT à la circulation atmosphérique.
\item Calcule la hauteur totale de pluie tombée sur le quartier pendant les trois mois de juillet, août et septembre, puis le volume d'eau correspondant tombé sur les \SI{2}{km^2} du quartier (rappel : $V = h \times S$).
\item Relie les inondations de Ndogpassi à leur origine dynamique (mousson + ZCIT) et identifie, à partir des documents 1 et 4, au moins trois facteurs aggravants locaux.
\item Rédige un plan de protection en quatre volets (zonage, drainage, alerte, évacuation), chiffré à l'aide du budget disponible (150 millions de FCFA), et argumenté scientifiquement.
\item Présente ton plan en 5 minutes devant le « conseil municipal » (la classe), qui votera le meilleur plan selon la grille : exactitude scientifique (4 pts), réalisme du chiffrage (4 pts), qualité de l'argumentation (4 pts), clarté de la présentation (4 pts), réponses aux questions (4 pts).
\end{enumerate}

\courssub{Ressources à mobiliser}

\begin{aretenir}[Ce qu'il faut avoir en tête]
\begin{itemize}
\item La \cle{ZCIT} est la zone de convergence des alizés des deux hémisphères ; l'air chaud et humide y subit une ascendance, se refroidit, et la vapeur d'eau se condense : ce sont les \cle{précipitations de convection}.
\item La \cle{mousson} est un renversement saisonnier des vents : en été boréal, le continent africain surchauffé (basses pressions thermiques) aspire l'air humide océanique (alizé de sud-ouest devenu mousson).
\item L'océan chaud ($\geq$ \SI{26}{\celsius}) est le réservoir d'eau et d'énergie : évaporation intense $\Rightarrow$ air chargé en humidité.
\item Une inondation résulte toujours de la combinaison d'un \cle{aléa} (pluies exceptionnelles, origine dynamique) et d'une \cle{vulnérabilité} (topographie basse, urbanisation anarchique, drainage défaillant).
\item La prévention repose sur quatre piliers : connaissance du risque (zonage), réduction de la vulnérabilité (aménagement), surveillance et alerte, préparation à la crise (évacuation).
\item Calculs utiles : $V = h \times S$ ; \SI{1}{mm} de pluie sur \SI{1}{km^2} $= \num{1000}$ m$^3$ d'eau.
\end{itemize}
\end{aretenir}

\courssub{Démarche et corrigé commenté}

\begin{exoresolu}[Corrigé commenté de la commission d'alerte]
\textbf{Étape 1 : description du mécanisme atmosphérique (consigne 1).}

L'océan Atlantique équatorial, chauffé à $27$--\SI{29}{\celsius}, est le siège d'une évaporation intense. L'air de surface s'y charge en vapeur d'eau. Sous l'effet de la différence de pression entre l'océan (relativement frais, pressions plus élevées) et le continent africain surchauffé (basses pressions thermiques), cet air humide s'écoule vers le nord-est : c'est la \cle{mousson de sud-ouest}. Arrivée sur le continent, cette masse d'air chaude et humide rencontre, au niveau de la ZCIT, l'air sec de l'harmattan. Moins dense, l'air humide est soulevé ; en s'élevant, il se refroidit (environ \SI{0,6}{\celsius} par 100 m) ; la vapeur d'eau se condense en gouttelettes formant des nuages convectifs (cumulonimbus) qui déversent des pluies abondantes et prolongées sur le littoral camerounais.

\tcblower

\textbf{Étape 2 : pourquoi un maximum en juillet-août ? (consigne 2).}

Le maximum de pluie correspond au moment où la ZCIT atteint sa position la plus septentrionale : la mousson humide pénètre alors le plus profondément dans les terres et Douala reste durablement sous l'influence du flux océanique humide. C'est aussi la période où la mer est la plus chaude, donc où l'évaporation et la recharge en humidité sont maximales. D'où les valeurs records de juillet (\SI{690}{mm}) et août (\SI{700}{mm}).

\textbf{Étape 3 : calculs hydrologiques (consigne 3).}

Hauteur totale juillet + août + septembre :
\[ h = 690 + 700 + 560 = \SI{1950}{mm} = \SI{1,95}{m}. \]
Surface du quartier : $S = \SI{2}{km^2} = \num{2000000}$ m$^2$.
Volume d'eau tombé :
\[ V = h \times S = 1,95 \times \num{2000000} = \num{3900000} \text{ m}^3. \]
Soit près de \textbf{4 millions de mètres cubes} d'eau — l'équivalent de plus de \num{1500} piscines olympiques — tombés en trois mois sur un quartier plat de \SI{2}{km^2}. Ce calcul justifie à lui seul l'ampleur du risque : même une fraction de ce volume suffit à submerger la zone basse si le drainage est insuffisant.

\textbf{Étape 4 : relation cause-effet et facteurs aggravants (consigne 4).}

\emph{Origine dynamique :} les pluies diluviennes résultent de la conjonction mousson humide + ZCIT (aléa climatique naturel, prévisible et saisonnier).

\emph{Facteurs aggravants locaux (vulnérabilité) :}
\begin{enumerate}
\item \textbf{Topographie :} la zone basse (cote $+2$ m) ne peut pas s'écouler gravitairement lors de fortes pluies, d'autant que les marées hautes dans l'estuaire bloquent l'exutoire vers le Wouri.
\item \textbf{Drainage défaillant :} 70 \% des caniveaux sont obstrués (déchets plastiques, sédiments) ; la capacité d'évacuation est réduite à moins d'un tiers de sa valeur nominale.
\item \textbf{Urbanisation anarchique :} constructions informelles sur le marigot comblé, qui servait autrefois de zone d'expansion naturelle des crues ; imperméabilisation des sols (béton, latérite compactée) qui supprime l'infiltration et accélère le ruissellement.
\item \textbf{Facteur sanitaire :} les eaux stagnantes mélangées aux déchets et aux latrines débordées ont déclenché l'épidémie de choléra (87 cas).
\end{enumerate}

\textbf{Étape 5 : plan de protection chiffré (consigne 5).}

Budget disponible : 150 millions de FCFA. Proposition de répartition :

\begin{center}
\begin{tabularx}{\linewidth}{|l|X|X|}
\hline
\rowcolor{popblueL}
Volet & Actions & Coût (millions FCFA) \\
\hline
1. Zonage & Cartographie des zones inondables ; interdiction de construire dans le marigot ; relogement progressif des 200 familles les plus exposées & 40 \\
\hline
2. Drainage & Curage des caniveaux avant juin ; élargissement du collecteur principal ; campagne de gestion des déchets plastiques & 60 \\
\hline
3. Alerte & Station pluviométrique scolaire + relais radio communautaire ; seuil d'alerte fixé à \SI{80}{mm} en 24 h ; météo locale diffusée par SMS & 15 \\
\hline
4. Évacuation & Aménagement de 2 sites-refuges en zone haute (écoles) ; exercice d'évacuation annuel en mai ; kits d'urgence (eau potable, pastilles de purification) & 35 \\
\hline
\textbf{Total} & & \textbf{150} \\
\hline
\end{tabularx}
\end{center}

\emph{Argumentation :} le volet drainage est prioritaire (coût le plus élevé) car il attaque directement le facteur aggravant n°2, le plus facilement corrigeable ; le zonage traite la cause structurelle mais ses effets sont à long terme ; l'alerte et l'évacuation ne suppriment pas le risque mais réduisent drastiquement le nombre de victimes (les 3 décès de l'an passé étaient évitables avec un seuil d'alerte et des sites-refuges).

\textbf{Étape 6 : présentation et vote (consigne 6).}

Chaque groupe présente son plan en 5 minutes. Le conseil municipal (la classe) évalue selon la grille sur 20 points. Le débat permet de comparer les arbitrages : faut-il privilégier le curage immédiat ou le relogement de long terme ? La meilleure réponse combine les deux, car la prévention efficace agit à la fois sur l'aléa (surveillance) et sur la vulnérabilité (aménagement).
\end{exoresolu}

\begin{attention}[Pièges fréquents dans ce type d'activité]
\begin{itemize}
\item Confondre \cle{aléa} (le phénomène naturel, qu'on ne peut pas supprimer) et \cle{vulnérabilité} (les facteurs humains, sur lesquels on peut agir) : les inondations de Douala ne sont pas une « fatalité ».
\item Oublier les unités dans le calcul de volume : convertir les mm en m et les km$^2$ en m$^2$ avant d'appliquer $V = h \times S$.
\item Affirmer que la ZCIT « produit » la pluie : c'est l'\emph{ascendance de l'air humide} au niveau de la convergence qui provoque condensation et précipitations.
\item Proposer des mesures sans les chiffrer ni les hiérarchiser : un plan de prévention crédible respecte une contrainte budgétaire.
\end{itemize}
\end{attention}

\courssub{Bilan de l'activité}

Cette activité t'a permis de mettre en évidence plusieurs idées essentielles :

\begin{itemize}
\item Les catastrophes liées aux mouvements atmosphériques et océaniques ont une \cle{origine dynamique identifiable} : la mousson humide née sur l'Atlantique équatorial chaud et la position de la ZCIT expliquent le rythme saisonnier des pluies diluviennes de Douala.
\item Une catastrophe n'est jamais purement « naturelle » : elle résulte de la rencontre d'un \cle{aléa} et d'une \cle{vulnérabilité}. À Douala, la faible altitude, le comblement du marigot, l'obstruction des caniveaux et l'urbanisation non planifiée transforment une pluie saisonnière prévisible en drame humain et sanitaire (noyades, choléra).
\item Les calculs simples (volume d'eau tombé) donnent une \cle{mesure objective du risque} et justifient l'ampleur des investissements de prévention.
\item La protection des populations repose sur une démarche organisée : \cle{zonage} (connaître), \cle{drainage} (aménager), \cle{alerte} (surveiller), \cle{évacuation} (se préparer). La prévention coûte bien moins cher que la réparation : 150 millions de FCFA de prévention contre 800 millions de dégâts constatés.
\item Enfin, tu as exercé des compétences transversales : analyser des documents, calculer, argumenter, débattre et voter — autant de savoir-faire du citoyen éclairé, acteur du développement durable de sa ville.
\end{itemize}

\newpage

\coursec{Méthodes et exercices résolus}

\coursec{Leçon ND12 : Sensibilisation dans le cadre des catastrophes liées aux mouvements atmosphériques et océaniques}

\begin{aretenir}[Cadre officiel -- Module 3 : Éducation à l'Environnement et au Développement Durable (1ère D -- 8 h)]
    \textbf{Famille de situations :} Dégâts causés par les mouvements atmosphériques et océaniques.
    \par\medskip
    \textbf{Savoirs (Programme officiel) :}
    \begin{itemize}
        \item Circulation générale de l'atmosphère (vents, cellules de convection).
        \item Circulation océanique (courants marins) et interactions avec l'atmosphère.
        \item Phénomènes climatiques dangereux liés à ces mouvements (tempêtes, cyclones, inondations).
    \end{itemize}
    \textbf{Savoir-faire (Programme officiel) :}
    \begin{itemize}
        \item Décrire un mouvement atmosphérique ou océanique à partir d'un schéma ou d'une carte.
        \item Relier un phénomène météorologique dangereux à son origine dynamique.
        \item Proposer des mesures de protection face à ces phénomènes.
    \end{itemize}
    \textbf{Savoir-être :} Conscience des risques naturels, esprit de prévention, responsabilité citoyenne face aux alertes, solidarité en situation de crise.
\end{aretenir}

\begin{aretenir}[Plan général de la leçon]
    \begin{enumerate}
        \item La circulation générale de l'atmosphère (Cellules de Hadley, Ferrel, Polaire ; ZCIT ; Alizés, Vents d'ouest).
        \item La circulation océanique : les courants marins (Gyres subtropicaux, Courants de bord ouest/est, Circulation thermohaline).
        \item Interactions océan-atmosphère (Couplage thermique, El Niño/La Niña, Mousson, Cyclogenèse).
        \item Les phénomènes climatiques dangereux (Cyclones tropicaux, Lignes de grain, Inondations, Sécheresses, Submersions marines).
        \item Prévention et protection face aux catastrophes (Gestion des risques, SAP, PCS, Adaptation au changement climatique, NbS).
    \end{enumerate}
\end{aretenir}

\courssub{Méthode 1 : Décrire la circulation générale de l'atmosphère à partir d'un schéma ou d'une carte}
\begin{methode}[Démarche pour analyser un schéma de la circulation atmosphérique générale]
    \textbf{Objectif :} Identifier les cellules de convection, les vents dominants, les zones de convergence/divergence et les ceintures de précipitations.
    \begin{enumerate}
        \item \textbf{Observer le cadre :} Repérer l'hémisphère (Nord/Sud), l'axe de l'équateur, les pôles, les continents et océans représentés.
        \item \textbf{Identifier les trois cellules par hémisphère :}
              \begin{itemize}
                  \item \textbf{Cellule de Hadley (0°--30°) :} Air chaud montant à l'équateur (convergence alizés, \cle{Zone de Convergence Intertropicale -- ZCIT}), divergence en altitude, subsidence aux latitudes des chevaux (30°), vents d'alizés au sol (NE hémisphère Nord, SE hémisphère Sud).
                  \item \textbf{Cellule de Ferrel (30°--60°) :} Circulation \textbf{thermiquement indirecte}. Montée vers 60° (front polaire), subsidence vers 30°. Vents d'ouest dominants au sol (\cle{vents d'ouest}).
                  \item \textbf{Cellule polaire (60°--90°) :} Air froid descendant au pôle, se dirigeant vers l'équateur au sol (\cle{vents d'est polaires}), montant vers 60°.
              \end{itemize}
        \item \textbf{Noter l'effet de Coriolis :} Déviation vers la \textbf{droite} dans l'hémisphère Nord, vers la \textbf{gauche} dans l'hémisphère Sud. Explique le sens de rotation des alizés et des vents d'ouest.
        \item \textbf{Localiser les zones climatiques clés :}
              \begin{itemize}
                  \item \textbf{ZCIT :} Basse pression, forts orages, pluies abondantes (forêt équatoriale).
                  \item \textbf{Subtropicaux (30°) :} Haute pression (anticyclones), temps sec, déserts chauds (Sahara, Kalahari, Atacama).
                  \item \textbf{60° :} Basse pression subpolaire, fronts polaires, temps perturbé (dépressions tempérées).
                  \item \textbf{Pôles :} Haute pression thermique, air très froid et sec.
              \end{itemize}
        \item \textbf{Rédiger la description :} Suivre le parcours d'une parcelle d'air : chauffage $\rightarrow$ montée $\rightarrow$ transport altitude $\rightarrow$ refroidissement $\rightarrow$ subsidence $\rightarrow$ retour au sol. Préciser les noms des vents (Alizés, Vents d'ouest, Vents polaires).
    \end{enumerate}
    \textbf{Réflexe Probatoire :} Sur une carte, les flèches rouges = air chaud/montant ; bleues = air froid/descendant. Au sol, les flèches indiquent le sens du vent (d'où il vient).
\end{methode}

\begin{exoresolu}[Analyse d'un schéma simplifié de la circulation atmosphérique (Type Probatoire)]
On considère le schéma ci-dessous représentant une coupe méridienne de l'atmosphère terrestre (hémisphère Nord). Les flèches pleines indiquent les mouvements d'air au sol et en altitude. Les lettres A, B, C, D correspondent à des latitudes caractéristiques.
\begin{popfigure}
    \begin{tikzpicture}[scale=0.9, every node/.style={font=\small}]
        % Terre
        \draw[popdark, thick] (0,0) arc (180:0:6cm);
        \draw[popdark, thick] (-6,0) -- (6,0);
        % Latitudes
        \foreach \lat/\x in {90/0, 60/3, 30/5.2, 0/6, -30/5.2, -60/3, -90/0} {
            \draw[popmuted, dashed] ({-6*cos(\lat)}, {6*sin(\lat)}) -- ({6*cos(\lat)}, {6*sin(\lat)});
        }
        % Equateur
        \draw[popblue, thick] (0,0) -- (0,6) node[above, pos=1.1] {Équateur};
        % Pôle Nord
        \node[popdark] at (0,6.5) {Pôle Nord (90°N)};
        % Cellule Hadley (0-30)
        \draw[poporange, thick, ->] (0.2,0.2) -- (0.2,3.5) node[midway, right] {Montée};
        \draw[poporange, thick, ->] (0.2,3.5) -- (5,3.5) node[midway, above] {Altitude};
        \draw[poporange, thick, ->] (5,3.5) -- (5,0.2) node[midway, right] {Subsidence};
        \draw[poporange, thick, ->] (5,0.2) -- (0.2,0.2) node[midway, below] {Alizés NE};
        % Cellule Ferrel (30-60)
        \draw[popgreen, thick, ->] (5,0.2) -- (5.2,2.5);
        \draw[popgreen, thick, ->] (5.2,2.5) -- (3,2.5) node[midway, above] {Altitude};
        \draw[popgreen, thick, ->] (3,2.5) -- (3,0.2) node[midway, right] {Subsidence};
        \draw[popgreen, thick, ->] (3,0.2) -- (5,0.2) node[midway, below] {Vents d'Ouest};
        % Cellule Polaire (60-90)
        \draw[popblue, thick, ->] (3,0.2) -- (2.8,1.5);
        \draw[popblue, thick, ->] (2.8,1.5) -- (0.5,1.5) node[midway, above] {Altitude};
        \draw[popblue, thick, ->] (0.5,1.5) -- (0.5,0.2) node[midway, right] {Subsidence};
        \draw[popblue, thick, ->] (0.5,0.2) -- (3,0.2) node[midway, below] {Vents d'Est polaires};
        % Labels latitudes
        \node[popdark] at (-6.5, 0) {A: 0°};
        \node[popdark] at (-6.5, 3.1) {B: 30°N};
        \node[popdark] at (-6.5, 5.2) {C: 60°N};
        \node[popdark] at (-6.5, 6) {D: 90°N};
        % Pression au sol
        \node[popred] at (1, -0.5) {BP Équatoriale};
        \node[popblue] at (4.5, -0.5) {HP Subtropicale};
        \node[popred] at (2.5, -0.5) {BP Subpolaire};
        \node[popblue] at (0.2, -0.5) {HP Polaire};
    \end{tikzpicture}
    \legende{Schéma simplifié de la circulation méridienne atmosphérique (Hémisphère Nord).}
\end{popfigure}
\textbf{Questions :}
\begin{enumerate}
    \item Nommez les trois cellules de convection représentées et donnez leurs limites latitudinales approximatives.
    \item Identifiez les vents dominants au niveau du sol pour chaque cellule (noms et sens).
    \item Expliquez pourquoi les zones B et D correspondent à des hautes pressions au sol, alors que A et C correspondent à des basses pressions.
    \item Le Cameroun se situe entre 2°N et 13°N. Dans quelle cellule se trouve-t-il ? Quels vents dominants y soufflent au sol ? Quelle est la conséquence sur le climat (saison sèche/pluvieuse) liée à la migration de la ZCIT ?
\end{enumerate}
\tcblower
\textbf{Solution détaillée :}
\begin{enumerate}
    \item \textbf{Cellule de Hadley :} entre l'équateur (A, 0°) et les subtropicales (B, ~30°N).
          \textbf{Cellule de Ferrel :} entre les subtropicales (B, ~30°N) et les subpolaires (C, ~60°N).
          \textbf{Cellule Polaire :} entre les subpolaires (C, ~60°N) et le pôle (D, 90°N).
    \item \begin{itemize}
              \item \textbf{Hadley (sol) :} \textbf{Alizés du Nord-Est} (vent venant du NE vers l'équateur). Déviation de Coriolis vers la droite d'un vent qui irait du Nord vers le Sud.
              \item \textbf{Ferrel (sol) :} \textbf{Vents d'Ouest} (vent venant de l'Ouest/SO vers le NE). Déviation vers la droite d'un vent allant du Sud vers le Nord.
              \item \textbf{Polaire (sol) :} \textbf{Vents d'Est polaires} (vent venant de l'Est/NE vers l'Ouest/SO). Air froid descendant au pôle s'écoulant vers l'équateur, dévié vers la droite.
          \end{itemize}
    \item \textbf{Physique des pressions :}
          \begin{itemize}
              \item \textbf{A (Équateur) :} Fort ensoleillement $\rightarrow$ fort réchauffement du sol $\rightarrow$ réchauffement air contact $\rightarrow$ dilatation, densité faible $\rightarrow$ \textbf{montée convection} $\rightarrow$ \textbf{Basse Pression (BP)} au sol. Convergence des alizés (ZCIT).
              \item \textbf{B (30°N) :} Air monté à l'équateur se refroidit en altitude, perd son humidité (précipitations équatoriales), descend par subsidence (mouvement descendant adiabatique $\rightarrow$ réchauffement par compression, assèchement) $\rightarrow$ accumulation d'air au sol $\rightarrow$ \textbf{Haute Pression (HP)} (Anticyclones subtropicaux, ex: Anticyclone des Açores). Déserts.
              \item \textbf{C (60°N) :} Rencontre de l'air doux/humide venu des subtropicales (vents d'Ouest) et de l'air froid polaire (vents d'Est polaires) $\rightarrow$ \textbf{Front polaire}. L'air chaud glisse sur l'air froid (front chaud) ou l'air froid soulève l'air chaud (front froid) $\rightarrow$ \textbf{montée forcée} $\rightarrow$ \textbf{Basse Pression (BP)} subpolaire. Dépressions tempérées.
              \item \textbf{D (Pôle) :} Rayonnement solaire nul/faible $\rightarrow$ fort refroidissement sol/air $\rightarrow$ air dense, lourd $\rightarrow$ \textbf{subsidence} $\rightarrow$ accumulation au sol $\rightarrow$ \textbf{Haute Pression (HP)} thermique polaire.
          \end{itemize}
    \item \textbf{Localisation Cameroun :} Latitudes 2°N--13°N $\rightarrow$ pleinement dans la \textbf{Cellule de Hadley}.
          \textbf{Vents dominants au sol :} \textbf{Alizés du Nord-Est} (Harmattan en saison sèche) et \textbf{Mousson du Sud-Ouest} (en saison des pluies).
          \textbf{Mécanisme saisonnier :} La ZCIT migre suivant le soleil zénithal.
          \begin{itemize}
              \item \textbf{Saison sèche (Nov--Mars) :} ZCIT au sud du Cameroun (ou au large). Le pays sous influence des \textbf{Alizés NE (Harmattan)} : air chaud, sec, poussiéreux, ciel dégagé. HP subtropicale influence le nord.
              \item \textbf{Saison des pluies (Avr--Oct) :} ZCIT remonte au nord du pays (vers 13°N en août). Le pays sous influence de la \textbf{Mousson SW} (air maritime tropical humide, instable) : orages, pluies abondantes. La ZCIT \emph{est} la limite entre Alizés NE et Mousson SW.
          \end{itemize}
          \textbf{Conclusion :} Le climat camerounais (soudanien au nord, équatorial au sud) est directement commandé par la dynamique de la cellule de Hadley et la migration latitudinale de la ZCIT.
\end{enumerate}
\end{exoresolu}

\courssub{Méthode 2 : Décrire la circulation océanique de surface (gyres) et profonde (thermohaline) à partir d'une carte}
\begin{methode}[Démarche pour analyser une carte des courants marins]
    \textbf{Objectif :} Identifier les grands gyres, le sens de rotation, la nature des courants (chauds/froids), et le tapis roulant mondial.
    \begin{enumerate}
        \item \textbf{Repérer les bassins océaniques :} Pacifique, Atlantique, Indien (hémisphères Nord et Sud).
        \item \textbf{Identifier les gyres subtropicaux (5 par bassin) :}
              \begin{itemize}
                  \item \textbf{Hémisphère Nord :} Rotation \textbf{horaire} (sens des aiguilles d'une montre). Ex: Gyre Nord-Atlantique, Gyre Nord-Pacifique.
                  \item \textbf{Hémisphère Sud :} Rotation \textbf{anti-horaire} (sens inverse). Ex: Gyre Sud-Atlantique, Gyre Sud-Pacifique, Gyre Indien.
              \end{itemize}
              \textbf{Origine :} Vent dominant (Alizés à l'Est $\rightarrow$ courant équatorial Ouest ; Vents d'Ouest au Nord $\rightarrow$ courant de retour Est) + Effet Coriolis + Fermeture par les continents.
        \item \textbf{Classifier les courants de bord :}
              \begin{itemize}
                  \item \textbf{Courants de bord ouest (Western Boundary Currents) :} Étroits, profonds, \textbf{rapides}, \textbf{chauds} (transportent chaleur vers pôles). Ex: \cle{Gulf Stream} (Atl N), \cle{Kuroshio} (Pac N), \cle{Brazil Current} (Atl S), \cle{Agulhas} (Ind S).
                  \item \textbf{Courants de bord est (Eastern Boundary Currents) :} Larges, peu profonds, \textbf{lents}, \textbf{froids} (remontée d'eau profonde / upwelling). Ex: \cle{Canaries} (Atl N), \cle{Californie} (Pac N), \cle{Benguela} (Atl S), \cle{Humboldt} (Pac S).
              \end{itemize}
        \item \textbf{Analyser la circulation thermohaline (Tapis roulant / MOC) :}
              \begin{itemize}
                  \item \textbf{Moteur :} Différences de densité ($\rho$) dues à la \textbf{Température (Thermo)} et \textbf{Salinité (Haline)}.
                  \item \textbf{Sites de formation d'eau profonde (NADW, AABW) :} Mer du Labrador, Mer du Groenland, Mer de Weddell (Antarctique). Refroidissement intense + augmentation salinité (formation glace de mer $\rightarrow$ sel rejeté) $\rightarrow$ eau très dense $\rightarrow$ \textbf{plongée (convection profonde)}.
                  \item \textbf{Parcours :} Eau profonde $\rightarrow$ circulation lente au fond (siècles) $\rightarrow$ \textbf{remontée (upwelling)} principalement dans l'Océan Austral et Pacifique Nord $\rightarrow$ retour en surface vers Atlantique Nord.
              \end{itemize}
        \item \textbf{Rédiger :} "Le courant X est un courant de bord ouest/est, chaud/froid, qui relie la zone Y à la zone Z. Il est entraîné par les vents [Alizés/Vents d'Ouest] et dévié par Coriolis/continents."
    \end{enumerate}
    \textbf{Réflexe Probatoire :} Sur une carte, flèche rouge = courant chaud ; flèche bleue = courant froid. Sens rotation gyre = indicateur hémisphère. Upwelling = richesse halieutique (pêche).
\end{methode}

\begin{exoresolu}[Interprétation d'une carte des courants de l'Atlantique Nord et liaison avec le climat européen]
On donne la carte simplifiée ci-dessous montrant les principaux courants de surface de l'Atlantique Nord.
\begin{popfigure}
    \begin{tikzpicture}[scale=0.85, every node/.style={font=\small}]
        % Contours simplifiés
        \draw[popdark, thick] (-5,0) .. controls (-6,2) and (-6,6) .. (-4,8) .. controls (-2,10) and (2,10) .. (4,8) .. controls (6,6) and (6,2) .. (5,0); % Am Nord
        \draw[popdark, thick] (-5,0) .. controls (-4,-1) and (0,-2) .. (4,-1) .. controls (5,0) and (5,2) .. (4,4); % Afrique/Europe partiel
        % Courants
        % Gulf Stream (Chaud, Rouge)
        \draw[poporange, very thick, ->] (1,1) .. controls (2,2) and (3,4) .. (3.5,6) node[above right, pos=0.8, poporange] {Gulf Stream};
        \draw[poporange, thick, ->] (3.5,6) .. controls (2,7) and (0,7.5) .. (-1,7) node[above, pos=0.5, poporange] {Dérive N-A};
        % Canaries (Froid, Bleu)
        \draw[popblue, thick, ->] (-1,7) .. controls (-2,5) and (-3,3) .. (-3,1) node[left, pos=0.5, popblue] {Canaries};
        % Equatorial Nord (Chaud)
        \draw[poporange, thick, ->] (-3,1) .. controls (-1,0.5) and (1,0.5) .. (2,1) node[below, pos=0.5, poporange] {Équatorial Nord};
        % Labrador (Froid)
        \draw[popblue, thick, ->] (3.5,6) .. controls (4,7) and (4.5,7.5) .. (3.5,8) node[above right, popblue] {Labrador};
        % Latitudes
        \draw[popmuted, dashed] (-5, 3.5) -- (5, 3.5) node[right] {30°N};
        \draw[popmuted, dashed] (-5, 6.5) -- (5, 6.5) node[right] {60°N};
        % Vent Alizés / Vents d'Ouest
        \draw[popgreen, ->, dashed] (-4, 2) -- (-2, 2) node[above] {Alizés NE};
        \draw[popgreen, ->, dashed] (0, 7) -- (2, 7) node[above] {Vents d'Ouest};
    \end{tikzpicture}
    \legende{Principaux courants de surface de l'Atlantique Nord.}
\end{popfigure}
\textbf{Questions :}
\begin{enumerate}
    \item Nommez les 4 courants principaux formant le gyre subtropical de l'Atlantique Nord et indiquez pour chacun s'il est chaud ou froid.
    \item Justifiez le sens de rotation de ce gyre (horaire) en invoquant les vents dominants et la force de Coriolis.
    \item Comparez les caractéristiques (largeur, vitesse, température, profondeur) du Gulf Stream et du Courant des Canaries. Expliquez l'origine de ces différences (physique des courants de bord ouest/est).
    \item Expliquez pourquoi l'Europe de l'Ouest (France, Royaume-Uni, Norvège) bénéficie d'un climat tempéré océanique doux pour sa latitude (comparé au Canada à même latitude), en lien avec le Gulf Stream et la Dérive Nord-Atlantique.
    \item Citez une conséquence économique majeure des courants de bord est (ex: Canaries, Benguela, Humboldt) pour les pays riverains (Cameroun inclus via le courant de Guinée/Benguela).
\end{enumerate}
\tcblower
\textbf{Solution détaillée :}
\begin{enumerate}
    \item \textbf{Gyre Subtropical Atlantique Nord (Rotation Horaire) :}
          \begin{enumerate}
              \item \textbf{Courant Équatorial Nord} (Chaud) : Poussé par les Alizés NE, va vers l'Ouest.
              \item \textbf{Gulf Stream} (Chaud) : Courant de bord Ouest, rapide, étroit, transporte eau chaude vers Nord.
              \item \textbf{Dérive Nord-Atlantique} (Chaud/Tiède) : Prolongement du Gulf Stream, poussé par les Vents d'Ouest, traverse l'Atlantique vers l'Est.
              \item \textbf{Courant des Canaries} (Froid) : Courant de bord Est, lent, large, ramène eau froide vers l'Équateur.
          \end{enumerate}
    \item \textbf{Mécanisme rotation horaire :}
          \begin{itemize}
              \item Au Sud (~10°-20°N) : \textbf{Alizés de Nord-Est} soufflent vers l'Ouest. Force de Coriolis (hémisphère Nord) dévie vers la \textbf{droite} $\rightarrow$ courant vers l'Ouest (Équatorial Nord).
              \item À l'Ouest (Côte Amérique) : Accumulation d'eau $\rightarrow$ pente pression $\rightarrow$ courant remonte vers Nord (Gulf Stream). Coriolis dévie vers la droite (vers l'Est).
              \item Au Nord (~40°-60°N) : \textbf{Vents d'Ouest} soufflent vers l'Est. Coriolis dévie vers la droite (vers le Sud) $\rightarrow$ courant vers l'Est (Dérive N-A).
              \item À l'Est (Côte Afrique/Europe) : Accumulation $\rightarrow$ pente pression $\rightarrow$ courant descend vers Sud (Canaries). Coriolis dévie vers la droite (vers l'Ouest). Boucle fermée.
          \end{itemize}
    \item \textbf{Comparaison Gulf Stream vs Canaries :}
          \begin{center}
              \begin{tabularx}{\linewidth}{|l|X|X|}\hline
                  \rowcolor{popblueL} \textbf{Caractéristique} & \textbf{Gulf Stream (Bord Ouest)} & \textbf{Courant des Canaries (Bord Est)} \\ \hline
                  \textbf{Température} & \textbf{Chaud} (eau tropicale) & \textbf{Froid} (eau subpolaire/upwelling) \\ \hline
                  \textbf{Vitesse} & \textbf{Très rapide} ($\sim$ 1--2 m/s, $\sim$ 100--150 Sv) & \textbf{Lent} ($\sim$ 0,1--0,5 m/s) \\ \hline
                  \textbf{Largeur/Profondeur} & \textbf{Étroit} (50--100 km), \textbf{Profond} (1000--2000 m) & \textbf{Large} (1000 km), \textbf{Superficiel} \\ \hline
                  \textbf{Origine dynamique} & \textbf{Intensification occidentale} (Stommel/Munk) : $\beta$-effet (variation Coriolis avec latitude) + frottement latéral concentre le transport vers l'ouest. & Retour diffus, large, compensant le transport ouest. Upwelling côtier (vent parallèle à côte + Coriolis $\rightarrow$ écartement eau de surface $\rightarrow$ remontée fond). \\ \hline
              \end{tabularx}
          \end{center}
    \item \textbf{Rôle climatique pour l'Europe :}
          \begin{itemize}
              \item Le \textbf{Gulf Stream} transporte une quantité massive de chaleur sensible ($Q = \rho C_p V \Delta T$) depuis les tropiques vers les moyennes/hautes latitudes.
              \item La \textbf{Dérive Nord-Atlantique} libère cette chaleur dans l'atmosphère (flux de chaleur latente et sensible) au-dessus de l'Atlantique Nord.
              \item Les \textbf{Vents d'Ouest} dominants (cellule de Ferrel) advectent cet air réchauffé et humidifié vers le continent européen.
              \item \textbf{Résultat :} Hivers beaucoup plus doux (ex: Brest 7°C en janv vs Québec -10°C à même latitude ~47°N), ports libres de glace (Norvège jusqu'au Cap Nord), pluviométrie abondante. Sans ce "radiateur", l'Europe aurait un climat continental rigoureux type Sibérie/Canada.
          \end{itemize}
    \item \textbf{Conséquence économique (Upwelling / Courants de bord Est) :}
          Les courants de bord est (Canaries, Benguela, Humboldt, Californie, \textbf{Guinée} pour le Golfe de Guinée/Cameroun) favorisent l'\textbf{upwelling côtier}.
          \begin{itemize}
              \item Remontée d'eaux profondes riches en sels nutritifs (nitrates, phosphates, silicates).
              \item Explosion de la \textbf{production primaire} (phytoplancton) $\rightarrow$ base chaîne alimentaire.
              \item \textbf{Zones de pêche les plus productives au monde} (ex: Pérou/Chili - Anchois ; Afrique du Sud/Namibie - Sardines ; Mauritanie/Sénégal ; Cameroun - crevettes, poissons pélagiques).
              \item Enjeu majeur pour la \textbf{sécurité alimentaire} et l'économie bleue du Cameroun (pêche artisanale/industrielle à Kribi, Douala, Limbé).
          \end{itemize}
\end{enumerate}
\end{exoresolu}

\courssub{Méthode 3 : Relier la genèse d'un cyclone tropical à l'interaction Océan-Atmosphère}
\begin{methode}[Démarche pour expliquer la formation et l'intensification d'un cyclone tropical (ouragan/typhon)]
    \textbf{Objectif :} Structurer l'explication causale : conditions initiales $\rightarrow$ processus thermodynamiques $\rightarrow$ organisation dynamique $\rightarrow$ stade mature.
    \begin{enumerate}
        \item \textbf{Rappeler les 6 conditions nécessaires (Check-list de Gray) :}
              \begin{enumerate}
                  \item \textbf{Température de surface de la mer (TSM) $\ge$ 26,5°C} sur une profondeur $\ge$ 50 m (réservoir d'énergie).
                  \item \textbf{Latitudes entre 5° et 20°} (force de Coriolis $\neq$ 0 pour initier rotation, mais pas trop fort pour ne pas cisailler).
                  \item \textbf{Instabilité atmosphérique} (gradient de température vertical fort) $\rightarrow$ convection profonde.
                  \item \textbf{Humidité relative élevée} dans la moyenne troposphère (500 hPa) $\rightarrow$ limite l'entrainement d'air sec qui étouffe la convection.
                  \item \textbf{Faible cisaillement vertical du vent} (différence vitesse/direction vent 200 hPa - 850 hPa $<$ 10 m/s) $\rightarrow$ permet l'alignement vertical du vortex.
                  \item \textbf{Perturbation préexistante} (onde de l'Est, creux barométrique, ZCIT active) $\rightarrow$ amorce la convergence de surface.
              \end{enumerate}
        \item \textbf{Décrire le moteur thermodynamique (Machine de Carnot) :}
              \begin{itemize}
                  \item \textbf{Source chaude :} Océan (TSM $\ge$ 26,5°C). Évaporation massive $\rightarrow$ flux de chaleur latente ($L_v E$).
                  \item \textbf{Puits froid :} Tropopause ($\sim$ -70°C à 15 km).
                  \item \textbf{Cycle :} Air chaud/humide converge au sol $\rightarrow$ monte en tourbillonnant (paroi de l'œil) $\rightarrow$ condensation $\rightarrow$ libération chaleur latente $\rightarrow$ réchauffement cœur $\rightarrow$ baisse pression centrale $\rightarrow$ accentue gradient pression $\rightarrow$ vents plus forts $\rightarrow$ plus évaporation (rétroaction positive \cle{WISHE} -- Wind-Induced Surface Heat Exchange).
              \end{itemize}
        \item \textbf{Structure mature :} \textbf{Œil} (air subsident, sec, calme, clair) -- \textbf{Paroi de l'œil} (vents max, pluies extrêmes, convection la plus forte) -- \textbf{Bandes spiralées} (orages, tornades).
        \item \textbf{Trajectoire et fin de vie :} Déplacement vers l'Ouest/Nord-Ouest (Alizés + dérive $\beta$) $\rightarrow$ Atterrissage (Landfall) : coupure source énergie (eau) + frottement sol $\rightarrow$ affaiblissement rapide. Ou remontée en latitudes moyennes : transition extratropicale (interaction front polaire).
    \end{enumerate}
    \textbf{Réflexe Probatoire :} Ne pas confondre "création" (perturbation + conditions) et "intensification" (WISHE). Le terme "cyclone tropical" est générique ; "ouragan" (Atlantique/Nord-Est Pacifique), "typhon" (Nord-Ouest Pacifique), "cyclone" (Pacifique Sud/Océan Indien).
\end{methode}

\begin{exoresolu}[Étude de cas : Le cyclone tropical et le couplage Océan-Atmosphère]
\begin{popfigure}
    \begin{tikzpicture}[scale=0.9, every node/.style={font=\small}]
        % Coupe transversale cyclone
        \draw[popdark, thick] (-5,0) -- (5,0); % Mer
        % Œil
        \draw[poporange, thick] (-0.5,0) -- (-0.5,5) -- (0.5,5) -- (0.5,0) -- cycle;
        \node[poporange] at (0, 2.5) {ŒIL};
        \node[popmuted, font=\tiny] at (0, 1.5) {Subsidence $\downarrow$};
        \node[popmuted, font=\tiny] at (0, 0.5) {Calme, Clair};
        % Paroi œil
        \draw[popred, thick, decorate, decoration={snake, amplitude=1pt, segment length=4pt}] (-1.5,0) -- (-0.5,6);
        \draw[popred, thick, decorate, decoration={snake, amplitude=1pt, segment length=4pt}] (1.5,0) -- (0.5,6);
        \node[popred, rotate=90] at (-1.2, 3) {PAROI DE L'ŒIL};
        \node[popred, font=\tiny] at (-1.2, 1.5) {Vents Max, Pluies diluviennes};
        \draw[popblue, thick, ->] (-2,0.5) -- (-1.5,0.5) node[left] {Air chaud/humide};
        \draw[popblue, thick, ->] (2,0.5) -- (1.5,0.5) node[right] {Convergence};
        \draw[popblue, thick, ->] (-0.5,5.5) -- (-0.5,6.5) node[above] {Divergence altitude};
        \draw[popblue, thick, ->] (0.5,5.5) -- (0.5,6.5);
        % Bandes spirales
        \draw[poppurple, thick, ->] (-4,0) .. controls (-3,2) and (-2,4) .. (-1.5,5.5);
        \draw[poppurple, thick, ->] (4,0) .. controls (3,2) and (2,4) .. (1.5,5.5);
        \node[poppurple] at (-4.5, 1) {Bandes};
        \node[poppurple] at (-4.5, 0.5) {spiralées};
        % Flux chaleur océan
        \draw[poporange, thick, <->] (-3,-0.5) -- (3,-0.5) node[midway, below] {Flux chaleur latente $L_vE$ (TSM > 26,5°C)};
        \draw[popblue, thick, ->] (0,-0.5) -- (0,-1.2) node[below] {Refroidissement mer (Brassage/Upwelling)};
    \end{tikzpicture}
    \legende{Coupe schématique d'un cyclone tropical mature : structure et transferts d'énergie.}
\end{popfigure}
\textbf{Questions :}
\begin{enumerate}
    \item En vous appuyant sur le schéma, expliquez pourquoi la température de surface de la mer (TSM) doit être supérieure à 26,5°C sur une couche d'au moins 50 m pour permettre la cyclogenèse.
    \item Décrivez le cycle thermodynamique (moteur de Carnot) qui alimente le cyclone. Identifiez la source chaude, le puits froid, le fluide de travail et le travail produit.
    \item Qu'est-ce que le rétrocontrôle positif \textbf{WISHE} (\emph{Wind-Induced Surface Heat Exchange}) ? Expliquez comment il conduit à l'intensification rapide.
    \item Pourquoi le cyclone s'affaiblit-il rapidement après le \emph{landfall} (toucher terre) ? Citez deux raisons physiques distinctes.
    \item Le Cameroun est-il concerné par les cyclones tropicaux ? Justifiez votre réponse par la latitude et la configuration océanique (Golfe de Guinée). Quels phénomènes violents de convection y observe-t-on à la place (citer le nom local) ?
\end{enumerate}
\tcblower
\textbf{Solution détaillée :}
\begin{enumerate}
    \item \textbf{Rôle de la TSM $\ge$ 26,5°C / 50 m :}
          \begin{itemize}
              \item C'est la \textbf{source d'énergie} (source chaude du moteur thermique). L'évaporation nécessite une eau chaude pour fournir un flux de chaleur latente ($L_v E$) suffisant pour entretenir la convection profonde.
              \item La condition de \textbf{profondeur (50 m)} est cruciale : les vents violents du cyclone brassent la couche de surface (mélange turbulent). Si la couche chaude est mince, l'upwelling induit par le cyclone lui-même fait remonter de l'eau froide ($\ll$ 26°C) $\rightarrow$ coupe l'alimentation en chaleur/humidité $\rightarrow$ \textbf{auto-limitation} (mort du cyclone). Une couche épaisse (ex: courant chaud, Eddy chaud) permet de maintenir le flux.
          \end{itemize}
    \item \textbf{Cycle de Carnot (Emanuel, 1986) :}
          \begin{itemize}
              \item \textbf{Fluide de travail :} Air humide.
              \item \textbf{Source chaude ($T_s \approx$ 300 K) :} Surface océan (évaporation isotherme $\rightarrow$ gain enthalpie $k^*_s$).
              \item \textbf{Puits froid ($T_o \approx$ 200 K) :} Tropopause (rayonnement infrarouge vers l'espace $\rightarrow$ perte enthalpie).
              \item \textbf{Processus :} 1. Expansion adiabatique humide (montée paroi œil) $\rightarrow$ condensation $\rightarrow$ libération chaleur. 2. Détente isotherme en altitude (rayonnement). 3. Compression adiabatique sèche (subsidence dans l'œil). 4. Compression isotherme au sol (frottement, gain chaleur mer).
              \item \textbf{Travail produit :} Énergie cinétique des vents ($V_{max}^2/2$). Rendement $\eta = (T_s - T_o)/T_s \approx 1/3$.
          \end{itemize}
    \item \textbf{WISHE (Échange de chaleur de surface induit par le vent) :}
          \begin{enumerate}
              \item Perturbation initiale $\rightarrow$ convergence basse couche $\rightarrow$ montée $\rightarrow$ condensation $\rightarrow$ baisse pression centre.
              \item Baisse pression $\rightarrow$ \textbf{Gradient de pression horizontal plus fort} $\rightarrow$ \textbf{Vents de surface plus forts} ($V$).
              \item Vents plus forts $\rightarrow$ \textbf{Turbulence accrue} à l'interface air-mer $\rightarrow$ \textbf{Coefficients d'échange $C_k, C_d$ augmentent} $\rightarrow$ \textbf{Flux de chaleur latente ($L_v E$) et sensible ($H$) augmentent} massivement.
              \item Plus d'énergie apportée au cœur $\rightarrow$ pression chute encore $\rightarrow$ vents augmentent... \textbf{Boucle de rétroaction positive}.
          \end{enumerate}
          C'est le mécanisme d'\textbf{intensification rapide} (ex: passage Cat 1 $\rightarrow$ Cat 4 en 24h).
    \item \textbf{Affaiblissement au Landfall :}
          \begin{enumerate}
              \item \textbf{Coupure de la source d'énergie :} Plus d'évaporation depuis l'océan (sol terre $\rightarrow$ faible humidité, capacité thermique faible, pas de réservoir 50 m). Flux $L_v E \rightarrow 0$.
              \item \textbf{Augmentation du frottement de surface :} Rugosité terre (forêts, villes, reliefs) $\gg$ rugosité mer. Frottement $\rightarrow$ convergence de masse accrue au centre $\rightarrow$ montée forcée mais \textbf{perte d'énergie cinétique} (dissipation) + entrée d'air sec en moyenne troposphère $\rightarrow$ étouffement convection.
          \end{enumerate}
    \item \textbf{Cameroun et Cyclones :}
          \begin{itemize}
              \item \textbf{NON.} Le Cameroun (2°--13°N) est \textbf{trop au Sud} pour les cyclones de l'Atlantique Nord (genèse typique 10°--20°N, trajectoire Ouest $\rightarrow$ Nord-Ouest vers Caraïbes/USA). Le Golfe de Guinée est un "cul-de-sac" océanique ; la ZCIT y est large mais le cisaillement vertical y est souvent fort (Jet d'Est africain) et la TSM, bien que chaude, ne génère pas de vortex organisés cycloniques (Coriolis faible à l'équateur, pas de place pour aller vers l'Ouest sans toucher terre immédiatement).
              \item \textbf{Phénomènes violents locaux :} Les \textbf{Lignes de Grain (Squall Lines)} ou \textbf{Orages Multicellulaires} (Mesoscale Convective Systems -- MCS). Localement appelés \textbf{"Orages de la ligne de grain"} ou \textbf{"Vents violents d'orage"}. Ce sont des systèmes convectifs organisés (souvent propagés par le Jet d'Est Africain - AEJ), apportant vents violents (rafales > 90 km/h), grêle, pluies diluviennes (inondations éclairs), foudre intense. Très fréquents avril-juin et sept-oct (passages ZCIT).
          \end{itemize}
\end{enumerate}
\end{exoresolu}

\courssub{Méthode 4 : Relier les inondations et crues extrêmes à leur origine dynamique (Mousson, ZCIT, Lignes de Grain)}
\begin{methode}[Démarche pour analyser l'origine dynamique d'inondations catastrophiques en zone tropicale]
    \textbf{Objectif :} Différencier les mécanismes pluviométriques (mousson, onde de l'Est, ligne de grain, cyclone extratropical) et leur lien avec la circulation générale.
    \begin{enumerate}
        \item \textbf{Identifier le système météorologique générateur :}
              \begin{itemize}
                  \item \textbf{Mousson (Saison des pluies) :} Flux de Sud-Ouest (Mousson) $\rightarrow$ air maritime tropical humide + convergence ZCIT + instabilité $\rightarrow$ pluies continues/orographiques (monts Cameroun, Adamawa). \textbf{Crue lente, montante, prévisible.}
                  \item \textbf{Ondes de l'Est (African Easterly Waves - AEW) :} Perturbations baroclines sur le Jet d'Est Africain (AEJ, ~600 hPa, 15°N). Propagation Ouest $\rightarrow$ convergence avant l'onde $\rightarrow$ orages organisés (MCS). \textbf{Pluies intenses, localisées, récurrentes.}
                  \item \textbf{Lignes de Grain (Squall Lines) :} Systèmes convectifs organisés (MCS) rapides (50 km/h), durée 6--12h. Front de rafale (gust front) $\rightarrow$ soulèvement air chaud $\rightarrow$ nouvelle convection. \textbf{Pluies extrêmes en < 3h (inondations éclairs), vents violents.}
                  \item \textbf{Cyclones extratropicaux / Dépressions coupées (Cut-off lows) :} Rare au Cameroun, possible au Nord (Sahélien) en fin de saison (sept-oct) si descente d'air polaire.
              \end{itemize}
        \item \textbf{Analyser les facteurs aggravants (Physique du bassin versant) :}
              \begin{itemize}
                  \item \textbf{Antécédents d'humidité :} Sol saturé (indice d'humidité antécédent IHA) $\rightarrow$ ruissellement coefficient $\uparrow$.
                  \item \textbf{Topographie :} Pentes fortes (Mont Cameroun, Mandara) $\rightarrow$ concentration rapide, crues éclairs. Plaines (Logone, Bénoué) $\rightarrow$ inondations lentes, vastes étendues.
                  \item \textbf{Urbanisation / Déforestation :} Imperméabilisation $\rightarrow$ pic de crue $\uparrow$, temps de concentration $\downarrow$.
              \end{itemize}
        \item \textbf{Relier à la circulation générale :} Position de la ZCIT (migration Nord/Sud), force du Jet d'Est Africain (AEJ), température Atlantique Sud (dipôle Atlantique), ENSO (La Niña $\rightarrow$ mousson plus forte au Sahel/Soudan).
    \end{enumerate}
    \textbf{Réflexe Probatoire :} "Inondation" $\neq$ "Crue". Crue = débit rivière > seuil. Inondation = eau hors lit majeur. Distinguer \textbf{crue lente} (plaine, mousson) et \textbf{crue éclair} (montagne, ligne de grain). Citer le bassin versant concerné (Logone, Bénoué, Sanaga, Wouri, Mungo).
\end{methode}

\begin{exoresolu}[Analyse des inondations d'août 2012 dans le bassin du Logone (Extrême-Nord Cameroun)]
\textbf{Contexte :} En août-septembre 2012, le département du Logone-et-Chari (Extrême-Nord) a connu des inondations catastrophiques (plus de 30 000 sinistrés, cultures détruites, épidémies choléra). Les précipitations cumulées ont dépassé 400 mm en 10 jours sur le bassin versant amont (Tchad, RCA, Cameroun).
\begin{popfigure}
    \begin{tikzpicture}[scale=0.8, every node/.style={font=\small}]
        % Carte simplifiée bassin Logone
        \draw[popblue, thick] (-2,0) .. controls (-1,2) and (1,3) .. (2,4) .. controls (3,5) and (3,6) .. (2,7); % Logone
        \draw[popblue, thick] (-3,1) .. controls (-1,2) and (0,3) .. (1,4); % Affluent
        \draw[popblue, thick] (-4,0) .. controls (-2,1) and (-1,2) .. (-1,3); % Affluent
        % Flèches vents Mousson
        \draw[popgreen, thick, ->] (-5, 3) -- (-2, 3) node[above] {Mousson SW};
        \draw[popgreen, thick, ->] (-5, 5) -- (-1, 5);
        % ZCIT
        \draw[poporange, dashed, thick] (-5, 6) -- (3, 6) node[right] {ZCIT (Position Août ~12-14°N)};
        % Jet d'Est (AEJ)
        \draw[poppurple, thick, ->, dashed] (4, 4) -- (1, 4) node[above, pos=0.3] {Jet d'Est (AEJ)};
        \draw[poppurple, thick, ->, dashed] (4, 2) -- (0, 2);
        % Ondes de l'Est
        \draw[popred, thick, ->] (3, 4.5) .. controls (2, 4) and (1, 3.5) .. (0.5, 3) node[below left] {Onde de l'Est};
        \draw[popred, thick, ->] (3, 2.5) .. controls (2, 2) and (1, 1.5) .. (0.5, 1);
        % Ville
        \node[popdark, circle, fill=popred, text=white, inner sep=2pt] at (2, 5.5) {K};
        \node[popdark] at (2.5, 5.5) {Kousseri};
        \node[popdark, circle, fill=popred, text=white, inner sep=2pt] at (-1, 3) {N};
        \node[popdark] at (-0.5, 3) {N'Djamena};
        % Légende
        \node[popmuted, font=\tiny] at (-4.5, -0.5) {Tchad / RCA (Amont)};
        \node[popmuted, font=\tiny] at (2.5, -0.5) {Lac Tchad (Aval)};
    \end{tikzpicture}
    \legende{Configuration synoptique type août 2012 : Mousson active, ZCIT au nord, Ondes de l'Est sur le Jet d'Est.}
\end{popfigure}
\textbf{Questions :}
\begin{enumerate}
    \item Identifiez sur le schéma les deux principaux systèmes dynamiques responsables des pluies diluviennes d'août 2012 sur le bassin du Logone. Expliquez le mécanisme physique de chacun.
    \item Pourquoi la position de la ZCIT au nord de sa position climatologique (vers 12°-14°N au lieu de 8°-10°N) est-elle un facteur aggravant pour l'Extrême-Nord ?
    \item Le Jet d'Est Africain (AEJ) joue un rôle clé. Décrivez sa structure (niveau, vitesse, position) et expliquez comment il génère les ondes de l'Est (instabilité barocline/barotrope).
    \item En quoi la morphologie du bassin du Logone (plaine d'inondation, pente très faible, connexion Lac Tchad) transforme-t-elle une crue en catastrophe humanitaire durable ?
    \item Proposez deux mesures de prévision/alerte précoce basées sur la surveillance de la circulation atmosphérique (indicateurs dynamiques) pour anticiper ce type d'événement 5 à 10 jours à l'avance.
\end{enumerate}
\tcblower
\textbf{Solution détaillée :}
\begin{enumerate}
    \item \textbf{1. La Mousson (Flux de Sud-Ouest) :}
          \begin{itemize}
              \item \textbf{Origine :} Anticyclone de Sainte-Hélène (Atlantique Sud) $\rightarrow$ flux équatorial $\rightarrow$ franchit l'équateur $\rightarrow$ dévié par Coriolis (gauche hémisphère Sud $\rightarrow$ devient SO) $\rightarrow$ \textbf{Mousson SW}.
              \item \textbf{Mécanisme :} Air maritime tropical (AMT) chaud, humide, instable. Convergence avec l'Harmattan (NE) au niveau de la ZCIT $\rightarrow$ convergence de masse $\rightarrow$ montée généralisée $\rightarrow$ \textbf{pluies stratiformes continues + orages} sur plusieurs jours. Apport d'humidité massif sur le bassin versant.
          \end{itemize}
          \textbf{2. Les Ondes de l'Est (African Easterly Waves - AEW) :}
          \begin{itemize}
              \item \textbf{Origine :} Instabilité du \textbf{Jet d'Est Africain (AEJ)} (niveau ~600-700 hPa, noyau ~15°N, vitesse 10-15 m/s). Instabilité barocline (cisaillement vertical + gradient température méridien) et barotrope (cisaillement horizontal).
              \item \textbf{Mécanisme :} Onde de longueur 2000-4000 km, période 3-5 jours, vitesse phase ~8 m/s vers Ouest. \textbf{Avant de l'onde (côté Est) :} Convergence basse couche + divergence altitude $\rightarrow$ \textbf{déclenchement/organisation convection profonde (MCS/Lignes de grain)}. Pluies intenses, localisées, superposées aux pluies mousson $\rightarrow$ \textbf{cumul extrême}.
          \end{itemize}
    \item \textbf{ZCIT au Nord (Position anormale) :}
          \begin{itemize}
              \item Climatologiquement en août, ZCIT $\approx$ 10°-12°N. En 2012, remontée exceptionnelle vers 14°-16°N (liée à Atlantique Nord chaud / Dipôle Atlantique positif + La Niña).
              \item \textbf{Conséquence :} La zone de convergence maximale (et donc les pluies les plus intenses) se déplace \textbf{sur le bassin versant du Logone} (situé ~10°-13°N) au lieu de rester plus au sud (bassin Bénoué/Sanaga). L'Extrême-Nord, zone sahélienne habituellement semi-aride (400-600 mm/an), reçoit des volumes pluviométriques soudano-guinéens (> 1000 mm en saison) en quelques semaines $\rightarrow$ sols non adaptés, absence de végétation dense $\rightarrow$ ruissellement massif.
          \end{itemize}
    \item \textbf{Jet d'Est Africain (AEJ) et Ondes de l'Est :}
          \begin{itemize}
              \item \textbf{Structure :} Vent zonal d'Est maximum à \textbf{600-700 hPa} ($\sim$ 3-4 km), noyau centré vers \textbf{15°N}. Vitesse 10-20 m/s. Séparé du Jet Tropical d'Est (TEJ, 200 hPa, Ouest) par la couche de la mousson (850-700 hPa, Ouest).
              \item \textbf{Génèse ondes :} \textbf{Instabilité barocline} : Fort gradient de température méridien (Sahara chaud au Nord, Golfe de Guinée plus frais/humide au Sud) $\rightarrow$ cisaillement vertical du vent (vent d'Est en haut, vent d'Ouest en bas) $\rightarrow$ énergie potentielle disponible. \textbf{Instabilité barotrope} : Cisaillement horizontal du jet (courbure). $\rightarrow$ Perturbations amplifiées $\rightarrow$ Ondes de l'Est (AEW). Ces ondes sont les "graines" de la plupart des cyclones atlantiques ET des MCS sahéliens.
          \end{itemize}
    \item \textbf{Morphologie bassin Logone (Facteurs aggravants) :}
          \begin{itemize}
              \item \textbf{Plaine d'inondation immense (Yaérés) :} Pente extrêmement faible ($<$ 0,1 \%). L'eau s'étale sur des dizaines de km de large. Volume stocké énorme.
              \item \textbf{Temps de concentration long :} L'onde de crue met des semaines à descendre de l'amont (Tchad/RCA) vers l'aval (Kousseri/Lac Tchad). Inondation \textbf{longue durée} (mois).
              \item \textbf{Effet "bathtub" (baignoire) :} Le Logone se jette dans le Lac Tchad (niveau bas). Si le lac est bas, il absorbe ; si le lac monte (apports Logone + Chari), le reflux bloque l'écoulement $\rightarrow$ montée eau généralisée.
              \item \textbf{Populations vulnérables :} Densité forte sur les rives (pêche, agriculture de décrue), habitats précaires (banco), absence de digues de protection. Coupure routes (RN1) $\rightarrow$ isolement total.
          \end{itemize}
    \item \textbf{Mesures de prévision/alerte (5-10 jours) :}
          \begin{enumerate}
              \item \textbf{Surveillance de la ZCIT et TSM Atlantique :} Prévision saisonnière (Prévisions COP/ACMAD) + Suivi temps réel (satellite METEOSAT, modèles NWP : ECMWF, ARPEGE, GFS). Indicateur : \textbf{Latitude de la ZCIT} (position temps réel vs climatologie). Si ZCIT $>$ 12°N en juillet $\rightarrow$ Alerte renforcée Extrême-Nord.
              \item \textbf{Tracking des Ondes de l'Est (AEW) :} Suivi du \textbf{Jet d'Est (AEJ)} à 700 hPa (vents, cisaillement). Détection des perturbations (vorticité relative, OLR - Outgoing Longwave Radiation basse) sortant de l'Éthiopie/Darfour. Modèles de prévision numérique (NWP) donnent la trajectoire et l'intensité des ondes 5-7 jours à l'avance. \textbf{Alerte :} "Onde de l'Est active + ZCIT au Nord + Sol saturé $\rightarrow$ Risque crue éclair/majorée 3-5 jours".
          \end{enumerate}
\end{enumerate}
\end{exoresolu}

\courssub{Méthode 5 : Proposer des mesures de protection et d'adaptation face aux aléas atmosphériques et océaniques}
\begin{methode}[Démarche pour proposer une stratégie de gestion des risques (Prévention, Protection, Gestion de crise)]
    \textbf{Objectif :} Structurer une réponse opérationnelle selon le cycle de la gestion des risques : \textbf{Connaissance $\rightarrow$ Prévention $\rightarrow$ Protection $\rightarrow$ Gestion de crise $\rightarrow$ Retour d'expérience}.
    \begin{enumerate}
        \item \textbf{Évaluer le risque (Aléa $\times$ Enjeux $\times$ Vulnérabilité) :}
              \begin{itemize}
                  \item \textbf{Aléa :} Carte d'aléa (zones inondables, vents cycloniques, submersion marine, érosion côtière). Scénarios (crue centennale, cyclone Cat 3).
                  \item \textbf{Enjeux :} Population, habitats, infrastructures (routes, hôpitaux, écoles, barrages, centrales), activités économiques (pêche, agriculture, port).
                  \item \textbf{Vulnérabilité :} Type d'habitat (dur/précaire), niveau de vie, conscience du risque, capacité d'évacuation.
              \end{itemize}
        \item \textbf{Mesures de Prévention (Amont - Réduire l'aléa ou l'exposition) :}
              \begin{itemize}
                  \item \textbf{Urbanisme / POS / PLU :} \textbf{Zone inconstructible} (lit majeur, bande côtière 100m, couloirs d'écoulement). \textbf{Zone constructible sous conditions} (surélévation, matériaux résistants, refuge au dernier étage).
                  \item \textbf{Aménagement du bassin versant :} Reboisement (captation eau, stabilisation sols), agriculture de conservation (couvert végétal), zones d'expansion de crues (ZEC) préservées/restaurées (zones humides).
                  \item \textbf{Littoral :} Gestion intégrée (GIEC). Préservation mangroves (brise-vagues naturel, nurserie), dunes, récifs coralliens. Recul stratégique (relocalisation) vs endiguement (effet digue $\rightarrow$ érosion aval).
              \end{itemize}
        \item \textbf{Mesures de Protection (Physiques - Réduire la vulnérabilité) :}
              \begin{itemize}
                  \item \textbf{Ouvrages :} Digues, barrages de rétention (ex: Lagdo, Memve'ele - gestion des lâchers), murs de soutènement, brise-lames, seuils de ralentissement.
                  \item \textbf{Bâtiment :} Normes parasismiques / paracycloniques (toitures ancrées, ouvertures protégées), surélévation habitations (pilotis), points d'eau/sanitation hors d'eau.
              \end{itemize}
        \item \textbf{Mesures de Gestion de Crise / Alerte (Organisationnelles) :}
              \begin{itemize}
                  \item \textbf{Système d'Alerte Précoce (SAP) :} Surveillance (météo, hydrométrie, satellite) $\rightarrow$ Prévision (modèles) $\rightarrow$ \textbf{Message d'alerte codifié} (Vigilance : Vert/Jaune/Orange/Rouge) $\rightarrow$ Diffusion (radio, SMS, sirènes, relais communautaires).
                  \item \textbf{Plans d'urgence :} PCS (Plan Communal de Sauvegarde), ORSEC. Itinéraires/aires d'évacuation, abris identifiés, stocks vivres/médicaments, exercices simulés.
                  \item \textbf{Assurance / Fonds de solidarité :} Assurance indiciaire (météo), fonds calamités.
              \end{itemize}
        \item \textbf{Adaptation au Changement Climatique (Long terme) :} Intégration scénarios GIEC (RCP/SSP) dans l'aménagement (élévation niveau mer +0,5-1m 2100, intensification pluies extrêmes, fréquence cyclones ?).
    \end{enumerate}
    \textbf{Réflexe Probatoire :} Distinguer \textbf{Prévention} (éviter le risque : ne pas construire) et \textbf{Protection} (se protéger du risque : digue). La "meilleure" protection est la prévention (non-exposition). Citer le \textbf{Plan National de Gestion des Risques de Catastrophes (PNGRC)} du Cameroun et la \textbf{Plateforme Nationale de Réduction des Risques de Catastrophes (PNRRC)}.
\end{methode}

\begin{exoresolu}[Élaboration d'un Plan Communal de Sauvegarde (PCS) face au risque d'inondation côtière à Limbé (Sud-Ouest)]
\textbf{Situation :} La ville de Limbé (capitale économique Sud-Ouest, $\sim$ 150 000 hab., raffinerie SONARA, port, plage touristique) est située au pied du Mont Cameroun, sur une côte basse exposée aux vagues de tempête, à l'érosion côtière (recul trait de côte $\sim$ 1-2 m/an) et aux crues torrentielles des rivières (Bota, Mokundi) lors des pluies extrêmes de mousson (juillet-septembre).
\begin{popfigure}
    \begin{tikzpicture}[scale=0.85, every node/.style={font=\small}]
        % Profil transversal Limbé
        \draw[popdark, thick] (-4,0) -- (4,0); % Niveau mer
        % Mer
        \fill[popblueL] (-4,0) rectangle (4,-2);
        \node[popblue] at (0,-1) {Océan Atlantique};
        % Plage / Dune
        \draw[popgold, thick] (-3,0) -- (-2,1) -- (-1,0);
        \node[popgold] at (-2, 0.5) {Plage/Dune};
        % Zone urbaine basse (Risque)
        \fill[poppinkL] (-1,0) rectangle (2,2);
        \node[popdark] at (0.5, 1) {Ville Basse (Centre, SONARA, Port)};
        \node[popred, font=\tiny] at (0.5, 2.2) {COTE 0-5 m : SUBMERSION / INONDATION};
        % Rivière
        \draw[popblue, thick] (1.5,0) .. controls (1.5,1) and (1,2) .. (1,3);
        \node[popblue] at (1.8, 1.5) {Rivière Bota};
        % Mont Cameroun (Pente forte)
        \draw[popgreen, thick] (2,0) .. controls (2.5,2) and (3,4) .. (3.5,6);
        \node[popgreen] at (3.5, 3) {Mont Cameroun (Pentes fortes)};
        \node[popred, font=\tiny] at (3, 2) {Ruissellement torrentiel / Lahars};
        % Mesures
        \draw[poporange, thick, ->] (-3.5, -1.5) -- (-2.5, 0.5) node[left] {Brise-lames / Récifs artificiels};
        \draw[poporange, thick, ->] (-1.5, 2.5) -- (-0.5, 2.5) node[above] {Digue de protection / Mur};
        \draw[poporange, thick, ->] (1, 3) -- (1.5, 2) node[right] {Bassin de rétention amont};
        \draw[poporange, thick, ->] (0.5, 2.5) -- (0.5, 3.5) node[above] {Évacuation vers Hauteurs (Buea/Mutengene)};
    \end{tikzpicture}
    \legende{Profil de risque côtier à Limbé : aléa multiple (submersion, crue torrentielle, érosion).}
\end{popfigure}
\textbf{Questions :}
\begin{enumerate}
    \item Identifiez les trois aléas majeurs couplés qui menacent Limbé (origine atmosphérique/océanique/géologique).
    \item Pour chaque aléa, proposez une mesure de \textbf{prévention} (urbanisme/aménagement) et une mesure de \textbf{protection} (ouvrage/technique).
    \item Décrivez l'organisation du \textbf{Système d'Alerte Précoce (SAP)} spécifique à Limbé : capteurs, modèles, seuils d'alerte (vigilance), canaux de diffusion adaptés au contexte local (populations hétérogènes, zone portuaire, raffinerie).
    \item Le recul du trait de côte (érosion) s'accélère avec l'élévation du niveau marin. Proposez une stratégie d'adaptation "Nature-based Solutions" (SbN) combinée à une mesure "grise" pour protéger la raffinerie SONARA et la plage touristique sur 50 ans.
    \item En cas de crue torrentielle soudaine (ligne de grain sur le Mont Cameroun), le temps de concentration du bassin versant est de 30 min. Quelle organisation de crise (PCS) permet de sauver des vies dans ce délai ultra-court ?
\end{enumerate}
\tcblower
\textbf{Solution détaillée :}
\begin{enumerate}
    \item \textbf{Trois aléas couplés :}
          \begin{enumerate}
              \item \textbf{Submersion marine / Vague de tempête :} Origine \textbf{océanique/atmosphérique}. Dépressions tropicales (ou cyclones passant au large), forts vents d'Ouest/SO $\rightarrow$ surcote + déferlement vagues $\rightarrow$ franchissement cordon dunaire.
              \item \textbf{Crue torrentielle (Crues éclairs) :} Origine \textbf{atmosphérique/orographique}. Lignes de grain / Orages stationnaires sur flanc Mont Cameroun (pentes > 30\%, sols volcaniques perméables mais minces) $\rightarrow$ ruissellement ultra-rapide, charriage blocs/laves (lahars froids) $\rightarrow$ débouché brutal en ville basse.
              \item \textbf{Érosion côtière chronique :} Origine \textbf{océanique/anthropique}. Houle de Sud-Ouest dominante + montée niveau mer + déficit sédimentaire (barrages amont, extraction sable, port de Limbé piégeant sédiments) $\rightarrow$ recul trait de côte $\rightarrow$ fragilisation digues/habitats.
          \end{enumerate}
    \item \textbf{Mesures Prévention / Protection :}
          \begin{center}
              \begin{tabularx}{\linewidth}{|l|X|X|}\hline
                  \rowcolor{popblueL} \textbf{Aléa} & \\textbf{Prévention (Exposition/Vulnérabilité)} & \textbf{Protection (Ouvrages/Résistance)} \\ \hline
                  \textbf{Submersion} & \textbf{Zonage :} Zone inconstructible bande 100-200m (Loi littoral). Relocalisation progressive quartiers précaires (Bota, Down Beach) vers hauteurs (Buea, Mutengene). & \textbf{Digue / Mur de soutènement} dimensionné surcote centennale + vague. \textbf{Brise-lames émergés/submergés} pour casser énergie houle avant digue. \\ \hline
                  \textbf{Crue torrentielle} & \textbf{Bassin versant :} Reboisement pentes Mont Cameroun (espèces racines pivotantes). \textbf{ZEC :} Préservation zones d'expansion naturelles en pied de pente (marais, plaines). Interdiction construction lits mineurs. & \textbf{Bassins de rétention / Dissipateurs} en amont zone urbaine. \textbf{Calibrage/Endiguement} rivières Bota/Mokundi (section hydraulique crue 100 ans). \textbf{Grilles/Barrages anti-charriage} (blocs volcaniques). \\ \hline
                  \textbf{Érosion côtière} & \textbf{Gestion sédimentaire :} Recharge artificielle plage (sable dragué port/large). Arrêt extraction sable plage. \textbf{Recul stratégique :} Acceptation recul zones non vitales. & \textbf{Enrochements / Géotextiles} protection points durs (SONARA, Port). \textbf{Restauration Mangroves} (si estuaire) / \textbf{Végétation dunaire} (Ipomoea, Spinifex) fixation sable. \\ \hline
              \end{tabularx}
          \end{center}
    \item \textbf{SAP Limbé (Organisation) :}
          \begin{itemize}
              \item \textbf{Capteurs :} Marégraphe temps réel (port), Pluviomètres automatiques (pentes Mont Cameroun + ville), Radars météo (Douala/Yaoundé), Satellite (MSG/OLR), Stations hydrométriques rivières (niveau/débit).
              \item \textbf{Modèles/Seuils :} Modèle houle/surcote (Météo-France/ACMAD) + Modèle pluie-débit (GRP/HEC-HMS) bassin Bota.
                  \begin{itemize}
                      \item \textbf{Vert :} Vigilance routine.
                      \item \textbf{Jaune :} Pluie > 50mm/3h OU Houle > 3m + Coeff marée > 90. Info population.
                      \item \textbf{Orange :} Pluie > 80mm/3h OU Niveau rivière > Côte alerte + Houle > 4m. Pré-positionnement secours, alerte SONARA/Port.
                      \item \textbf{Rouge :} Pluie > 100mm/3h OU Rupture digue imminente OU Submersion en cours. \textbf{ÉVACUATION IMMÉDIATE} zones basses.
                  \end{itemize}
              \item \textbf{Diffusion :} Sirènes (quartiers denses), SMS Cell Broadcast (tous opérateurs), Radio locale (Ocean City Radio, CBS), Relais chefs de quartiers/blocs (tradition orale), WhatsApp groupes leaders jeunesse/femmes. \textbf{Langues :} Français, Pidgin, Bakweri.
          \end{itemize}
    \item \textbf{Stratégie Adaptation 50 ans (SbN + Gris) :}
          \begin{itemize}
              \item \textbf{SbN (Fondamentale) :} \textbf{Restauration massive de la mangrove} dans l'estuaire du Wouri/Moungo (zone tampon) et \textbf{Revégétalisation du cordon dunaire} (espèces pionnières fixatrices). Création d'\textbf{îlots artificiels/récifs} au large pour dissiper l'énergie de la houle avant la côte (création habitat halieutique bonus).
              \item \textbf{Gris (Ponctuel/Stratégique) :} \textbf{Mur de soutènement / Enrochement "dur" UNIQUEMENT} pour la protection de la \textbf{Raffinerie SONARA} (enjeu vital économique/sécurité nationale, non déplaçable à court terme) et du \textbf{Port}. Intégration d'échelles à poissons dans les ouvrages.
              \item \textbf{Gouvernance :} Observatoire littoral régional (suivi trait de côte LiDAR/satellite annuel). Fonds de compensation pour relocalisation progressive populations "Down Beach".
          \end{itemize}
    \item \textbf{Gestion de crise Crue Éclair (T < 30 min) :}
          \begin{itemize}
              \item \textbf{Impossible d'évacuer massivement par la route} (embouteillages, ponts submergés).
              \item \textbf{Stratégie "Refuge Vertical" (Vertical Evacuation) :} Identification/Construction de \textbf{bâtiments-refuges} en dur (écoles, mairies, églises, stations-service SONARA) \textbf{surélevés} (R+1 ou R+2 accessibles, toits plats accessibles), signalisés, équipés (eau, trousse secours, radio VHF).
              \item \textbf{Alerte "Dernière minute" :} Sirènes à tonalité distincte (ex: 3 sons longs = Crue Éclair). Message unique : "\textbf{Montez au point haut le plus proche MAINTENANT ! Ne prenez pas la voiture !}".
              \item \textbf{Exercices mensuels} (saison des pluies) avec écoles, marchés, quartiers. Désignation "Veilleurs de quartier" (jeunes formés) pour guider vers refuges.
              \item \textbf{Coordination PCS :} Mairie + SONARA + Protection Civile + Forces de l'ordre. Poste de commandement mobile (véhicule 4x4 radio).
          \end{itemize}
\end{enumerate}
\end{exoresolu}

\courssub{Méthode 6 : Analyser la vulnérabilité d'un territoire face aux aléas climatiques (Diagnostic territorial)}
\begin{methode}[Démarche pour réaliser un diagnostic de vulnérabilité territoriale (Approche par les compétences)]
    \textbf{Objectif :} Croiser données physiques (aléas), données socio-économiques (enjeux) et gouvernance pour hiérarchiser les secteurs d'action.
    \begin{enumerate}
        \item \textbf{Cartographier l'Aléa (Aléa = Probabilité $\times$ Intensité) :}
              \begin{itemize}
                  \item Collecte données historiques (archives, journaux, mémoire orale), données instrumentales (stations météo/hydro, satellite), modélisation (hydraulique 1D/2D, submersion, vent).
                  \item Production \textbf{Cartes d'aléa} par phénomène (Inondation, Glissement, Sécheresse, Érosion, Canicule) avec scénarios de retour (20, 50, 100 ans) et scénarios climat futurs (RCP 4.5 / 8.5).
              \end{itemize}
        \item \textbf{Cartographier les Enjeux (Exposition) :}
              \begin{itemize}
                  \item \textbf{Résidentiel :} Densité, type habitat (dur/précaire), étages.
                  \item \textbf{Activités :} Agricoles (cultures, élevage), Industrielles (usines, stockages dangereux), Commerciales, Touristiques.
                  \item \textbf{Infrastructures vitales :} Hôpitaux, écoles, réseaux (eau, élec, télécom, routes, ponts), barrages.
                  \item \textbf{Patrimoine :} Culturel, biodiversité (aires protégées).
              \end{itemize}
        \item \textbf{Évaluer la Vulnérabilité (Sensibilité + Manque de capacité d'adaptation) :}
              \begin{itemize}
                  \item \textbf{Physique :} Âge bâti, matériaux, fondations, conformité normes.
                  \item \textbf{Sociale :} Pauvreté, taux dépendance (personnes âgées, handicapés, enfants), niveau éducation, genre (femmes cheffes famille), migrants/PDI (Personnes Déplacées Internes - contexte NO/SO/Extrême-Nord).
                  \item \textbf{Économique :} Assurance, épargne, diversification revenus, accès crédit.
                  \item \textbf{Institutionnelle :} Existence PCS, PCDR (Plan Communal Développement Rural), budget risque, personnel formé, coordination intercommunale.
              \end{itemize}
        \item \textbf{Calculer l'Indice de Risque (par maille ou entité administrative) :}
              \[ R = A \times E \times V \]
              (Souvent score qualitatif : Faible / Moyen / Fort / Très Fort). Priorisation : "Hotspots" (Risque Très Fort).
        \item \textbf{Proposer un Plan d'Action Priorisé (PAP) :} Mesures "Sans regret" (bénéfiques quel que soit le climat futur) $\rightarrow$ Mesures "Faible regret" $\rightarrow$ Mesures structurelles lourdes. Intégration dans le PCD (Plan Communal de Développement) et le PNGRC.
    \end{enumerate}
    \textbf{Réflexe Probatoire :} La vulnérabilité n'est pas l'aléa. Un aléa fort sur zone désertique = Risque faible. Aléa modéré sur ville dense précaire = Risque très fort. Au Cameroun, intégrer le contexte \textbf{conflit/insécurité} (régions NO, SO, Extrême-Nord) qui détruit la résilience (déplacements, destruction infrastructures, perte moyens existence).
\end{methode}

\begin{exoresolu}[Diagnostic de vulnérabilité de la ville de Maroua (Extrême-Nord) face aux inondations et sécheresses]
\textbf{Données contextuelles :} Maroua (chef-lieu Région Extrême-Nord, $\sim$ 350 000 hab.). Climat soudano-sahélien (600-900 mm/an, 3-4 mois pluies). Rivière Kaliao (traverse la ville, lit mineur étroit, lit majeur urbanisé). Nappe phréatique superficielle. Croissance urbaine anarchique. Contexte insécurité Boko Haram (afflux PDI). Sol latéritique durci (crusting) $\rightarrow$ ruissellement fort.
\begin{popfigure}
    \begin{tikzpicture}[scale=0.85, every node/.style={font=\small}]
        % Matrice Risque simplifiée (Aléa x Vulnérabilité)
        \def\cellsize{2.2cm}
        \foreach \i/\al in {1/Faible, 2/Moyen, 3/Fort, 4/Très Fort} {
            \node[popmuted, rotate=90, anchor=south] at (-\cellsize*0.6, \i*\cellsize - \cellsize/2) {\al};
        }
        \node[popmuted, rotate=90] at (-\cellsize*0.6, 0) {ALÉA};
        \foreach \j/\vul in {1/Faible, 2/Moyenne, 3/Forte, 4/Très Forte} {
            \node[popmuted] at (\j*\cellsize - \cellsize/2, -\cellsize*0.6) {\vul};
        }
        \node[popmuted] at (0, -\cellsize*0.6) {VULNÉRABILITÉ};
        % Grille
        \foreach \i in {1,2,3,4} {
            \foreach \j in {1,2,3,4} {
                \pgfmathsetmacro{\riskval}{(\i+\j)/2}
                \pgfmathtruncatemacro{\risksum}{\i+\j}
                \ifnum\risksum<6 \def\col{popgreenL}\fi
                \ifnum\risksum=6 \def\col{popgoldL}\fi
                \ifnum\risksum>6 \def\col{poporangeL}\fi
                \ifnum\risksum>7 \def\col{poppinkL}\fi
                \fill[\col] (\j*\cellsize, \i*\cellsize) rectangle (\j*\cellsize+\cellsize, \i*\cellsize+\cellsize);
                \node[popdark, font=\tiny] at (\j*\cellsize+\cellsize/2, \i*\cellsize+\cellsize/2) {R\pgfmathprintnumber[fixed, precision=1]{\riskval}};
            }
        }
        % Légende
        \node[popdark, font=\small] at (5*\cellsize, 4*\cellsize) {Risque Faible};
        \node[popdark, font=\small] at (5*\cellsize, 3*\cellsize) {Risque Modéré};
        \node[popdark, font=\small] at (5*\cellsize, 2*\cellsize) {Risque Élevé};
        \node[popdark, font=\small] at (5*\cellsize, 1*\cellsize) {Risque Critique};
        % Position Maroua Inondation / Sécheresse
        \draw[popred, very thick] (2*\cellsize, 4*\cellsize) rectangle (3*\cellsize, 5*\cellsize); % Aléa Fort, Vuln Très Forte -> Inondation
        \node[popred, font=\small] at (2.5*\cellsize, 4.5*\cellsize) {INONDATION};
        \draw[poporange, very thick] (3*\cellsize, 3*\cellsize) rectangle (4*\cellsize, 4*\cellsize); % Aléa Très Fort, Vuln Forte -> Sécheresse
        \node[poporange, font=\small] at (3.5*\cellsize, 3.5*\cellsize) {SÉCHERESSE};
    \end{tikzpicture}
    \legende{Matrice de risque qualitative pour Maroua : Positionnement Inondation (Aléa Fort, Vuln Très Forte) et Sécheresse (Aléa Très Fort, Vuln Forte).}
\end{popfigure}
\textbf{Questions :}
\begin{enumerate}
    \item En vous appuyant sur la matrice, justifiez pourquoi le risque d'inondation à Maroua est classé "Critique" (Aléa Fort, Vulnérabilité Très Forte) alors que la pluviométrie est modérée (700 mm/an).
    \item Analysez la vulnérabilité "Très Forte" de Maroua face à l'inondation en distinguant les 4 dimensions : Physique, Sociale, Économique, Institutionnelle. Donnez un exemple concret marouan pour chacune.
    \item La sécheresse a un Aléa "Très Fort" (fréquence/intensité croissante) mais une Vulnérabilité "Forte" (pas "Très Forte"). Expliquez cette différence de vulnérabilité par rapport à l'inondation (rôle de l'adaptation agricole/pastorale traditionnelle).
    \item Proposez 3 actions prioritaires (court/moyen/long terme) pour faire baisser le risque d'inondation de "Critique" à "Élevé" en agissant sur la vulnérabilité (pas sur l'aléa climatique).
    \item Comment l'insécurité (Boko Haram) et l'afflux de PDI modifient-ils la matrice de risque (Aléa, Vulnérabilité, Capacité d'adaptation) ?
\end{enumerate}
\tcblower
\textbf{Solution détaillée :}
\begin{enumerate}
    \item \textbf{Justification Risque Inondation Critique :}
          \begin{itemize}
              \item \textbf{Aléa Fort :} Bien que pluviométrie modérée, \textbf{concentration temporelle} extrême (80\% en 3 mois : juil-sept). Orages violents (Lignes de grain) $\rightarrow$ intensités > 50 mm/h. \textbf{Bassin versant Kaliao petit, pentu, imperméabilisé} (latérite + urbanisation) $\rightarrow$ temps de concentration très court ($<$ 1h) $\rightarrow$ \textbf{crues éclairs violentes} (débits spécifiques énormes). Fréquence : quasi-annuelle (retour 2-5 ans pour crue dommageable).
              \item \textbf{Vulnérabilité Très Forte :} Convergence de tous les facteurs d'aggravation (voir Q2).
          \end{itemize}
          Produit Aléa $\times$ Vulnérabilité = Maximum.
    \item \textbf{Dimensions Vulnérabilité Inondation Maroua :}
          \begin{enumerate}
              \item \textbf{Physique :} \textbf{Habitat en banco (terre crue)} non stabilisé, fondations inexistantes, murs non étanches. \textbf{Ex :} Quartiers Doualaré, Domayo, Mairi : maisons effondrées chaque année (2012, 2019, 2022). Réseau drainage inexistant ou bouché (déchets plastiques, sable).
              \item \textbf{Sociale :} \textbf{Pauvreté extrême} (région la plus pauvre Cameroun). \textbf{Femmes cheffes de ménage} nombreuses (hommes partis/élevages/morts). \textbf{PDI (Boko Haram)} : camps précaires en zone inondable (bords Kaliao), pas de réseaux sociaux locaux, barrière langue. \textbf{Ex :} Camp de Minawao (proche) ou sites spontanés Maroua : latrines inondées $\rightarrow$ choléra.
              \item \textbf{Économique :} \textbf{Aucune assurance} habitation/récolte. Épargne nulle. Revenus journaliers (petit commerce, moto-taxi) stoppés par crue. Perte stock mil/sorgho/arachide en magasins au sol. Endettement usurier post-crue.
              \item \textbf{Institutionnelle :} \textbf{PCS inexistant ou non opérationnel} (manque budget, personnel non formé). Mairie débordée. Coordination avec MINAT/MINEE/Protection Civile faible. Pas de sirènes, pas de refuges identifiés. Gestion déchets défaillante (bouche caniveaux).
          \end{enumerate}
    \item \textbf{Sécheresse : Vulnérabilité "Forte" < "Très Forte" (Inondation) :}
          \begin{itemize}
              \item \textbf{Adaptation traditionnelle forte :} Systèmes agropastoraux résilients : \textbf{Cultures associées} (mil/niébé/sorgho), \textbf{Variétés courtes cycles} (mil 90 jours), \textbf{Agriculture de récession} (bords Mayo Tsanaga, Logone), \textbf{Transhumance} (mobilité troupeaux), \textbf{Greniers traditionnels} (stockage inter-annuel).
              \item \textbf{Filets de sécurité :} Réseaux solidarité familiale/diaspora, programmes ONG/Etat (distribution semences améliorées, intrants, alimentation bétail).
              \item \textbf{Différence :} L'inondation est \textbf{soudaine, destructrice d'actifs physiques (maisons, routes), contaminante (eau sale)}. La sécheresse est \textbf{lente, prévisible} (alertes saisonnières), gérable par ajustement progressif des pratiques. L'inondation détruit le capital physique (maison = actif principal pauvre) ; la sécheresse réduit le flux de revenus/récolte mais l'actif foncier reste. D'où vulnérabilité perçue/réelle plus aiguë pour l'inondation.
          \end{itemize}
    \item \textbf{3 Actions prioritaires pour baisser Risque Inondation (Critique $\rightarrow$ Élevé) :}
          \begin{enumerate}
              \item \textbf{Court terme (Urgence - Vulnérabilité Sociale/Physique) :} \textbf{Opération "Zéro habitat en lit majeur Kaliao"}. Délimitation stricte zone rouge (GPS + bornes). \textbf{Relocalisation assistée} populations PDI/précaires vers sites sûrs (aménagés eau/latrines/écoles) \textbf{AVANT} saison des pluies (mars-avril). Nettoyage mécanique d'urgence caniveaux existants (cash-for-work jeunes/PDI). Kits d'urgence familles (bidons, pastilles chlore, moustiquaires) pré-positionnés.
              \item \textbf{Moyen terme (Vulnérabilité Physique/Institutionnelle) :} \textbf{Plan de Drainage Urbain Intégré (PDUI) Maroua}. Construction \textbf{collecteurs principaux} dimensionnés crue 20 ans (béton armé) + bassins de rétention en amont ville (zones non bâties). \textbf{Réglementation :} Permis de construire conditionné à raccordement drainage + surélévation fondation (0,5m). \textbf{Opérationnalisation PCS :} Recrutement/formation 2 agents risque mairie, achat sirènes (solaire), simulation annuelle, convention Croix-Rouge/Protection Civile.
              \item \textbf{Long terme (Vulnérabilité Économique/Structurelle) :} \textbf{Assurance indiciaire inondation communale} (paramétrique : déclenchement si pluviométrie station Maroua > seuil 100mm/24h ou niveau Kaliao > cote alerte). Fonds mutualisé Etat/Commune/Partenaires pour indemnisation rapide (reconstruction banco stabilisé, relance commerce). \textbf{Aménagement bassin versant Kaliao amont :} Reboisement pentes (Acacia, Jatropha), seuils de ralentissement (pierres), agriculture de conservation (zaï, demi-lunes) $\rightarrow$ réduction pic crue à la source.
          \end{enumerate}
    \item \textbf{Impact Insécurité / PDI sur la Matrice de Risque :}
          \begin{itemize}
              \item \textbf{Aléa :} \textbf{Indirectement accru}. Destruction couverture végétale (bois de chauffe, défrichement survie) $\rightarrow$ ruissellement accru, érosion amont $\rightarrow$ crue plus violente/rapide. Pollution cours d'eau (déchets, cadavres) $\rightarrow$ aléa sanitaire couplé.
              \item \textbf{Vulnérabilité :} \textbf{Explosion "Très Forte $\rightarrow$ Critique"}.
                  \begin{itemize}
                      \item \textbf{Physique :} Habitats de fortune (paillotes, tôles récupérées) en zone inondable (seul foncier dispo/gratuit). Aucune résistance.
                      \item \textbf{Sociale :} Traumatismes, perte repères, barrières linguistiques/culturelles, isolement (pas de chef de quartier légitime). Femmes/Enfants non accompagnés majoritaires.
                      \item \textbf{Économique :} Zéro capital, dépendance aide humanitaire (vivres), pas d'accès terre/emploi. Endettement immédiat.
                      \item \textbf{Institutionnelle :} Mairie débordée (mandat/expertise/budget). ONG humanitaires (OCHA, UNHCR, IOM) gèrent camps mais coordination risque catastrophe (DRR) souvent absente de leur mandat initial. \textbf{Conflit d'usage sol} : zones refuges = zones inondables.
                  \end{itemize}
              \item \textbf{Capacité d'adaptation / Résilience :} \textbf{Effondrement}. Mécanismes traditionnels (solidarité clanique, chefferies) brisés. Dépendance aide externe. Perte savoir-faire local. \textbf{Risque de piège à pauvreté :} Crue $\rightarrow$ perte tout $\rightarrow$ aide $\rightarrow$ reconstruction même lieu $\rightarrow$ crue suivante...
          \end{itemize}
          \textbf{Conclusion :} L'insécurité transforme un risque naturel gérable en \textbf{risque complexe systémique (Naturel + Anthropique + Humanitaire)} nécessitant une approche \textbf{Nexus Humanitaire-Développement-Paix (HDP)}.
\end{enumerate}
\end{exoresolu}

\courssub{Méthode 7 : Calculer et interpréter des indicateurs climatiques extrêmes (Indices ETCCDI) pour la détection de tendance}
\begin{methode}[Démarche pour calculer et analyser les indices d'extrêmes climatiques (Précipitations / Température)]
    \textbf{Objectif :} Passer de la donnée brute (journalière) à l'indicateur synthétique robuste, détecter la tendance (test statistique) et l'attribuer.
    \begin{enumerate}
        \item \textbf{Choisir les indices pertinents (ETCCDI - Expert Team on Climate Change Detection and Indices) :}
              \begin{itemize}
                  \item \textbf{Précipitations :}
                        \begin{itemize}
                            \item \textbf{RX1day / RX5day :} Précipitation max 1 jour / 5 jours consécutifs (Intensité extrême).
                            \item \textbf{R95p / R99p :} Précipitation totale jours $>$ 95ème / 99ème percentile (Part des extrêmes).
                            \item \textbf{CDD / CWD :} \textbf{Consecutive Dry Days / Wet Days} (Durée max sèches/humides $\rightarrow$ Sécheresse/Inondation).
                            \item \textbf{PRCPTOT :} Total annuel (Moyenne).
                            \item \textbf{SDII :} Simple Daily Intensity Index (PRCPTOT / Nb jours pluie $\ge$ 1mm).
                        \end{itemize}
                  \item \textbf{Température :}
                        \begin{itemize}
                            \item \textbf{TXx / TNn :} Max des max / Min des min (Extrêmes absolus).
                            \item \textbf{TX90p / TN10p :} \% jours $>$ 90ème percentile TX / $<$ 10ème percentile TN (Fréquence canicules/nuits froides).
                            \item \textbf{WSDI / CSDI :} \textbf{Warm / Cold Spell Duration Indicator} (Durée vagues chaleur/froid).
                            \item \textbf{DTR :} Diurnal Temperature Range (TX - TN moyen).
                        \end{itemize}
              \end{itemize}
        \item \textbf{Calculer sur série homogène (min 30 ans, idéalement 50+) :}
              \begin{itemize}
                  \item Contrôle qualité (valeurs aberrantes, ruptures d'homogénéité -- test SNHT, Pettitt).
                  \item Calcul percentiles sur période de référence climatologique (ex: 1981-2010 ou 1991-2020) \textbf{fenêtre glissante 5 jours} pour lisser cycle annuel.
                  \item Calcul indices annuels (ou saisonniers : DJF, MAM, JJA, SON).
              \end{itemize}
        \item \textbf{Analyser la tendance :}
              \begin{itemize}
                  \item \textbf{Pente de Sen (Theil-Sen) :} Estimateur robuste de la magnitude (mm/an, jours/an, °C/an). Non paramétrique, résistant aux outliers.
                  \item \textbf{Test de Mann-Kendall (MK) :} Test non paramétrique de significativité de la tendance monotone ($H_0$ : pas de tendance). Donne valeur $p$ (seuil 0,05) et $\tau$ de Kendall (force liaison).
                  \item \textbf{Rupture de tendance :} Test de Pettitt ou Buishand pour détecter année de basculement (ex: 1998, 2010).
              \end{itemize}
        \item \textbf{Interpréter et Attribuer :} Lier tendance observée aux forçages (GHG, Aérosols, Variabilité naturelle : AMO, ENSO, AMO) via modèles (CMIP6) ou analyses physiques (ex: augmentation TX90p + TXx + WSDI = signature forçage radiatif GES).
    \end{enumerate}
    \textbf{Réflexe Probatoire :} Ne pas confondre \textbf{RX1day} (valeur max absolue) et \textbf{R95p} (somme sur jours extrêmes). Unités : mm, jours, °C. Significativité $\neq$ magnitude forte. Toujours citer la période de référence pour les percentiles.
\end{methode}

\begin{exoresolu}[Analyse de tendance des précipitations extrêmes à Douala (1960-2020) -- Type Probatoire]
\textbf{Données :} Série journalière pluviométrie aéroport Douala (1960-2020). Calcul des indices annuels RX1day et R95p (référence percentiles 1981-2010). Résultats partiels :
\begin{center}
    \begin{tabular}{|c|c|c|c|c|}\hline
        \rowcolor{popblueL} \textbf{Période} & \textbf{RX1day moyen (mm)} & \textbf{Pente Sen RX1day (mm/an)} & \textbf{R95p moyen (mm)} & \textbf{Pente Sen R95p (mm/an)} \\ \hline
        1960-1990 & 85,2 & +0,12 & 420 & +1,8 \\ \hline
        1991-2020 & 98,7 & +0,45 & 580 & +4,2 \\ \hline
        \textbf{1960-2020 (Total)} & \textbf{91,5} & \textbf{+0,28} & \textbf{495} & \textbf{+2,9} \\ \hline
    \end{tabular}
\end{center}
\textbf{*} Significatif au seuil 5\% (Test Mann-Kendall $p < 0,05$).
\textbf{Questions :}
\begin{enumerate}
    \item Définissez précisément les indices \textbf{RX1day} et \textbf{R95p}. Quelle information complémentaire apporte R95p par rapport à RX1day ?
    \item Interprétez l'évolution des moyennes et des pentes entre les deux sous-périodes (1960-1990 vs 1991-2020) pour ces deux indices. Que cela signifie-t-il pour le risque d'inondation à Douala ?
    \item Le test de Mann-Kendall sur la période totale donne $p < 0,05$ pour les deux indices. Expliquez ce que cela signifie statistiquement et pourquoi le test de Sen est préféré à une régression linéaire classique (moindres carrés) pour la magnitude.
    \item En vous basant sur la physique de l'atmosphère (relation Clausius-Clapeyron), expliquez le mécanisme thermodynamique liant le réchauffement global (+1,2°C depuis l'ère préindustrielle) à l'augmentation observée de RX1day et R95p sous les tropiques humides.
    \item Proposez une mesure d'adaptation urbaine concrète pour Douala dérivée de ce diagnostic (intensité + fréquence extrêmes en hausse).
\end{enumerate}
\tcblower
\textbf{Solution détaillée :}
\begin{enumerate}
    \item \textbf{Définitions :}
          \begin{itemize}
              \item \textbf{RX1day (Max 1-day precipitation amount) :} \textbf{Valeur maximale} de la précipitation journalière enregistrée sur l'année. Unité : mm. Représente l'\textbf{intensité maximale ponctuelle} (pic de crue potentielle).
              \item \textbf{R95p (Very wet days precipitation) :} \textbf{Somme des précipitations} tombées les jours où la pluviométrie journalière \textbf{dépasse le 95ème percentile} de la distribution des jours de pluie ($\ge$ 1 mm) calculé sur la période de référence (1981-2010). Unité : mm. Représente le \textbf{volume total apporté par les événements extrêmes} (charge hydrique extrême).
          \end{itemize}
          \textbf{Complémentarité :} RX1day = "Pic instantané" (dimensionnement ouvrage ponctuel : dalot, caniveau). R95p = "Volume extrême cumulé" (dimensionnement bassin rétention, saturation sol, risque inondation généralisée). R95p capture la \textbf{fréquence} des jours extrêmes, RX1day seulement le \textbf{max}.
    \item \textbf{Interprétation évolution :}
          \begin{itemize}
              \item \textbf{Moyennes en hausse :} RX1day passe de 85 $\rightarrow$ 99 mm (+16\%). R95p passe de 420 $\rightarrow$ 580 mm (+38\%). Les événements extrêmes sont \textbf{plus intenses} ET \textbf{plus fréquents/volumineux}.
              \item \textbf{Pentes en forte accélération :} Pente Sen RX1day $\times$ 3,7 (0,12 $\rightarrow$ 0,45 mm/an). Pente Sen R95p $\times$ 2,3 (1,8 $\rightarrow$ 4,2 mm/an). L'intensification s'\textbf{accélère} récemment.
              \item \textbf{Conséquence Risque Douala :} Douala (ville basse, mangrove, drainage déficient, urbanisation galopante) : \textbf{Fréquence inondations éclairs multipliée}. Seuils de saturation des réseaux d'assainissement atteints plus souvent. Crues de plus en plus hautes (RX1day) et longues (R95p $\uparrow$ = sols saturés plus longtemps). Risque "inondation majeure" (retour 20 ans) devient retour 5-10 ans.
          \end{itemize}
    \item \textbf{Signification tests statistiques :}
          \begin{itemize}
              \item \textbf{Mann-Kendall ($p < 0,05$) :} Rejet de l'hypothèse nulle $H_0$ "pas de tendance monotone". La tendance à la hausse est \textbf{statistiquement significative} (probabilité < 5\% qu'elle soit due au hasard/variabilité interne). Il y a un \textbf{signal climatique détectable}.
              \item \textbf{Sen vs Moindres Carrés (OLS) :} OLS sensible aux \textbf{valeurs aberrantes (outliers)}. Une seule année exceptionnelle (ex: 2020) fausse la pente OLS. L'estimateur de \textbf{Sen (Theil-Sen)} calcule la médiane de toutes les pentes entre paires de points $\rightarrow$ \textbf{Robuste} (point de rupture 29\%). Donne la \textbf{magnitude réelle} de la tendance centrale, non biaisée par les extrêmes de la série elle-même.
          \end{itemize}
    \item \textbf{Mécanisme Thermodynamique (Clausius-Clapeyron - C-C) :}
          \begin{itemize}
              \item \textbf{Loi C-C :} Pression de saturation vapeur d'eau $e_s(T)$ croît exponentiellement avec $T$. $\frac{1}{e_s}\frac{de_s}{dT} \approx \frac{L_v}{R_v T^2} \approx \textbf{7\% par °C}$ près de 25-30°C.
              \item \textbf{Application atmosphérique :} Atmosphère plus chaude ($\Delta T \approx +1,2°C$ global, +1,5 à +2°C local Douala/UHI) $\rightarrow$ \textbf{Capacité de rétention d'humidité accrue de $\sim$ 7-14\%}.
              \item \textbf{Disponibilité humidité :} Douala = zone de convergence (Mousson + ZCIT + Brises de mer), Océan chaud (TSM $\uparrow$). \textbf{Humidité relative quasi constante} (proche saturation). $\rightarrow$ \textbf{Humidité spécifique $q$ augmente $\sim$ 7\%/°C}.
              \item \textbf{Précipitation extrême :} $P_{extr} \approx -\omega \cdot q$ (convergence verticale $\times$ humidité). Si dynamique ($\omega$) inchangée ou légèrement accrue (instabilité $\uparrow$), $P_{extr}$ suit $q$ $\rightarrow$ \textbf{Augmentation $\sim$ 7\%/°C (Scaling C-C)}.
              \item \textbf{Observé vs Théorique :} Pente RX1day totale +0,28 mm/an sur 61 ans $\approx$ +17 mm $\approx$ +20\% sur 85 mm. Réchauffement local $\sim$ +1,5°C $\rightarrow$ Attendu C-C $\sim$ +10,5\%. Observé > C-C $\rightarrow$ \textbf{Super-CC scaling} possible (rétroaction dynamique : réchauffement latent $\rightarrow$ convergence plus forte $\rightarrow$ $\omega$ plus négatif). R95p augmente encore plus (+38\% volume) $\rightarrow$ \textbf{Fréquence jours extrêmes augmente} (seuil 95p fixe dépassé plus souvent car distribution se décale et s'aplatit).
          \end{itemize}
    \item \textbf{Mesure d'adaptation urbaine Douala :}
          \textbf{"Ville Éponge" (Sponge City) à l'échelle des bassins versants urbains (Bassa, Bonabéri, Akwa, Deido).}
          \begin{itemize}
              \item \textbf{Objectif :} Restaurer le cycle de l'eau naturel : \textbf{Infiltrer, Stocker, Évaporer, Réutiliser, Drainer (lentement)}.
              \item \textbf{Actions concrètes :}
                  \begin{enumerate}
                      \item \textbf{Désimperméabilisation massive :} Remplacement dallages/asphalte par \textbf{revêtements perméables} (béton drainant, pavés gazonnés, gravier stabilisé) parkings, cours écoles, trottoirs, places marchés.
                      \item \textbf{Infrastructures Vertes/Bleues :} Création \textbf{jardins de pluie (rain gardens)}, \textbf{noues végétalisées (swales)} le long voiries, \textbf{toitures végétalisées} obligatoires nouveaux bâtiments > 500 m², \textbf{bassins de rétention paysagers} multifonctions (parc sec $\rightarrow$ bassin crue).
                      \item \textbf{Restauration zones humides urbaines :} Protection/restauration \textbf{mangroves arrière-littorales} (Wouri, Dibamba) et \textbf{marais intérieurs} (Mboalé, Ngwè) comme \textbf{zones d'expansion de crues naturelles} (stockage gratuit, épuration, biodiversité).
                      \item \textbf{Gestion intégrée eaux pluviales (GIEP) :} Séparation réseaux eaux usées / eaux pluviales (nouveaux quartiers). Règlementation : \textbf{Coefficient de ruissellement maxi 0,3} pour toute nouvelle construction (vs 0,9 actuel). Incitation fiscale (taxe foncière réduite) pour existants.
                  \end{enumerate}
                  \textbf{Bénéfice co-bénéfices :} Réduction pic crue (RX1day impact), recharge nappe (sécheresse), fraîcheur urbaine (îlot de chaleur), cadre de vie, pêche urbaine.
          \end{itemize}
\end{enumerate}
\end{exoresolu}

\begin{attention}[Les 5 erreurs qui coûtent des points au Probatoire -- Leçon ND12]
    \begin{enumerate}
        \item \textbf{Confondre "Cellule de Hadley" et "ZCIT" :} La ZCIT est la \textbf{zone de convergence au sol} (branche ascendante) de la cellule de Hadley. Elle \textbf{migre} (Nord/Sud). La cellule de Hadley est la \textbf{circulation méridienne complète} (montée équateur, altitude, subsidence 30°, alizés sol). \textbf{Piège :} Dire "La ZCIT monte à 30°" (Faux : c'est l'air qui monte \emph{à} la ZCIT, la ZCIT elle reste au sol). Dire "Les alizés soufflent vers les pôles" (Faux : ils soufflent vers l'équateur).
        \item \textbf{Inverser le sens de rotation des gyres / Courants de bord :} \textbf{Hémisphère Nord = Horaire. Hémisphère Sud = Anti-horaire.} \textbf{Bord Ouest = Chaud, Rapide, Étroit (Gulf Stream, Kuroshio, Brésil, Agulhas). Bord Est = Froid, Lent, Large (Canaries, Californie, Benguela, Humboldt).} \textbf{Erreur classique :} Placer le Gulf Stream à l'Est ou le courant des Canaries à l'Ouest. Oublier l'intensification occidentale (effet $\beta$).
        \item \textbf{Mélanger Cyclone Tropical et Ligne de Grain (Orage violent) au Cameroun :} \textbf{Il n'y a PAS de cyclone tropical au Cameroun (Golfe de Guinée).} Les vents violents + pluies diluviennes = \textbf{Lignes de Grain (Squall Lines) / MCS} propagées par le Jet d'Est Africain. \textbf{Erreur éliminatoire :} Expliquer une inondation à Douala ou Maroua par "le passage d'un cyclone". Citer "Cyclone" pour le Golfe de Guinée = 0 point à la question "Origine dynamique".
        \item \textbf{Confondre Prévention et Protection (Gestion du risque) :} \textbf{Prévention = Agir sur l'aléa ou l'exposition (Ne pas construire en zone inondable, Reboisement, Recul trait de côte). Protection = Agir sur la vulnérabilité face à l'aléa inévitable (Digue, Mur, Surélévation maison, Sirène).} \textbf{Consigne :} "Proposez une mesure de prévention" $\rightarrow$ Répondre "Construire une digue" = \textbf{Hors sujet (0 pt)}. "Proposez une mesure de protection" $\rightarrow$ Répondre "Interdire la construction" = \textbf{Hors sujet}.
        \item \textbf{Négliger le contexte camerounais (Données locales) dans les exemples :} Le programme exige des exemples du Cameroun. \textbf{Erreur :} Citer seulement Katrina (USA), Haiyan (Philippines), ou des théories abstraites sans ancrage local. \textbf{Réflexe gagnant :} Pour "Circulation atmosphérique" $\rightarrow$ ZCIT/Harmattan/Mousson sur le Cameroun. "Courants océaniques" $\rightarrow$ Courant de Guinée / Upwelling / Pêche Kribi/Limbé. "Catastrophes" $\rightarrow$ Inondations Logone 2012/2022, Douala 2023, Glissements Mont Cameroun/Buea, Sécheresse Extrême-Nord, Érosion Limbé/Kribi. "Prévention" $\rightarrow$ PNGRC, PCS, ONPC, ACMAD.
    \end{enumerate}
\end{attention}

\newpage

\coursec{Exercices}

\courssub{Je vérifie mes connaissances}

\begin{exercice}[QCM : la circulation atmosphérique et océanique]
Pour chaque question, choisir la (ou les) bonne(s) réponse(s).

\textbf{1.} La cellule de Hadley est une cellule de convection :
\begin{itemize}
\item[a)] située entre le pôle et la latitude $60^\circ$ ;
\item[b)] située entre l'équateur et la latitude $30^\circ$ environ ;
\item[c)] où l'air chaud monte à l'équateur et redescend vers les tropiques ;
\item[d)] responsable des vents d'ouest des moyennes latitudes.
\end{itemize}

\textbf{2.} Les alizés sont des vents :
\begin{itemize}
\item[a)] soufflant des tropiques vers l'équateur ;
\item[b)] soufflant de l'équateur vers les pôles ;
\item[c)] déviés par la force de Coriolis ;
\item[d)] qui n'existent que dans l'hémisphère Nord.
\end{itemize}

\textbf{3.} Le Gulf Stream est :
\begin{itemize}
\item[a)] un courant froid longeant les côtes de l'Angola ;
\item[b)] un courant chaud de l'Atlantique Nord ;
\item[c)] un courant qui adoucit le climat de l'Europe occidentale ;
\item[d)] un courant de surface uniquement.
\end{itemize}

\textbf{4.} Un cyclone tropical se forme obligatoirement :
\begin{itemize}
\item[a)] sur une mer dont la température dépasse environ $26,5\,^\circ$C ;
\item[b)] exactement à l'équateur ;
\item[c)] à une latitude suffisante pour que la force de Coriolis soit efficace ;
\item[d)] au-dessus des continents.
\end{itemize}
\end{exercice}

\begin{exercice}[Vrai ou faux, en justifiant]
Répondre par \textbf{vrai} ou \textbf{faux}, puis justifier chaque réponse en une ou deux phrases.

\textbf{1.} L'air chaud étant plus dense que l'air froid, il descend au niveau de l'équateur.

\textbf{2.} La force de Coriolis dévie les vents vers la droite dans l'hémisphère Nord et vers la gauche dans l'hémisphère Sud.

\textbf{3.} Le courant de Benguela est un courant chaud qui réchauffe les côtes de la Namibie et de l'Angola.

\textbf{4.} Les inondations ne peuvent avoir que des causes naturelles.
\end{exercice}

\begin{exercice}[Définitions]
Définir en une phrase claire et scientifiquement correcte chacun des termes suivants :
\begin{enumerate}
\item cellule de convection ;
\item alizé ;
\item courant marin ;
\item cyclone tropical ;
\item El Niño.
\end{enumerate}
\end{exercice}

\begin{exercice}[Calculs directs : énergie et vitesse]
\textbf{1.} Un cyclone tropical de catégorie 3 présente des vents de \SI{180}{km.h^{-1}}. Convertir cette vitesse en mètres par seconde (arrondir à l'unité).

\textbf{2.} La pression au centre (l'\oe il) d'un cyclone est de \SI{950}{hPa}, alors que la pression normale au niveau de la mer est de \SI{1013}{hPa}. Calculer l'écart de pression en hectopascals, puis exprimer cet écart en pourcentage de la pression normale (arrondir à $0,1\,\%$).

\textbf{3.} Lors d'une crue à Douala, il est tombé \SI{120}{mm} de pluie en 6 heures. Calculer l'intensité moyenne de cette pluie en \si{mm.h^{-1}}. Sachant qu'une pluie est dite diluvienne au-delà de \SI{15}{mm.h^{-1}}, conclure.
\end{exercice}

\courssub{Je m'entraîne}

\begin{exercice}[Schéma muet de la circulation générale]
On donne un schéma muet d'un hémisphère terrestre, de l'équateur ($0^\circ$) au pôle ($90^\circ$), avec les latitudes $30^\circ$ et $60^\circ$ repérées.
\begin{enumerate}
\item Reproduire le schéma et y placer les trois cellules de convection (cellule de Hadley, cellule de Ferrel, cellule polaire) en indiquant par des flèches le sens de déplacement de l'air.
\item Placer les zones de hautes pressions (H) et de basses pressions (B) aux latitudes correspondantes, en justifiant chaque position par le mouvement vertical de l'air.
\item Dessiner le sens des alizés entre $30^\circ$ et l'équateur, puis expliquer pourquoi ces vents ne soufflent pas directement du nord vers le sud mais sont déviés.
\end{enumerate}
\end{exercice}

\begin{exercice}[Carte des courants de l'Atlantique]
Sur une carte muette de l'océan Atlantique figurent quatre flèches numérotées de 1 à 4, correspondant aux courants suivants : Gulf Stream, courant des Canaries, courant de Benguela, courant de Guinée.
\begin{enumerate}
\item Identifier chaque courant en justifiant votre choix par sa position géographique.
\item Préciser la nature (chaude ou froide) de chaque courant.
\item Expliquer l'influence du courant de Guinée sur le climat du littoral camerounais (températures, humidité, précipitations).
\item Citer une conséquence climatique du courant de Benguela sur la côte sud-ouest de l'Afrique.
\end{enumerate}
\end{exercice}

\begin{exercice}[Pourquoi pas de cyclones dans le golfe de Guinée ?]
Chaque année, les Caraïbes et le golfe du Mexique sont frappés par des cyclones tropicaux, alors qu'aucun cyclone n'a jamais été observé dans le golfe de Guinée, pourtant bordé par des eaux chaudes.
\begin{enumerate}
\item Rappeler les trois conditions nécessaires à la formation d'un cyclone tropical.
\item Montrer, en raisonnant sur la latitude du golfe de Guinée (entre $0^\circ$ et $5^\circ$N environ) et sur la force de Coriolis, pourquoi l'une de ces conditions n'y est jamais remplie.
\item En déduire pourquoi les côtes camerounaises ne craignent pas les cyclones, mais restent exposées à d'autres phénomènes dangereux que vous citerez.
\end{enumerate}
\end{exercice}

\begin{exercice}[Chaîne causale : El Niño]
El Niño est un phénomène qui se manifeste par un réchauffement anormal des eaux de surface du Pacifique équatorial oriental, au large du Pérou.
\begin{enumerate}
\item Décrire brièvement la situation normale (eaux, alizés, courant de Humboldt) dans le Pacifique équatorial.
\item Construire une chaîne causale complète, sous la forme : modification océanique $\rightarrow$ conséquence atmosphérique $\rightarrow$ phénomène climatique $\rightarrow$ dégâts, pour deux régions différentes (côte ouest de l'Amérique du Sud ; Indonésie/Australie).
\item Expliquer pourquoi El Niño illustre bien la notion d'interaction océan-atmosphère.
\end{enumerate}
\end{exercice}

\courssub{Je me prépare au Probatoire}

\begin{exercice}[Sujet type Probatoire : les cyclones tropicaux]
\textbf{Document 1.} Le cyclone est un système tourbillonnaire qui naît au-dessus des océans tropicaux. Le tableau ci-dessous donne les caractéristiques de deux cyclones observés dans l'Atlantique Nord.

\begin{center}
\begin{tabularx}{\linewidth}{|l|X|X|}
\hline
\rowcolor{popblueL} \textbf{Caractéristique} & \textbf{Cyclone A} & \textbf{Cyclone B} \\
\hline
Température des eaux de surface & $29\,^\circ$C & $24\,^\circ$C \\
\hline
Latitude de formation & $15^\circ$N & $2^\circ$N \\
\hline
Pression au centre & \SI{945}{hPa} & --- \\
\hline
Vitesse maximale des vents & \SI{210}{km.h^{-1}} & --- \\
\hline
\end{tabularx}
\end{center}

\textbf{Document 2.} On donne un schéma en coupe verticale d'un cyclone mûr, sur lequel figurent en lettres muettes : l'\oe il, le mur de l'\oe il, les bandes nuageuses spiralées, le sens de rotation des vents et les mouvements verticaux de l'air.

\begin{enumerate}
\item À l'aide de vos connaissances, rappeler les conditions de formation d'un cyclone tropical et vérifier, pour chacun des cyclones A et B du document 1, si ces conditions sont réunies. Conclure sur le sort probable du système B. \emph{(4 points)}
\item Reproduire le schéma du document 2 et le compléter en légendant les cinq éléments muets. \emph{(3 points)}
\item Expliquer, à partir des mouvements de l'air, pourquoi le ciel est calme et dégagé dans l'\oe il alors que les vents et les pluies sont violents dans le mur de l'\oe il. \emph{(3 points)}
\item Convertir la vitesse maximale des vents du cyclone A en \si{m.s^{-1}}, puis citer trois types de dégâts qu'un tel cyclone peut provoquer à son arrivée sur les côtes. \emph{(3 points)}
\item Le Cameroun n'est jamais touché par les cyclones : proposer une explication scientifique à cette particularité. \emph{(2 points)}
\end{enumerate}
\end{exercice}

\begin{exercice}[Sujet type Probatoire : inondations à Yaoundé]
\textbf{Document 1.} Extrait d'un article de presse : « Pluies diluviennes à Yaoundé : plusieurs quartiers inondés. Les eaux ont envahi les habitations construites le long des marécages et sur les berges des ruisseaux. Les caniveaux, obstrués par des déchets plastiques, n'ont pas pu évacuer les eaux de ruissellement. Les pentes déboisées des sept collines de la ville ont favorisé l'écoulement rapide des eaux vers les bas-fonds. On déplore des pertes en vies humaines et d'importants dégâts matériels. »

\textbf{Document 2.} Rappel : Yaoundé connaît une saison des pluies marquée ; la ville est bâtie sur un relief accidenté de collines séparées par des vallées marécageuses ; sa population croît rapidement, avec un habitat spontané dans les zones non aedificandi (zones inconstructibles).

\begin{enumerate}
\item Dégager du document 1 les causes \textbf{naturelles} et les causes \textbf{humaines} de cette catastrophe, en les classant dans un tableau à deux colonnes. \emph{(4 points)}
\item Expliquer, en vous appuyant sur vos connaissances sur le ruissellement et l'infiltration, comment le déboisement des pentes aggrave les inondations en ville. \emph{(3 points)}
\item Relier les pluies diluviennes de Yaoundé aux mouvements atmosphériques (convection, zones de basses pressions, mousson) qui caractérisent le climat de la région. \emph{(3 points)}
\item Proposer \textbf{cinq mesures concrètes} de prévention et de protection adaptées au contexte de Yaoundé, en précisant pour chacune si elle relève de l'aménagement du territoire, de l'information des populations ou des secours. \emph{(5 points)}
\end{enumerate}
\end{exercice}

\begin{exercice}[Sujet type Probatoire : circulation océanique et climat du littoral camerounais]
\textbf{Document 1.} Sur une carte de l'Atlantique tropical figurent : le courant de Guinée (courant chaud s'écoulant d'ouest en est le long des côtes du golfe de Guinée), le courant de Benguela (courant froid remontant le long des côtes de l'Angola et de la Namibie) et les alizés de sud-est.

\textbf{Document 2.} Données climatiques moyennes annuelles de deux villes côtières situées à des latitudes voisines :

\begin{center}
\begin{tabularx}{\linewidth}{|l|X|X|}
\hline
\rowcolor{popblueL} \textbf{Ville} & \textbf{Température moyenne} & \textbf{Précipitations annuelles} \\
\hline
Douala (Cameroun, $4^\circ$N) & $26,5\,^\circ$C & environ \SI{4000}{mm} \\
\hline
Namibe (Angola, $15^\circ$S) & $21\,^\circ$C & moins de \SI{50}{mm} \\
\hline
\end{tabularx}
\end{center}

\begin{enumerate}
\item Décrire, à partir du document 1, le trajet et la nature du courant de Guinée et du courant de Benguela. \emph{(3 points)}
\item Expliquer comment le courant de Guinée, associé à la mousson, contribue aux fortes précipitations enregistrées à Douala. \emph{(4 points)}
\item Expliquer, en vous appuyant sur le document 2, pourquoi la côte angolaise, bordée par le courant de Benguela, est au contraire aride (désert côtier du Namibe). \emph{(4 points)}
\item Les pluies abondantes du littoral camerounais provoquent régulièrement des inondations et des glissements de terrain (région de Limbe, Buea). Citer deux facteurs aggravants liés aux activités humaines et proposer trois mesures de protection des populations. \emph{(4 points)}
\end{enumerate}
\end{exercice}

\newpage

\coursec{Corrigés des exercices}

\coursec{Exercices corrigés et sujets type Probatoire — Leçon ND12}

\begin{corrige}[Exercice 1 — QCM : la circulation atmosphérique et océanique]
\textbf{1. Réponses b et c.} La cellule de \cle{Hadley} est la cellule de convection tropicale, située entre l'équateur et environ $30^\circ$ de latitude : l'air chaud et humide s'élève à l'équateur (zone de basses pressions), se déplace en altitude vers les tropiques, redescend vers $30^\circ$ (zone de hautes pressions), puis revient en surface vers l'équateur. La réponse a décrit la cellule polaire, et la d est fausse : les vents d'ouest des moyennes latitudes sont liés à la cellule de Ferrel.

\textbf{2. Réponses a et c.} Les \cle{alizés} sont des vents de surface qui soufflent des hautes pressions subtropicales (vers $30^\circ$) vers les basses pressions équatoriales. Ils sont déviés par la \cle{force de Coriolis} : vers la droite dans l'hémisphère Nord (alizés de nord-est), vers la gauche dans l'hémisphère Sud (alizés de sud-est). Ils existent donc dans les deux hémisphères (d est faux).

\textbf{3. Réponses b et c.} Le \cle{Gulf Stream} est un courant marin chaud de l'Atlantique Nord, qui transporte des eaux tropicales vers les côtes européennes et adoucit le climat de l'Europe occidentale. Le courant froid longeant l'Angola est le courant de Benguela (a est faux). Le Gulf Stream n'est pas qu'un courant de surface : il s'inscrit dans la circulation thermohaline globale (d est faux).

\textbf{4. Réponses a et c.} Un cyclone tropical exige une mer à plus de $26,5\,^\circ$C environ (source d'énergie : chaleur latente de condensation) et une latitude suffisante (en général au-delà de $5^\circ$) pour que la force de Coriolis mette la perturbation en rotation. À l'équateur, la force de Coriolis est nulle (b est faux), et le cyclone se forme au-dessus des océans, jamais des continents (d est faux).
\end{corrige}

\begin{corrige}[Exercice 2 — Vrai ou faux, en justifiant]
\textbf{1. Faux.} L'air chaud est \emph{moins} dense que l'air froid (il est dilaté) : il s'élève donc au niveau de l'équateur, créant une zone de basses pressions (la zone de convergence intertropicale, ZCIT).

\textbf{2. Vrai.} La force de Coriolis, due à la rotation de la Terre, dévie tout mouvement (vents, courants) vers la droite dans l'hémisphère Nord et vers la gauche dans l'hémisphère Sud. C'est elle qui donne aux alizés leur orientation nord-est/sud-est et aux cyclones leur sens de rotation.

\textbf{3. Faux.} Le courant de \cle{Benguela} est un courant \emph{froid} qui remonte le long des côtes de la Namibie et de l'Angola. En refroidissant l'air de surface, il stabilise l'atmosphère et empêche la convection : il est responsable de l'aridité du désert côtier du Namibe.

\textbf{4. Faux.} Les inondations ont souvent des causes humaines aggravantes ou déterminantes : déboisement des pentes (accélération du ruissellement), urbanisation des zones inondables, imperméabilisation des sols, obstruction des caniveaux par les déchets. La catastrophe résulte presque toujours de la combinaison d'un aléa naturel et d'une vulnérabilité humaine.
\end{corrige}

\begin{corrige}[Exercice 3 — Définitions]
\begin{enumerate}
\item \textbf{Cellule de convection} : circuit fermé de mouvements de l'air dans la troposphère, dans lequel l'air chaud s'élève dans une zone, se déplace en altitude, redescend dans une autre zone et revient en surface, sous l'effet des différences de température et de pression.
\item \textbf{Alizé} : vent régulier de surface soufflant des hautes pressions subtropicales (vers $30^\circ$ de latitude) vers les basses pressions équatoriales, dévié par la force de Coriolis (de nord-est dans l'hémisphère Nord, de sud-est dans l'hémisphère Sud).
\item \textbf{Courant marin} : déplacement permanent et dirigé d'une masse d'eau de mer, en surface ou en profondeur, provoqué par les vents dominants, la force de Coriolis et les différences de température et de salinité.
\item \textbf{Cyclone tropical} : système tourbillonnaire de très basses pressions se formant au-dessus des océans tropicaux chauds (plus de $26,5\,^\circ$C), caractérisé par des vents violents en rotation autour d'un \oe il central et des pluies diluviennes.
\item \textbf{El Niño} : anomalie climatique périodique (tous les 2 à 7 ans) caractérisée par un réchauffement anormal des eaux de surface du Pacifique équatorial oriental, lié à un affaiblissement des alizés, et qui bouleverse les précipitations à l'échelle du globe.
\end{enumerate}
\end{corrige}

\begin{corrige}[Exercice 4 — Calculs directs : énergie et vitesse]
\textbf{1.} Conversion de \SI{180}{km.h^{-1}} en \si{m.s^{-1}} :
\[ v = \frac{180 \times 1000}{3600} = \frac{180}{3,6} = \SI{50}{m.s^{-1}}. \]

\textbf{2.} Écart de pression :
\[ \Delta P = 1013 - 950 = \SI{63}{hPa}. \]
En pourcentage de la pression normale :
\[ \frac{63}{1013} \times 100 \approx 6,2\,\%. \]
L'écart est de \SI{63}{hPa}, soit environ $6,2\,\%$ de la pression normale : une dépression très creuse, signe d'un cyclone puissant.

\textbf{3.} Intensité moyenne de la pluie :
\[ I = \frac{120\ \text{mm}}{6\ \text{h}} = \SI{20}{mm.h^{-1}}. \]
Conclusion : $20 > 15$, donc cette pluie est \textbf{diluvienne} ; à Douala, ville basse et marécageuse, une telle intensité déborde rapidement les réseaux d'évacuation et provoque des inondations.
\end{corrige}

\begin{attention}[Points de vigilance — Exercice 4]
Pour convertir des \si{km.h^{-1}} en \si{m.s^{-1}}, on divise par $3,6$ (et non l'inverse). Pour le pourcentage, on divise bien l'écart par la valeur de référence ($1013$ hPa), pas par la valeur mesurée.
\end{attention}

\begin{corrige}[Exercice 5 — Schéma muet de la circulation générale]
\textbf{1.} De l'équateur au pôle, on place trois cellules :
\begin{itemize}
\item \textbf{Cellule de Hadley} ($0^\circ$--$30^\circ$) : air ascendant à l'équateur, flux d'altitude vers le pôle, air descendant vers $30^\circ$, flux de surface vers l'équateur ;
\item \textbf{Cellule de Ferrel} ($30^\circ$--$60^\circ$) : cellule inverse (indirecte) : air descendant vers $30^\circ$, flux de surface vers $60^\circ$, air ascendant vers $60^\circ$, retour en altitude vers $30^\circ$ ;
\item \textbf{Cellule polaire} ($60^\circ$--$90^\circ$) : air descendant au pôle, flux de surface vers $60^\circ$, air ascendant vers $60^\circ$, retour en altitude vers le pôle.
\end{itemize}

\textbf{2.} Zones de pression :
\begin{itemize}
\item \textbf{B à l'équateur} ($0^\circ$) : l'air chaud s'élève, la pression en surface diminue (basses pressions équatoriales) ;
\item \textbf{H vers $30^\circ$} : l'air redescend et s'accumule en surface (hautes pressions subtropicales) ;
\item \textbf{B vers $60^\circ$} : convergence des flux de surface et ascendance de l'air (basses pressions subpolaires) ;
\item \textbf{H au pôle} ($90^\circ$) : l'air froid et dense descend (hautes pressions polaires).
\end{itemize}

\textbf{3.} Entre $30^\circ$ et l'équateur, le vent de surface souffle des H vers les B, donc vers l'équateur. Mais la \cle{force de Coriolis}, due à la rotation terrestre, le dévie vers la droite dans l'hémisphère Nord (alizé de nord-est) et vers la gauche dans l'hémisphère Sud (alizé de sud-est) : les alizés soufflent donc obliquement, du nord-est ou du sud-est, et non directement du pôle vers l'équateur.
\end{corrige}

\begin{corrige}[Exercice 6 — Carte des courants de l'Atlantique]
\textbf{1. et 2.} Identification et nature des courants :
\begin{center}
\begin{tabularx}{\linewidth}{|l|X|X|}
\hline
\rowcolor{popblueL} \textbf{Courant} & \textbf{Position géographique} & \textbf{Nature} \\
\hline
Gulf Stream & Longe la côte est de l'Amérique du Nord puis traverse l'Atlantique Nord vers l'Europe & Chaud \\
\hline
Courant des Canaries & Descend le long des côtes du Maroc et des îles Canaries & Froid \\
\hline
Courant de Benguela & Remonte le long des côtes de la Namibie et de l'Angola (Atlantique Sud-Est) & Froid \\
\hline
Courant de Guinée & S'écoule d'ouest en est le long des côtes du golfe de Guinée & Chaud \\
\hline
\end{tabularx}
\end{center}

\textbf{3.} Le courant de Guinée réchauffe les eaux du golfe de Guinée ; l'air de la mousson qui passe au-dessus se charge en chaleur et en vapeur d'eau. En arrivant sur le littoral camerounais, cet air chaud et humide s'élève (convection, relief du Mont Cameroun), la vapeur se condense : températures élevées et régulières, forte humidité et précipitations abondantes (environ \SI{4000}{mm} par an à Douala, plus de \SI{10000}{mm} à Debundscha).

\textbf{4.} Le courant froid de Benguela refroidit l'air de surface, ce qui stabilise l'atmosphère et empêche l'ascendance de l'air humide : la côte sud-ouest africaine connaît une aridité extrême (désert du Namibe), malgré la proximité de l'océan.
\end{corrige}

\begin{corrige}[Exercice 7 — Pourquoi pas de cyclones dans le golfe de Guinée ?]
\textbf{1.} Trois conditions nécessaires à la formation d'un cyclone tropical :
\begin{itemize}
\item une mer dont la température de surface dépasse environ $26,5\,^\circ$C sur une épaisseur suffisante (source d'énergie : évaporation puis condensation libérant de la chaleur latente) ;
\item une latitude suffisante (en général au-delà de $5^\circ$) pour que la \cle{force de Coriolis} mette la perturbation en rotation ;
\item une faible variation du vent avec l'altitude (faible cisaillement vertical), pour que le système ne soit pas déchiré.
\end{itemize}

\textbf{2.} Le golfe de Guinée s'étend entre $0^\circ$ et $5^\circ$N environ. Or la force de Coriolis est \emph{nulle à l'équateur} et très faible à proximité : elle est proportionnelle au sinus de la latitude, et $\sin 0^\circ = 0$. Une perturbation qui s'y formerait ne pourrait donc pas être mise en rotation : la deuxième condition n'y est jamais remplie, même si les eaux y sont bien plus chaudes que $26,5\,^\circ$C.

\textbf{3.} Les côtes camerounaises sont donc naturellement protégées des cyclones par leur position quasi équatoriale. En revanche, elles restent exposées à d'autres phénomènes dangereux : les \textbf{inondations} (pluies diluviennes de la mousson, marées, montée des eaux), les \textbf{glissements de terrain} sur les pentes du Mont Cameroun (région de Limbe et Buea), l'\textbf{érosion côtière} et les fortes houles, ainsi que les orages violents accompagnés de vents destructeurs (lignes de grains).
\end{corrige}

\begin{corrige}[Exercice 8 — Chaîne causale : El Niño]
\textbf{1.} Situation normale dans le Pacifique équatorial : les alizés de sud-est poussent les eaux chaudes de surface vers l'ouest (Indonésie, Australie), où la mer est chaude et la convection active (pluies abondantes). À l'est, le long du Pérou, le courant froid de Humboldt et la remontée d'eaux profondes (upwelling) maintiennent des eaux froides, un temps sec et des pêches riches (nutriments remontés du fond).

\textbf{2.} Chaînes causales lors d'un épisode El Niño :

\textbf{Côte ouest de l'Amérique du Sud (Pérou, Équateur) :} affaiblissement des alizés $\rightarrow$ eaux chaudes de surface qui reviennent vers l'est et réchauffement anormal du Pacifique oriental $\rightarrow$ évaporation et convection accrues au-dessus des eaux chaudes $\rightarrow$ pluies diluviennes sur des côtes normalement arides $\rightarrow$ inondations, glissements de terrain, destructions ; effondrement de la pêche (disparition de l'upwelling et des nutriments).

\textbf{Indonésie / Australie :} déplacement des eaux chaudes vers l'est $\rightarrow$ refroidissement relatif des eaux de l'ouest du Pacifique et affaiblissement de la convection $\rightarrow$ sécheresse exceptionnelle $\rightarrow$ mauvaises récoltes, feux de forêt gigantesques, pénuries d'eau.

\textbf{3.} El Niño illustre parfaitement l'interaction océan-atmosphère : une modification de la température de surface de l'océan (réchauffement) modifie l'évaporation, les pressions et les vents (atmosphère), ce qui renforce à son tour le déplacement des eaux chaudes (océan). Océan et atmosphère fonctionnent comme un système couplé : un changement dans l'un entraîne des réponses dans l'autre, avec des conséquences climatiques à l'échelle planétaire.
\end{corrige}

\begin{corrige}[Exercice 9 — Sujet type Probatoire : les cyclones tropicaux]
\textbf{1. Conditions de formation (4 pts).} Un cyclone tropical exige : (a) des eaux de surface à plus de $26,5\,^\circ$C ; (b) une latitude suffisante (au-delà d'environ $5^\circ$) pour que la force de Coriolis soit efficace ; (c) un faible cisaillement vertical des vents.
\begin{itemize}
\item \textbf{Cyclone A} : eaux à $29\,^\circ$C ($> 26,5\,^\circ$C) et formation à $15^\circ$N (Coriolis efficace) : les conditions sont réunies, ce qui est confirmé par la dépression très creuse (\SI{945}{hPa}) et les vents violents (\SI{210}{km.h^{-1}}).
\item \textbf{Système B} : eaux à $24\,^\circ$C ($< 26,5\,^\circ$C) et formation à $2^\circ$N (Coriolis quasi nulle) : \emph{aucune} des deux conditions essentielles n'est remplie.
\end{itemize}
Conclusion : le système B ne peut pas se développer en cyclone ; il restera au stade de simple perturbation tropicale ou se dissipera.

\textbf{2. Schéma (3 pts).} Légendes attendues : l'\textbf{\oe il} (zone centrale calme, air descendant) ; le \textbf{mur de l'\oe il} (paroi nuageuse verticale entourant l'\oe il, air ascendant violent) ; les \textbf{bandes nuageuses spiralées} (enroulées autour du centre) ; le \textbf{sens de rotation des vents} (antihoraire dans l'hémisphère Nord) ; les \textbf{mouvements verticaux} (ascendance dans le mur, descente dans l'\oe il).

\textbf{3. (3 pts).} Dans le mur de l'\oe il, l'air chaud et humide s'élève très rapidement : la condensation massive libère de la chaleur latente qui entretient l'ascendance, d'où des vents extrêmes et des pluies diluviennes. Dans l'\oe il, au contraire, l'air \emph{descend} ; en descendant, il se réchauffe et se dessèche, les nuages s'évaporent : le ciel y est dégagé, les vents faibles et le calme règne.

\textbf{4. (3 pts).} Conversion : $v = \dfrac{210}{3,6} \approx \SI{58}{m.s^{-1}}$ (valeur exacte $58,3$). Dégâts possibles : destruction des habitations et des infrastructures par les vents violents ; inondations côtières par onde de tempête (surélévation brutale du niveau de la mer) ; inondations intérieures et glissements de terrain dus aux pluies diluviennes (aussi : rupture des réseaux électriques, pertes en vies humaines, contamination de l'eau potable).

\textbf{5. (2 pts).} Le Cameroun borde le golfe de Guinée, situé entre $0^\circ$ et $5^\circ$N environ : à ces latitudes quasi équatoriales, la force de Coriolis est nulle ou trop faible pour mettre une perturbation en rotation. Les cyclones ne peuvent donc jamais s'y former ni s'y maintenir.

\textbf{Barème indicatif} : 1. conditions 1,5 pt + analyse A 1 pt + analyse B 1 pt + conclusion 0,5 pt ; 2. 0,5 pt par légende correcte (5 légendes + rotation) ; 3. ascendance/condensation 1,5 pt + descente/assèchement 1,5 pt ; 4. conversion 1 pt + 3 dégâts 2 pts ; 5. Coriolis nulle à l'équateur 2 pts.
\end{corrige}

\begin{attention}[Points de vigilance — Exercice 9]
Citer explicitement le seuil de $26,5\,^\circ$C et le rôle de la force de Coriolis : ce sont les deux justifications attendues. Ne pas écrire que le cyclone B « ira ailleurs » : sans énergie (eau froide) et sans rotation (Coriolis), il se dissipe. Dans la conversion, diviser par $3,6$.
\end{attention}

\begin{corrige}[Exercice 10 — Sujet type Probatoire : inondations à Yaoundé]
\textbf{1. (4 pts).} Classement des causes :
\begin{center}
\begin{tabularx}{\linewidth}{|X|X|}
\hline
\rowcolor{popblueL} \textbf{Causes naturelles} & \textbf{Causes humaines} \\
\hline
Pluies diluviennes (aléa météorologique) & Habitats construits dans les marécages et sur les berges (zones inondables) \\
\hline
Relief accidenté : pentes des sept collines accélérant l'écoulement vers les bas-fonds & Caniveaux obstrués par les déchets plastiques \\
\hline
Vallées marécageuses naturellement inondables & Déboisement des pentes (suppression de la protection végétale) \\
\hline
 & Urbanisation rapide et non planifiée (habitat spontané) \\
\hline
\end{tabularx}
\end{center}

\textbf{2. (3 pts).} Sur une pente boisée, la végétation intercepte une partie de la pluie, les racines maintiennent le sol et favorisent l'\cle{infiltration} de l'eau : le ruissellement est lent et limité. Après déboisement, l'eau frappe directement le sol, l'infiltration diminue fortement et le \cle{ruissellement} s'accélère sur la pente : un volume d'eau considérable arrive brutalement dans les bas-fonds urbains, où les caniveaux (de surcroît obstrués) ne peuvent l'évacuer. Le déboisement transforme donc une pluie forte en crue éclair.

\textbf{3. (3 pts).} Yaoundé, proche de l'équateur, est sous l'influence de la zone de convergence intertropicale (ZCIT), zone de \textbf{basses pressions} où convergent les alizés et où l'air chaud et humide s'élève en permanence (\textbf{convection}). Pendant la saison des pluies, la \textbf{mousson} (flux de sud-ouest chargé d'humidité venant de l'océan Atlantique et du golfe de Guinée) alimente cette convection : l'air humide s'élève, se refroidit, la vapeur d'eau se condense en nuages convectifs (cumulonimbus) qui déversent des pluies intenses et prolongées, parfois diluviennes.

\textbf{4. (5 pts).} Cinq mesures (une par catégorie au minimum) :
\begin{itemize}
\item \textbf{Aménagement} : interdire et déguerpir les constructions dans les zones non aedificandi (marécages, berges) et reboiser les pentes des collines ;
\item \textbf{Aménagement} : curer, élargir et entretenir régulièrement les caniveaux et les ouvrages d'évacuation des eaux ;
\item \textbf{Information} : campagnes de sensibilisation contre le jet des déchets dans les caniveaux et sur les risques d'habitat en zone inondable ;
\item \textbf{Information} : mettre en place un système d'alerte météorologique (radiodiffusion des avis de fortes pluies, messages d'alerte) ;
\item \textbf{Secours} : organiser des plans d'évacuation et des centres d'accueil, avec des équipes de secours (protection civile, Croix-Rouge) pré-positionnées pendant la saison des pluies.
\end{itemize}

\textbf{Barème indicatif} : 1. 0,5 pt par cause correctement classée (8 attendues) ; 2. infiltration 1 pt + ruissellement accéléré 1 pt + lien avec la crue 1 pt ; 3. convection/basses pressions 1,5 pt + mousson humide 1,5 pt ; 4. 1 pt par mesure concrète et correctement catégorisée.
\end{corrige}

\begin{attention}[Points de vigilance — Exercice 10]
Distinguer clairement l'\emph{aléa} (la pluie, naturelle) de la \emph{vulnérabilité} (habitat, déchets, déboisement, humaine) : la catastrophe naît de leur rencontre. Les mesures doivent être \emph{concrètes} et \emph{adaptées au contexte de Yaoundé}, et chacune doit être rattachée à sa catégorie, comme demandé.
\end{attention}

\begin{corrige}[Exercice 11 — Sujet type Probatoire : circulation océanique et climat du littoral camerounais]
\textbf{1. (3 pts).} Le \textbf{courant de Guinée} est un courant \emph{chaud} qui s'écoule d'ouest en est le long des côtes du golfe de Guinée, jusqu'au littoral camerounais. Le \textbf{courant de Benguela} est un courant \emph{froid} qui remonte du sud vers le nord le long des côtes de la Namibie et de l'Angola, alimenté par les alizés de sud-est qui chassent les eaux de surface (avec upwelling d'eaux profondes froides).

\textbf{2. (4 pts).} Au-dessus du courant chaud de Guinée, l'évaporation est intense : l'air de la \textbf{mousson} (flux maritime de sud-ouest) se charge en chaleur et en vapeur d'eau. En atteignant le littoral camerounais, cet air chaud et humide est forcé de s'élever (convection liée aux basses pressions équatoriales, accentuée par le relief du Mont Cameroun) ; en altitude, la vapeur d'eau se condense, formant des nuages convectifs qui déversent des précipitations abondantes : environ \SI{4000}{mm} par an à Douala, et plus de \SI{10000}{mm} sur le versant exposé du Mont Cameroun (Debundscha, l'un des lieux les plus arrosés d'Afrique).

\textbf{3. (4 pts).} À Namibe ($21\,^\circ$C, moins de \SI{50}{mm} de pluie), le courant froid de Benguela refroidit l'air de surface. Cet air froid, plus dense, reste plaqué au sol : l'atmosphère est \emph{stable}, la convection est bloquée, et la faible humidité de l'air ne peut se condenser en nuages pluvieux. Il en résulte une aridité quasi totale malgré la proximité de l'océan : c'est le désert côtier du Namibe. La comparaison avec Douala montre que c'est la \emph{nature du courant} (chaud ou froid) qui commande le climat côtier, à latitude comparable.

\textbf{4. (4 pts).} Facteurs aggravants humains (deux attendus) : le \textbf{déboisement des pentes} du Mont Cameroun (agriculture, bois de chauffe), qui accélère le ruissellement et déstabilise les sols ; les \textbf{constructions anarchiques} dans les bas-fonds, le long des ravins et sur les pentes raides (aussi : caniveaux obstrués par les déchets).

Mesures de protection (trois attendues) : \textbf{reboiser} les pentes et protéger les versants raides de toute culture ; \textbf{aménager} les ouvrages de drainage et interdire l'habitat dans les zones à risque (relocalisation des populations exposées, comme à Limbe après les glissements de terrain) ; mettre en place un système d'\textbf{alerte précoce} et d'information des populations pendant la saison des pluies, avec des plans d'évacuation et des secours organisés.

\textbf{Barème indicatif} : 1. trajet 1 pt + nature 1 pt par courant (3 pts) ; 2. évaporation 1 pt + mousson humide 1 pt + ascendance/relief 1 pt + condensation/pluies 1 pt ; 3. refroidissement 1 pt + stabilité/blocage de la convection 2 pts + aridité 1 pt ; 4. 0,5 pt par facteur (1 pt) + 1 pt par mesure (3 pts).
\end{corrige}

\begin{attention}[Points de vigilance — Exercice 11]
Le raisonnement doit toujours enchaîner : nature du courant $\rightarrow$ température de l'air $\rightarrow$ stabilité ou convection $\rightarrow$ précipitations. Ne pas se contenter d'affirmer « le courant chaud donne de la pluie » sans expliquer le mécanisme (évaporation puis condensation lors de l'ascendance). Pour la question 4, distinguer nettement facteurs aggravants et mesures de protection.
\end{attention}

\newpage

\coursec{Fiche bilan}

\begin{popfigure}
\begin{center}
\begin{tikzpicture}[
  mindmap,
  grow cyclic,
  every node/.style={concept, text=white, font=\small\bfseries},
  root concept/.append style={concept color=popdark, font=\bfseries, minimum size=3.2cm},
  level 1 concept/.append style={level distance=4.6cm, sibling angle=72, font=\footnotesize\bfseries},
  level 2 concept/.append style={level distance=2.6cm, font=\scriptsize},
  concept connection/.append style={line width=1.2pt}
]
\node[root concept]{Catastrophes liées aux mouvements atmosphériques et océaniques}
  child[concept color=popblue]{ node{Circulation atmosphérique}
    child[concept color=popblue!70]{ node{Cellules de convection (Hadley, Ferrel, polaire)} }
    child[concept color=popblue!70]{ node{Vents planétaires : alizés, ouest, est} }
  }
  child[concept color=popgreen]{ node{Circulation océanique}
    child[concept color=popgreen!70]{ node{Courants de surface (vents)} }
    child[concept color=popgreen!70]{ node{Circulation thermohaline} }
  }
  child[concept color=poporange]{ node{Interactions océan--atmosphère}
    child[concept color=poporange!70]{ node{Échanges chaleur et humidité} }
    child[concept color=poporange!70]{ node{El Ni\~no / La Ni\~na} }
  }
  child[concept color=popink]{ node{Phénomènes dangereux}
    child[concept color=popink!70]{ node{Cyclones, tempêtes} }
    child[concept color=popink!70]{ node{Inondations, sécheresses} }
  }
  child[concept color=popgold]{ node{Prévention et protection}
    child[concept color=popgold!70]{ node{Surveillance, alerte} }
    child[concept color=popgold!70]{ node{Aménagement, éducation} }
  };
\end{tikzpicture}
\end{center}
\legende{Carte mentale de la leçon : les cinq grands axes à maîtriser, de la circulation générale de l'atmosphère jusqu'aux mesures de prévention des catastrophes.}
\end{popfigure}

\begin{aretenir}[L'essentiel de la leçon : définitions clés]
\begin{itemize}
  \item \cle{Cellule de convection} : boucle fermée de circulation de l'air (air chaud qui monte à l'équateur, air froid qui descend aux zones subtropicales ou polaires). On en compte trois par hémisphère : cellule de \textbf{Hadley} (entre l'équateur et environ $30^{\circ}$ de latitude), cellule de \textbf{Ferrel} (entre $30^{\circ}$ et $60^{\circ}$), cellule \textbf{polaire} (entre $60^{\circ}$ et le pôle).
  \item \cle{Alizés} : vents réguliers des zones intertropicales, soufflant des hautes pressions subtropicales vers les basses pressions équatoriales, déviés par la force de Coriolis (vers la droite dans l'hémisphère Nord, vers la gauche dans l'hémisphère Sud).
  \item \cle{Courant marin} : déplacement d'ensemble des eaux océaniques ; les courants de surface sont provoqués par les vents dominants, la circulation \cle{thermohaline} par les différences de température et de salinité (donc de densité) entre les masses d'eau.
  \item \cle{Cyclone tropical} : tourbillon dépressionnaire violent (vents supérieurs à \SI{118}{km/h}) né sur une mer chaude ($T > \SI{26}{\celsius}$) ; appelé \emph{ouragan} dans l'Atlantique Nord, \emph{typhon} dans le Pacifique Nord-Ouest.
  \item \cle{El Ni\~no} : anomalie chaude des eaux superficielles du Pacifique équatorial Est qui bouleverse les précipitations à l'échelle planétaire (sécheresses ici, inondations là) ; son inverse est \cle{La Ni\~na} (anomalie froide).
\end{itemize}
\end{aretenir}

\begin{aretenir}[Mécanismes à connaître par c\oe ur]
\begin{center}
\begin{tabularx}{\linewidth}{|l|X|}
\hline
\rowcolor{popblueL} \textbf{Phénomène} & \textbf{Origine dynamique} \\
\hline
Vents planétaires & Gradient de pression (air chaud à l'équateur, air froid aux pôles) + déviation par la force de Coriolis (rotation terrestre) \\
\hline
Courants de surface & Entraînement par les vents dominants + force de Coriolis + position des continents (formation de grands tourbillons appelés gyres) \\
\hline
Circulation thermohaline & Eaux froides et salées, plus denses, plongent aux hautes latitudes (<< tapis roulant >> océanique) \\
\hline
Cyclone tropical & Mer chaude ($>\SI{26}{\celsius}$) $\rightarrow$ évaporation intense $\rightarrow$ air chaud humide qui monte $\rightarrow$ dépression $\rightarrow$ rotation amplifiée par la force de Coriolis \\
\hline
Inondation & Pluies diluviennes (mousson, cyclone) + relief + imperméabilisation des sols + déforestation \\
\hline
\end{tabularx}
\end{center}
\end{aretenir}

\begin{methode}[Méthodes express]
\begin{enumerate}
  \item \textbf{Décrire un mouvement sur un schéma ou une carte} : nommer la zone (latitude), le sens du mouvement (horizontal ou vertical), la cause (pression, température, vent), la conséquence climatique (pluie, sécheresse).
  \item \textbf{Relier un danger à son origine} : toujours remonter la chaîne \emph{cause dynamique} $\rightarrow$ \emph{mécanisme} $\rightarrow$ \emph{phénomène observé} $\rightarrow$ \emph{dégâts}.
  \item \textbf{Proposer des protections} : penser aux trois volets \emph{prévision} (météorologie, systèmes d'alerte), \emph{aménagement} (digues, reboisement, urbanisme raisonné), \emph{éducation} (plans d'évacuation, comportements à adopter).
\end{enumerate}
\end{methode}

\begin{attention}[Pièges fréquents au Probatoire]
\begin{itemize}
  \item Ne pas confondre \textbf{courant de surface} (causé par les vents) et \textbf{circulation thermohaline} (causée par les différences de densité).
  \item La force de Coriolis \textbf{dévie} les vents et les courants, elle ne les \textbf{crée} pas : le moteur initial est le gradient de pression (pour l'air) ou l'action du vent (pour l'eau de surface).
  \item Un cyclone ne se forme jamais exactement à l'équateur : la force de Coriolis y est nulle, la rotation ne peut pas s'enclencher.
  \item Ne pas écrire << ouragan >> pour un phénomène de l'océan Indien : employer le terme général \textbf{cyclone tropical}.
  \item Le Cameroun n'est pas touché par les cyclones tropicaux, mais par des \textbf{pluies extrêmes}, des \textbf{inondations} et des \textbf{glissements de terrain} : adapter ses exemples au contexte local.
\end{itemize}
\end{attention}

\begin{aretenir}[Checklist avant le Probatoire]
\begin{itemize}
  \item[$\square$] Je sais nommer et situer les \textbf{trois cellules de convection} de chaque hémisphère (Hadley, Ferrel, polaire).
  \item[$\square$] Je sais expliquer le rôle de la \textbf{force de Coriolis} dans la déviation des vents et des courants.
  \item[$\square$] Je sais citer les grands \textbf{vents planétaires} : alizés, vents d'ouest des moyennes latitudes, vents d'est polaires.
  \item[$\square$] Je distingue \textbf{courants de surface} (dus aux vents) et \textbf{circulation thermohaline} (due à la densité : température + salinité).
  \item[$\square$] Je sais expliquer comment un océan chaud \textbf{alimente un cyclone} (évaporation, puis condensation qui libère de la chaleur et entretient la dépression).
  \item[$\square$] Je connais les conditions de formation d'un cyclone : eau à plus de $\SI{26}{\celsius}$ ; je sais qu'on parle de cyclone lorsque les vents dépassent $\SI{118}{km/h}$.
  \item[$\square$] Je sais définir \textbf{El Ni\~no} et citer une conséquence (sécheresse ou inondations anormales).
  \item[$\square$] Je sais décrire un mouvement atmosphérique ou océanique \textbf{à partir d'un schéma ou d'une carte} en quatre étapes.
  \item[$\square$] Je sais relier une \textbf{inondation} à ses causes dynamiques (pluies intenses, relief, déforestation, urbanisation).
  \item[$\square$] Je sais citer un exemple local : pluies extrêmes et inondations à Douala ou Yaoundé, glissements de terrain dans la région de l'Ouest camerounais.
  \item[$\square$] Je sais proposer au moins \textbf{trois mesures de protection} : surveillance météorologique, reboisement, plans d'évacuation et sensibilisation des populations.
  \item[$\square$] Je rédige mes réponses avec le \textbf{vocabulaire scientifique exact} (dépression, gradient de pression, convection, thermohaline).
\end{itemize}
\end{aretenir}

\findecours

\end{document}
